\documentclass[11pt,a4paper]{article}
\usepackage[utf8]{inputenc}
\usepackage[T1]{fontenc}
\usepackage{lmodern}
\usepackage{amsmath, amssymb}
\usepackage{graphicx}
\usepackage{booktabs}
\usepackage{geometry}
\usepackage{hyperref}
\usepackage{tcolorbox}
\usepackage{enumitem}
\usepackage{xcolor}
\usepackage{float}
\usepackage{multicol}

\geometry{margin=1in}
\hypersetup{
    colorlinks=true,
    linkcolor=blue,
    filecolor=magenta,      
    urlcolor=cyan,
    pdftitle={Natural Language Processing Study Guide},
}

% Global spacing for lists
\setlist[itemize]{itemsep=0.5em, topsep=0.4em, parsep=0.1em}
\setlist[enumerate]{itemsep=0.5em, topsep=0.4em, parsep=0.1em}

\title{\textbf{Natural Language Processing (NLTP)\\Complete Study Material}}
\author{Comprehensive Study Guide}
\date{}

\begin{document}

\maketitle
\tableofcontents
\newpage


\section{Module 1: Foundations of NLP}


\subsection{1 Introduction to Natural language processing (1)}

\textbf{Slide 1:}

Natural Language Text Processing

\textbf{Slide 2:}

Natural Language

Languages that  are spoken  naturally in  human being

\textbf{Slide 3:}

Languages

Human to Human

Machine to machine

\textbf{Slide 4:}

\textbf{Perspective on  NLP:} Areas of AI and their inter-dependencies.

\textbf{Slide 5:}

Layer 1

Search

Search algorithms try to find out the best possible strategy, the optimal strategy for computer.

Logic

Logic is a vehicle for reasoning and inferencing. In logic we are concerned with several constructs like if x is true then y is true.

C. Knowledge Representation

knowledge must be extracted and embedded in the Machine.

\textbf{Slide 6:}

Layer 2

Machine Learning

NLP is using lots of Statistical Techniques. Statistical Techniques are Machine Learning techniques; they make use of the knowledge content in the data.

Planning

Planning in AI is the process of producing a series of actions or procedures to accomplish a particular goal. Planning in AI is often used for tasks that require reasoning, problem-solving, and decision-making, such as robotics, games, or scheduling

\textbf{Slide 7:}

\textbf{Perspective on  NLP:} Areas of AI and their inter-dependencies.

NLP-concerned with computer being able to process human language like Hindi, Marathi,Gujrathi,French,English , and understand .

Computer Vision-NLP is followed by Computer Vision where machine processes seen and understand how to operate in seen.

Robotics-there is an embedded software inside  robots that asking it to perform certain actions like navigating etc.

\textbf{Slide 8:}

\textbf{Perspective on  NLP:} Areas of AI and their inter-dependencies.

Expert System-Expert System is concerned with ,the expert level performance of the software on a specific task.

For example –The task could be diagnosis of diseases and curing it.

A doctor is known to operate with number of rules ,a very large no of rules obtained by years of education and practice on patience .

So, the expert system is concerned with emulating this behavior of the expert.

\textbf{Slide 9:}

\textbf{2 nd layer :} feeding Layer

Machine Learning and planning feed into several layers in the outer most category.

For example, Natural Language Processing is fed by Machine Learning and Natural Language Processing is also fed by Knowledge Representation.

The reason for this is that in current world, Natural Language Processing is using lots of Statistical Techniques.

Statistical Techniques are Machine Learning techniques; they make use of

the knowledge content in the data.

\textbf{Slide 10:}

Introduction to Natural Language  Processing

branch of artificial intelligence  that helps computers understand, interpret and manipulate human  language.

it deals with the interaction between  computers and humans using the natural language.

objective  is to read, decipher, understand, and  make sense of the human languages in a manner that is valuable.

is an interdisciplinary subfield of linguistics, computer science, and artificial intelligence concerned with the interactions between computers and human language

\textbf{Slide 11:}

Natural language processing

It refers to the branch of computer science—and more specifically, the branch of artificial intelligence or AI

concerned with giving computers the ability to understand text and spoken words in much the same way human beings can.

is an interdisciplinary subfield of linguistics, computer science, and artificial intelligence concerned with the interactions between computers and human language

\textbf{Slide 12:}

INPUT Text  (Document

/Paragraph

/Sentense

/word}

OUTPUT

(Text/  Rating

/Graph

/Audio)

NLP System

\textbf{Slide 13:}

Text  (Document

/Paragraph

/Sen tense

/word}

Speech

Processing

Audio

OUTPUT

(Text/  Rating

/Graph)

NLP System

\textbf{Slide 14:}

Text  (Document

/Paragraph

/Sen tense

/word}

Speech

Processing

OCR

Processing

Audio

Image

OUTPUT

(Text/  Rating

/Graph)

NLP System

\textbf{Slide 15:}

Stages of Natural language Processing

To converse with human a program must understand the syntax(grammar),semantics(word meaning) and morphology (word level analysis),and pragmatics(conversation).

There are certain phases in which Natural language processing is performed .

\textbf{Slide 16:}

Steps in NLP

Morphological Analysis

Syntactic analysis

Semantic Analysis

Discourse Analysis

Pragmatic Analysis

Word-Level Analysis

Sentence-Level Analysis

Sentence-Level Analysis

Sentence-Level Analysis

Sentence-Level Analysis

\textbf{Slide 17:}

Morphological Analysis

studies the structure of words or formation of the words.

How words are built from smaller pieces

Identification ,analysis of root words , affixes (suffixes and prefixes)

Example

Washing – wash+ ​ing

Browser – Browse + er

Incomplete- In+ Complete

\textbf{Slide 18:}

Morphological Analysis

1. Tokenization​

John ate the pizza ! !

2. Stop Word Removal  (removing the words that occur commonly across all the documents ,typically, articles and pronouns are generally classified as stop words )​

​

Example

Search Engines

News article categorization & Spam Detection

\textbf{Chatbots \& Virtual Assistants:} 

“I want to know the status of my order”→ know status order

\textbf{Slide 19:}

Morphological Analysis

Syntactic analysis

Semantic Analysis

Discourse Analysis

Pragmatic Analysis

3. Stemming

Stemming is a process of reducing words  into its base form (Root form/stem form).

Process of  converting infected word to their word  stem.

John->John  Ate -> eat  Pizza->Pizza

\textbf{Slide 20:}

4. Lemmatization

text normalization technique used for Natural Language Processing (NLP).

It can convert any word's inflections to the base root form.

\textbf{For example:} 

Playing, Plays, Played ------- Play (Common root form "play")

\textbf{Slide 21:}

Stemming vs Lemmatization

\textbf{Stemming Real-world Uses:} 

Search engines

Information retrieval systems

Text mining

Recommendation systems

\textbf{Lemmatization Real-world Uses:} 

Chatbots

Machine translation

Medical & legal NLP

Sentiment analysis

\textbf{Slide 22:}

Stemming vs Lemmatization

\textbf{Slide 23:}

Morphological Analysis

Syntactic analysis

Semantic Analysis

Discourse Analysis

Pragmatic Analysis

4. N-Gram Language Model

Continuous sequence of N-Items from a given sample text.

1-gram-

John

Ate

the

Pizza

Bigram -

Trigram -

4-Gram -

John	Ate	the	Pizza

John Ate the ?

John

Ate

the

Pizza

\textbf{Slide 24:}

Syntactic analysis

John Ate the Apple

Ate	the Apple John

analyzing the grammatical syntax of a sentence to understand its meaning

\textbf{Slide 25:}

Syntactic analysisset of rules needed to ensure a sentence is grammatically correct

John Ate the Apple

Ate	the Apple John

\textbf{Slide 26:}

Morphological Analysis

Semantic Analysis

Discourse Analysis

Pragmatic Analysis

She drank Some Milk

She drank Some books

semantics refers to meaning

a computer understands the meaning of a text by analyzing the text as a whole and not just looking at individual words

context in which a word is used is very important

“Does it all sound like a joke to you?”

Example 1

Example 2

\textbf{Slide 27:}

Morphological Analysis

Syntactic analysis

Semantic Analysis

Discourse Analysis

Pragmatic Analysis

Monkeys Eat Banana,  when they Wake up.

Who is they here?

-Monkey

Monkeys eat Banana,

when they are ripe.

Who is they here?

-Banana

Resolving the reference as per the context .

\textbf{Slide 28:}

Morphological Analysis

Syntactic analysis

Semantic Analysis

Discourse Analysis

Pragmatic Analysis

Close the Door

-Order

Please Close the Door

- Request ,affirmation

knowledge of the relationship of meaning to the goals and intentions of the speaker.

\textbf{Slide 29:}

Ambiguity in Natural language

Text is ambiguous if there are  multiple  alternative Ambiguous linguistic structures that can be built for it

situation where a word or a sentence may have more than one meaning.

There are different types of ambiguities

Lexical

Syntactical

Semantic

Discourse Ambiguity

Pragmatic Ambiguity

\textbf{Slide 30:}

Lexical Ambiguity

is the ambiguity of a single word

It can be resolved by parts of speech tagging

Word has more than one meaning/category

E.g.

Book – Noun –Textbook/Novel

Book – Verb - Book ticket/seat

Bank-Noun –Financial institute

Bank-Noun - River Bank

Bank- Verb-  Banking Transaction

\textbf{Slide 31:}

Syntactical Ambiguity(Grammar or rules are ambiguous)

it exists in the presence of two or more possible meanings within the sentence.

Specify the possible arrangements of  words in sentence used

Example

I saw the girl with the binocular.

In the above example, did I have the binoculars? Or did the girl have the binoculars?

\textbf{Slide 32:}

Semantic Ambiguity

When a sentence has more than one meaning then it is called as semantic ambiguity

Rahul loves his cat and Dipesh does too.

Whether Dipesh loves his cat or Rahul's cat

\textbf{Slide 33:}

Anaphoric Ambiguity

Occurs due to use of Anaphoric entities in discourse.

Anaphora-When same beginning of sentence is repeated in the sentence several times , we makes  the pronoun instead of noun .

For example, My mother liked the house very much , but she couldn't purchase it

\textbf{Example:} -

•Monkeys Eat Banana, when they  Wake up. -Who is they here? Monkey

•Monkeys eat Banana, when they  are ripe.-Who is they here?-Banana

Identify When, Where, by whom  occurrence was said

\textbf{Slide 34:}

Pragmatic Ambiguity

Understanding speaker's intention.

Eg. You are Late

\textbf{Slide 35:}

Challenges in NLP

Elongated words

I am sooooo sorry.

I t was toooo yummy.

Shortcuts

Pls	-

BTW	-

Wat	-

K	-

Please

By the way  What

OK

\textbf{Slide 36:}

Challenges in NLP

Emojis

\textbf{Slide 37:}

Challenges in NLP

Mix Use of Languages

I liked that movie. Salman khan ka acting was  Lajabab. Too good. जबरदस्त.

\textbf{Slide 38:}

Challenges in NLP

Ellipsis

Peter worked hard and passed exam, Kevin too.

Interpretation

Kevin worked hard

Kevin passed exam

Kevin worked hard also he passed exam

\textbf{Slide 39:}

Challenges in NLP

Punctuational Ambiguity

Women, without her man, is nothing.

Women! without her, man is nothing.

\textbf{Slide 40:}

Applications of NLP

1. Question Answering

Question Answering focuses on building systems that automatically answer the questions asked by humans in a natural language.

\textbf{Slide 41:}

Applications of NLP

2. Spam Detection

Spam detection is used to detect unwanted e-mails getting to a user's inbox.

\textbf{Slide 42:}

3. Sentiment Analysis

A combination of NLP (Natural Language Processing) and statistics by assigning the values to the text (positive, negative, or natural), identify the mood of the context (happy, sad, angry, etc.)​

\textbf{Slide 43:}

\textbf{4. Machine Translation (Example:}  Google Translator)

Machine translation is used to translate text or speech from one natural language to another natural language.

\textbf{Slide 44:}

5. Spelling correction

Microsoft Corporation provides word processor software like MS-word, PowerPoint for the spelling correction.

Chek's the spelling while typing.

Example Wrong Speling will be underlined with the red color as done in this sentence

\textbf{Slide 45:}

6. Chatbot

Implementing the Chatbot is one of the important applications of NLP. It is used by many companies to provide the customer's chat services.

\textbf{Slide 46:}

Practice Questions

Identify the Type Of Ambiguity for Given Sentence

Identify The Type of Preprocessing with respect to given scenario-based corpus

Steps in NLP

Challenges in NLP

Ambiguities in NLP

\textbf{Slide 47:}

Lexical Ambiguity

Occurs when a word has multiple meanings.

\textbf{Example:}  "The fisherman went to the bank."

"Bank" could mean a financial institution or the riverbank (edge of a river).

This ambiguity arises due to polysemy (multiple meanings of a word).

Syntactic (Structural) Ambiguity

Occurs when a sentence can have multiple grammatical interpretations due to different ways of parsing it.

\textbf{Example:}  "The manager told the employee that he was fired."

Who is fired? The manager or the employee?

The sentence structure allows both interpretations, leading to confusion.

Semantic Ambiguity

Occurs when the meaning of a sentence is unclear due to vague or multiple possible meanings.

\textbf{Example:}  "Visiting relatives can be annoying."

Does it mean relatives who visit are annoying or the act of visiting relatives is annoying?

Pragmatic Ambiguity

Occurs when the meaning depends on context, tone, or external knowledge.

\textbf{Example:}  "Can you pass the salt?"

Literally, this asks if the person has the ability to pass salt.

But in a real-world context, it is a request to pass the salt.

Referential Ambiguity

Occurs when pronouns (he, she, it, they) refer to multiple possible antecedents.

\textbf{Example:}  "John met Mike at the café, and he ordered a coffee."

Who ordered the coffee—John or Mike?

\textbf{Slide 48:}

Practice Questions

Q.  Identify the Type Of Ambiguity for Given Sentence

Consider the following sentences and identify type of ambiguity

The manager told the employee that he was fired.

The fisherman went to the bank.

\textbf{Solution:} 

"The manager told the employee that he was fired."

\textbf{Type of Ambiguity:}  Syntactic Ambiguity / Referential Ambiguity

Unclear whether "he" refers to the manager or the employee.

"The fisherman went to the bank."

\textbf{Type of Ambiguity:}  Lexical Ambiguity

The word "bank" can mean a financial institution or a riverbank.

\textbf{Slide 49:}

Consider the following sentences and identify type of ambiguity

He saw the seal at the zoo.

John told Mike that he won the award.

\textbf{Slide 50:}

Identify The Type of Preprocessing with respect to given scenario-basedcorpus

Q1 .A search engine ignores common words such as “the”, “is”, “and”, “of” while indexing web pages to improve search speed and relevance. Which text preprocessing technique is required?

Q2. A customer review sentiment analysis system removes symbols like !, ?, \#, and . before extracting features from text. Which text preprocessing technique is required?

Q3. A medical chatbot converts words such as “diagnoses”, “diagnosed”, and “diagnosing” into their base dictionary form “diagnose” to preserve meaning. Which text preprocessing technique is required?

Q4. A news classification system treats “COVID”, “Covid”, and “covid” as the same term during model training. Which text preprocessing technique is required?

Q5. A document retrieval system reduces words like “playing”, “played”, and “player” to “play” without checking grammatical correctness. Which text preprocessing technique is required?

\textbf{Slide 51:}

Q6

A social media analytics platform removes URLs, HTML tags, emojis, and special characters from tweets before analysis. Which text preprocessing technique is required?

Q7

A plagiarism detection system breaks paragraphs into individual words before computing similarity measures. Which text preprocessing technique is required?

Q8

An e-commerce search engine automatically corrects misspelled queries such as “samsng” to “samsung”. Which text preprocessing technique is required?

Q9

A financial report analysis system replaces values like “₹10,000”, “2023”, and “5%” with a generic token to reduce vocabulary size. Which text preprocessing technique is required?

Q10

A sentiment analysis system captures expressions like “not satisfied” and “very good” as meaningful units instead of individual words. Which text preprocessing technique is required?

\textbf{Slide 52:}

Answers


\subsection{2 Language Model_KD}

\textbf{Slide 1:}

Language Modeling

Introduction to N-grams

\textbf{Slide 2:}

Why?

\textbf{Machine Translation:} 

P(high winds tonite) > P(large winds tonite)

Spell Correction

The office is about fifteen minuets from my house

P(about fifteen minutes from) > P(about fifteen minuets from)

\textbf{Slide 3:}

Why ?

Speech Recognition

P(I saw a van) >> P(eyes awe of an)

+ Summarization, question-answering, etc., etc.!!

OCR(optical character recognition) & Hand Recognition

In, OCR LM can be used to predict a word that is not easily readable in the given image .

\textbf{Slide 4:}

Why ?

Suggestions while typing

We get suggestion on our cell phone while messaging

Context Sensitive Spelling Correction

Eg. "college is closed on Saturday and Sunday

\textbf{Spelling error :}  college has been typed as collage

\textbf{College :} Education institution

\textbf{Collage :}  Artwork

Suppose both words college and collage are present in corpus .

Spelling Mistake detection and correction will not work as both words are present in corpus.

Need Spelling Correction based on local context.

P(College is closed on Saturday and Sunday )> p(Collage is closed on Saturday and Sunday )

\textbf{Slide 5:}

Probabilistic Language Modeling

\textbf{compute the probability of a sentence or sequence of words occurring together:} 

P(W) = P(w1,w2,w3,w4,w5…wn)   Joint Probability

\textbf{probability of an upcoming word:} 

P(w5|w1,w2,w3,w4) Conditional Probability

\textbf{Slide 6:}

How to estimate these probabilities

Could we just count and divide?

No!  Too many possible sentences!

We’ll never see enough data for estimating these

\textbf{Slide 7:}

Drawback of estimating probabilities directly from counts

Large corpora (like the web) allow probability estimation using word counts.

This approach often fails because even huge datasets are not large enough.

Language is productive and constantly creates new sentences.

Many valid sentences or phrases never appear in the data.

Unseen sentences or phrases get zero counts, leading to poor probability estimates.

Solution is join probability

\textbf{Slide 8:}

Joint Probability  P(W)

\textbf{How to compute this joint probability:} 

P(its, water, is, so, transparent, that)

\textbf{Intuition:}  let’s rely on the Chain Rule of Probability

\textbf{Slide 9:}

\textbf{Reminder:}  The Chain Rule

Recall the definition of conditional probabilities

\textbf{p(B|A) = P(A,B)/P(A)	Rewriting:}    P(A,B) = P(A)P(B|A)

\textbf{More variables:} 

P(A,B,C,D) = P(A)P(B|A)P(C|A,B)P(D|A,B,C)

The Chain Rule in General

P(x1,x2,x3,…,xn) = P(x1)P(x2|x1)P(x3|x1,x2)…P(xn|x1,…,xn-1)

\textbf{Slide 10:}

Markov Assumption

In other words, we approximate each component in the product

\textbf{Slide 11:}

Markov Assumption

\textbf{Simplifying assumption:} 

Or maybe

Andrei Markov

\textbf{Slide 12:}

N-gram model

A model that applies markov assumption to text

For Example, bigram

The bigram model, approximates  the conditional probability of the preceding word P(wn|wn−1).

\textbf{Slide 13:}

Example

instead of computing the probability

P(the|Walden Pond’s water is so transparent that)

we approximate it with the probability

P(the|that)

\textbf{Slide 14:}

How do we estimate these bigram or n-gram probabilities?

The Maximum Likelihood Estimate

\textbf{Slide 17:}

An example

\textbf{Corpus:} 

<s> I am Sam </s>

<s> Sam I am </s>

<s> I do not like green eggs and ham </s>

\textbf{Slide 18:}

\textbf{Exercise 1:} 

Estimating Bi-gram probabilities

What is the most probable next word predicted by the model for the following word sequence?

<S> Do ?

<S> I like Henry ?

\textbf{Slide 19:}

1) <S> Do ?

Next word prediction probability Wᵢ₋₁ = do

I is most Probable

\textbf{Slide 20:}

<S> I like Henry ?

Next word prediction probability Wᵢ₋₁ = Henry

</S> is more probable

\textbf{Slide 21:}

Questions

\textbf{Slide 22:}

Solution

\textbf{Slide 23:}

Problems with Using Raw Probabilities

1. Numerical Underflow

2. Computational Inefficiency

\textbf{Slide 24:}

Practical Issues

We do everything in log space

Avoid underflow

(also adding is faster than multiplying)

\textbf{Slide 25:}

Probabilities ≤ 1 shrink when multiplied repeatedly

Multiplying many n-gram probabilities causes numerical underflow

Using log probabilities avoids very small numbers

Multiplication in probability space becomes addition in log space

Addition is faster and more stable than multiplication

Convert back to probabilities only when final results are needed

log(p1 X p2 X p3 X p4 ) =log p1 +log p2 + log p3 + log p4

Log Probabilities

\textbf{Slide 26:}

Example

\textbf{Slide 27:}

Example

\textbf{Slide 28:}

Thank You

\textbf{Slide 29:}

Practice Question

\textbf{Extract bigrams from the given corpus:} 

\textbf{Corpus C:} 

there is a big garden

children play in a garden

they play inside beautiful garden

\textbf{Slide 30:}

Practice Question 1

\textbf{Extract bigrams from the given corpus:} 

\textbf{Corpus C:} 

there is a big garden

children play in a garden

they play inside beautiful garden

Preprocess

\textbf{Adding sentence boundaries <s> and </s>:} 

<s> there is a big garden </s>

<s> children play in a garden </s>

<s> they play inside beautiful garden </s>

\textbf{Slide 31:}

Practice Question

\textbf{extract bigrams from the given corpus:} 

\textbf{Corpus C:} 

there is a big garden

children play in a garden

they play inside beautiful garden

\textbf{Adding sentence boundaries <s> and </s>:} 

<s> there is a big garden </s>

<s> children play in a garden </s>

<s> they play inside beautiful garden </s>

\textbf{Bigram Frequency Table:} 

\textbf{Slide 32:}

Practice Question 2

Compute P(in|play) using a Bigram Model

\textbf{Corpus C:} 

there is a big garden

children play in a garden

they play inside beautiful garden

\textbf{Slide 33:}

P(in|play)

Practice Question 2

Compute P(in|play) using a Bigram Model

\textbf{Corpus C:} 

there is a big garden

children play in a garden

they play inside beautiful garden

\textbf{Slide 34:}

Practice Question-3

Compute P(they play in a big garden) using a Bigram Model

\textbf{Corpus C:} 

there is a big garden

children play in a garden

they play inside beautiful garden

\textbf{Adding sentence boundaries <s> and </s>:} 

<s> there is a big garden </s>

<s> children play in a garden </s>

<s> they play inside beautiful garden </s>

\textbf{Unigram Frequency Table:} 

\textbf{Bigram Frequency Table:} 

\textbf{Slide 35:}

Practice Question

Compute P(they play in a big garden) using a Bigram Model

\textbf{Corpus C:} 

there is a big garden

children play in a garden

they play inside beautiful garden

\textbf{Unigram Frequency Table:} 

\textbf{Bigram Frequency Table:} 

\textbf{Slide 37:}

Calculations are made easy by a.b.c=exp.log(a.b.c)

a⋅b⋅c=exp(loga+logb+logc)

\textbf{Slide 38:}

Practice Question 6

"there is a big garden"

\textbf{Trigrams:}  (there, is, a), (is, a, big), (a, big, garden)

"children play in a garden"

\textbf{Trigrams:}  (children, play, in), (play, in, a), (in, a, garden)

"they play inside beautiful garden"

Trigrams: (they, play, inside), (play, inside, beautiful), (inside, beautiful, garden)

Total Trigrams

Total unique trigrams = 9

Total trigrams including repetitions = 9 (No duplicates here)

\textbf{Trigram Counts:} 

\textbf{Bigram Counts:} 


\subsection{3 Evaluating Language Models_}

\textbf{Slide 1:}

Evaluating Language Models

Natural Language Processing

1

\textbf{Slide 2:}

Natural Language Processing

2

Does our language model prefer good sentences to bad ones?

Assign higher probability to “real” or “frequently observed” sentences than  “ungrammatical” or “rarely observed” sentences?

We train parameters of our model on a training set.

We test the model’s performance on data we haven’t seen.

A test set is an unseen dataset that is different from our training set, totally unused.

An evaluation metric tells us how well our model does on the test set.

Evaluating Language Models

\textbf{Slide 3:}

Natural Language Processing

3

Evaluating Language Models

Extrinsic Evaluation(Application Based)

Extrinsic Evaluation of a N-gram language model is to use it in an application and  measure how much the application improves.

\textbf{To compare two language models A and B:} 

Use each of language model in a task such as spelling corrector, MT(Machine Translation) system.

Get an accuracy for A and for B

How many misspelled words corrected properly

How many words translated correctly

Compare accuracy for A and B

The model produces the better accuracy is the better model.

Extrinsic evaluation can be time-consuming.

\textbf{Slide 4:}

Natural Language Processing

4

Evaluating Language Models

Intrinsic Evaluation

An intrinsic evaluation metric is one that measures the quality of a model  independent of any application.

When a corpus of text is given and to compare two different n-gram models,

Divide the data into training and test sets,

Train the parameters of both models on the training set, and

Compare how well the two trained models fit the test set.

Whichever model assigns a higher probability to the test set

In practice, probability-based metric called perplexity is used instead of raw  probability as our metric for evaluating language models.

\textbf{Slide 5:}

Evaluating Language Models

Perplexity

The best language model is one that best predicts an unseen test set

–	Gives the highest P(testset)

The perplexity of a language model on a test set is the inverse probability of the test  set, normalized by the number of words.

Minimizing perplexity is the same as maximizing probability

The perpelexity PP for a test set W=w1w2…wn is

PP(W)	by chain rule

\textbf{The perpelexity PP for bigrams:}   PP(W)

Natural Language Processing

5

\textbf{Slide 6:}

Evaluating Language Models

Perplexity as branching factor

Perplexity can be seen as the weighted average branching factor of a language.

–	The branching factor of a language is the number of possible next words that can follow any word.

Let’s suppose a sentence consisting of random digits

What is the perplexity of this sentence according to a model that assign P=1/10 to each  digit?

6

\textbf{Slide 7:}

Example

<s>I am a human</s>

<s>I am not a stone </s>

</s>I I live in Mumbai</s>

\textbf{Using a bigram language model compute:} P("I I am not")

\textbf{Slide 8:}

Natural Language Processing

8

The n-gram model, like many statistical models, is dependent on the training corpus. the probabilities often encode specific facts about a given training corpus.

N-grams only work well for word prediction if the test corpus looks like the training corpus In real life, it often doesn’t We need to train robust models that generalize!

One kind of generalization: Getting rid of Zeros! Things that don’t ever occur in the training set, but occur in the test set Zeros: things that don’t ever occur in the training set but do occur in the test set causes problem for two reasons. First, underestimating the probability of all sorts of words that might occur, Second, if probability of any word in test set is 0, entire probability of test set is 0.

Generalization and Zeros

\textbf{Slide 9:}

Natural Language Processing

9

N-gram models depend on the training corpus

Perform well only if test data matches training data

Zero probabilities occur for unseen words

One zero makes the entire test probability zero

Generalization and Zeros

\textbf{Slide 10:}

Natural Language Processing

10

We have to deal with words we haven’t seen before, which we call unknown words.

We can model these potential unknown words in the test set by adding a pseudo-word

called <UNK> into our training set too.

\textbf{One way to handle unknown words is:} 

Replace words in the training data by <UNK> based on their frequency.

For example we can replace by <UNK> all words that occur fewer than n times in the  training set, where n is some small number, or

Equivalently select a vocabulary size V in advance (say 50,000) and choose the top V  words by frequency and replace the rest by <UNK>.

Proceed to train the language model as before, treating <UNK> like a regular word.

Unknown Words

\textbf{Slide 11:}

11

Real data contains unknown words not seen in training

Unknown words are modeled using a pseudo-word <UNK>

<UNK> is added to the training vocabulary

Rare words in training are replaced with <UNK>

Or fix vocabulary size V and keep top V frequent words

Replace all remaining words with <UNK>

Train the language model treating <UNK> as a normal word

Unknown Words

\textbf{Slide 12:}

<s>I am Henry</s>

<s> I like college </s>

<s>Do Henry like college </s>

<s>Henry I am</s>

<s>Do I like Henry</s>

<s>Do I like college </s>

<s>I do like Henry </s>

<s> like college</s>

<s>Do I like Henry</s>

Consider the Following Corpus and Identify the most Probable statement from the given sentence

\textbf{Slide 13:}

<s> like college </s>

= P(like | <s>) × P(college | like) × P(</s> | college)

= 0/7 × 3/5 × 3/3 = 0

<s> Do I like Henry </s>

= P(do | <s>) × P(I | do) × P(like | I) × P(Henry | like) × P(</s> | Henry)

= 3/7 × 2/4 × 3/6 × 2/5 × 3/5

= 9/350 = 0.0257

\textbf{Slide 14:}

Natural Language Processing

14

We lower the probability of common word pairs a little and use that leftover probability to give rare or unseen word pairs a small non-zero probability.

This modification is called smoothing (or discounting).

\textbf{With smoothing:} 

Unseen words get a small chance

The model becomes more realistic and useful

Smoothing

\textbf{Slide 15:}

15

\textbf{There are many ways to do smoothing, and some of themare:} 

Add-1 smoothing (Laplace Smoothing)

Add-k smoothing,

Smoothing

\textbf{Slide 16:}

Natural Language Processing

16

The simplest way to do smoothing is to add one to all the counts, before we  normalize them into probabilities.

–	All the counts that used to be zero will now have a count of 1, the counts of 1 will be 2, and  so on.

This algorithm is called Laplace smoothing (Add-1 Smoothing).

We pretend that we saw each word one more time than we did, and we just add one to

all the counts!,

Laplace Smoothing

\textbf{Slide 17:}

17

The unsmoothed maximum likelihood estimate of the unigram probability of the word wi is its count ci normalized by the total number of word tokens N:

P(wi) = ci / N

Laplace smoothing adds one to each count. Since there are V words in the vocabulary  and each one was incremented, we also need to adjust the denominator to take into  account the extra V observations.

PLaplace(wi) = (ci + 1) / (N + V)

Laplace Smoothing  ( Add-1 Smoothing )

Where

N = total word tokens in corpus,

V = number of unique words in the vocabulary.

\textbf{Slide 18:}

The normal bigram probabilities are computed by normalizing each bigram counts  by the unigram count:

\textbf{Add-one smoothed bigram probabilities:} 

Laplace Smoothing for Bigrams

Natural Language Processing

18

\textbf{Slide 19:}

\textbf{Unique words are:}  <S>, </S>, I, Henry, do, like, am, college

\textbf{Total unique words:}  8

But we exclude </S>

\textbf{Total unique words:}  7

Give the following bi-gram probabilities estimated by Laplace model.

\textbf{Slide 20:}

1. <S> like college </S>

= P(like | <S>) × P(college | like) × P(</S> | college)

= (0 + 1) / (7 + 7) × (3 + 1) / (5 + 7) × (3 + 1) / (3 + 7)

= 1/14 × 4/12 × 4/10

= 0.0095

2. <S> Do I like Henry </S>

= P(do | <S>) × P(I | do) × P(like | I) × P(Henry | like) × P(</S> | Henry)

= (3 + 1) / (7 + 7) × (2 + 1) / (4 + 7) × (3 + 1) / (6 + 7) × (2 + 1) / (5 + 7) × (3 + 1) / (5 + 7)

= 4/14 × 3/11 × 4/13 × 3/12 × 4/12

= 0.0020

\textbf{Slide 21:}

Now consider a Bi-gram Language model (without smoothing). Which of the following sentences has the highest probability estimate?

a. S1: <s> michael played at the playground </s>b. S2: <s> bob went to the school </s>c. S3: <s> the school was huge </s>d. S4: <s> zack went to the playground </s>

P(S1) = 0P(S2) = 0P(S3) = 2/5 × 3/5 × 2/5 × 1/2 × 1 = 2/25P(S4) = 0

\textbf{Slide 22:}

Laplace-smoothed Bigrams

\textbf{Corpus:}  Berkeley Restaurant Project Sentences

Natural Language Processing

22

\textbf{Slide 23:}

Laplace-smoothed Bigrams

\textbf{Corpus:}  Berkeley Restaurant Project Sentences: Adjusted counts

Natural Language Processing

23

\textbf{Slide 24:}

Laplace-smoothed Bigrams

\textbf{Corpus:}  Berkeley Restaurant Project Sentences: Adjusted counts

Natural Language Processing

24

\textbf{Slide 25:}

Problem with Add-1 Smoothing

\textbf{We fix zeros by pretending every possible bigram happened once:} 

Where V = vocabulary size.

\textbf{But here’s the issue:} 

Adding 1 to every possible word is huge when V is big.

\textbf{From the slide:} 

C(want → to) changed 608 → 238 after smoothing

P(to | want) dropped 0.66 → 0.26

\textbf{So we:} 

Steal probability from common events

Give too much to rare/unseen events

This is called over-smoothing.

\textbf{Slide 26:}

The sharp change in counts and probabilities occurs because too much probability

mass is moved to all the zeros.

One alternative to add-one smoothing is to move a bit less of the probability mass  from the seen to the unseen events.

Instead of adding 1 to each count, we add a fractional count k (.5? .05? .01?).

This algorithm is called add-k smoothing.

Add-k Smoothing

26

\textbf{What this does:} 

\textbf{Slide 27:}

Numerical

Consider the following Corpus

<s> I love natural language processing </s>

<s> language models learn patterns </s>

<s> I love machine learning </s>

<s> natural language models are powerful </s>

<s> machine learning uses data </s>


\subsection{4 EditDistance}

\textbf{Slide 1:}

Minimum Edit Distance

Definition of Minimum Edit Distance

\textbf{Slide 2:}

How similar are two strings?

Spell correction

The user typed “graffe”

Which is closest?

graf

graft

grail

giraffe

Computational Biology

Align two sequences of nucleotides

\textbf{Resulting alignment:} 

Also for Machine Translation, Information Extraction, Speech Recognition

AGGCTATCACCTGACCTCCAGGCCGATGCCC

TAGCTATCACGACCGCGGTCGATTTGCCCGAC

-AGGCTATCACCTGACCTCCAGGCCGA--TGCCC---

TAG-CTATCAC--GACCGC--GGTCGATTTGCCCGAC

\textbf{Slide 3:}

Edit Distance

Edit distance gives us a way to quantify intuitions about string similarity.

The minimum edit distance between two strings

Is the minimum number of editing operations

Insertion

Deletion

Substitution

Needed to transform one into the other

\textbf{Slide 4:}

Minimum Edit Distance

\textbf{Two strings and their alignment:} 

\textbf{Slide 5:}

Minimum Edit Distance

If each operation has cost of 1

Distance between these is 5

If substitutions cost 2 (Levenshtein)

Distance between them is 8

\textbf{Slide 6:}

Alignment in Computational Biology

Given a sequence of bases

\textbf{An alignment:} 

Given two sequences, align each letter to a letter or gap

-AGGCTATCACCTGACCTCCAGGCCGA--TGCCC---

TAG-CTATCAC--GACCGC--GGTCGATTTGCCCGAC

AGGCTATCACCTGACCTCCAGGCCGATGCCC

TAGCTATCACGACCGCGGTCGATTTGCCCGAC

\textbf{Slide 7:}

Alignment in Computational Biology

Given a sequence of bases

\textbf{An alignment:} 

Given two sequences, align each letter to a letter or gap

-AGGCTATCACCTGACCTCCAGGCCGA--TGCCC---

TAG-CTATCAC--GACCGC--GGTCGATTTGCCCGAC

AGGCTATCACCTGACCTCCAGGCCGATGCCC

TAGCTATCACGACCGCGGTCGATTTGCCCGAC

\textbf{Slide 8:}

Other uses of Edit Distance in NLP

Evaluating Machine Translation and speech recognition

R Spokesman confirms    senior government adviser was appointed

H Spokesman said    the senior            adviser was appointed

S      I              D                        I

Named Entity Extraction and Entity Coreference

IBM Inc. announced today

IBM profits

Stanford Professor Jennifer Eberhardt announced yesterday

for Professor Eberhardt…

\textbf{Slide 9:}

How to find the Min Edit Distance?

Searching for a path (sequence of edits) from the start string to the final string:

\textbf{Initial state:}  the word we’re transforming

\textbf{Operators:}  insert, delete, substitute

\textbf{Goal state:}   the word we’re trying to get to

\textbf{Path cost:}  what we want to minimize: the number of edits

\textbf{Slide 10:}

Minimum Edit as Search

But the space of all edit sequences is huge!

We can’t afford to navigate naïvely

Lots of distinct paths wind up at the same state.

We don’t have to keep track of all of them

Just the shortest path to each of those revisted states.

\textbf{Slide 11:}

Defining Min Edit Distance

For two strings

X of length n

Y of length m

We define D(i,j)

the edit distance between X[1..i] and Y[1..j]

i.e., the first i characters of X and the first j characters of Y

The edit distance between X and Y is thus D(n,m)

\textbf{Slide 12:}

Minimum Edit Distance

Definition of Minimum Edit Distance

\textbf{Slide 13:}

Dynamic Programming for Minimum Edit Distance

\textbf{Dynamic programming:}  A tabular computation of D(n,m)

Solving problems by combining solutions to subproblems.

Bottom-up

We compute D(i,j) for small i,j

And compute larger D(i,j) based on previously computed smaller values

i.e., compute D(i,j) for all i (0 < i < n)  and j (0 < j < m)

\textbf{Slide 14:}

Defining Min Edit Distance (Levenshtein)

Initialization

D(i,0) = i

D(0,j) = j

\textbf{Recurrence Relation:} 

For each  i = 1…M

For each  j = 1…N

D(i-1,j) + 1

D(i,j)= min   D(i,j-1) + 1

D(i-1,j-1) +   2;  if X(i) ≠ Y(j)

0;  if X(i) = Y(j)

\textbf{Termination:} 

D(N,M) is distance

\textbf{Slide 15:}

The Edit Distance Table

j

i

\textbf{Slide 16:}

The Edit Distance Table

\textbf{Slide 17:}

Edit Distance

\textbf{Slide 18:}

The Edit Distance Table

\textbf{Slide 19:}

Computing alignments

Edit distance isn’t sufficient

We often need to align each character of the two strings to each other

We do this by keeping a “backtrace”

Every time we enter a cell, remember where we came from

When we reach the end,

Trace back the path from the upper right corner to read off the alignment

\textbf{Slide 20:}

Edit Distance

\textbf{Slide 21:}

Exercise

\textbf{Slide 22:}

Solution

\textbf{Slide 23:}

MinEdit with Backtrace

\textbf{Slide 24:}

Adding Backtrace to Minimum Edit Distance

Base conditions:                                                        Termination:

D(i,0) = i         D(0,j) = j         D(N,M) is distance

\textbf{Recurrence Relation:} 

For each  i = 1…M

For each  j = 1…N

D(i-1,j) + 1

D(i,j)=     min   D(i,j-1) + 1

D(i-1,j-1) +  2; if X(i) ≠ Y(j)

0; if X(i) = Y(j)

LEFT

ptr(i,j)=   DOWN

DIAG

insertion

deletion

substitution

insertion

deletion

substitution

\textbf{Slide 25:}

The Distance Matrix

Slide adapted from Serafim Batzoglou with Permission

y0 ………………………………  yM

x0 ……………………  xN

Every non-decreasing path

from (0,0) to (M, N)

corresponds to

an alignment

of the two sequences

An optimal alignment is composed of optimal subalignments

\textbf{Slide 26:}

Result of Backtrace

\textbf{Two strings and their alignment:} 

\textbf{Slide 27:}

Minimum Edit Distance

Weighted Minimum Edit Distance

\textbf{Slide 28:}

Weighted Edit Distance

Why would we add weights to the computation?

\textbf{Spell Correction:}  some letters are more likely to be mistyped than others

\textbf{Biology:}  certain kinds of deletions or insertions are more likely than others

\textbf{Slide 29:}

Confusion matrix for spelling errors

\textbf{Slide 31:}

Weighted Min Edit Distance

\textbf{Initialization:} 

D(0,0) = 0

D(i,0) = D(i-1,0) + del[x(i)];    1 < i ≤ N

D(0,j) = D(0,j-1) + ins[y(j)];    1 < j ≤ M

\textbf{Recurrence Relation:} 

D(i-1,j)   + del[x(i)]

D(i,j)= min  D(i,j-1)   + ins[y(j)]

D(i-1,j-1) + sub[x(i),y(j)]

\textbf{Termination:} 

D(N,M) is distance

\textbf{Slide 32:}

Question

Find the minimum edit distance in transforming the word DOG to COW using Levenshtein distance, ie., insertion = deletion =1 and substitution = 2.

\textbf{Slide 33:}

Solution

\textbf{Slide 34:}

Example


\subsection{Chapter 1 Ngram}

\textbf{Page 1:}

Speech and Language Processing. Daniel Jurafsky & James H. Martin. Copyright © 2026. All
rights reserved. Draft of January 6, 2026.
CHAPTER
3 N-gram Language Models
“Youareuniformlycharming!”criedhe,withasmileofassociatingandnow
andthenIbowedandtheyperceivedachaiseandfourtowishfor.
RandomsentencegeneratedfromaJaneAustentrigrammodel
Predicting is difficult—especially about the future, as the old quip goes. But how
aboutpredictingsomethingthatseemsmucheasier, likethenextwordsomeoneis
goingtosay? W

\textbf{Page 2:}

2 CHAPTER3 • N-GRAMLANGUAGEMODELS
language model. An n-gram is a sequence of n words: a 2-gram (which we’ll call
bigram)isatwo-wordsequenceofwordslikeThe water,orwater of,anda3-
gram(atrigram)isathree-wordsequenceofwordslikeThe water of,orwater
of Walden. But we also (in a bit of terminological ambiguity) use the word ‘n-
gram’tomeanaprobabilisticmodelthatcanestimatetheprobabilityofawordgiven
then-1previouswords,andtherebyalsotoassignprobabilitiestoentiresequences.
In later chapters we will intr

\textbf{Page 3:}

3.1 • N-GRAMS 3
bility:
P(X ...X ) = P(X )P(X |X )P(X |X )...P(X |X )
1 n 1 2 1 3 1:2 n 1:n−1
n
(cid:89)
= P(X k |X 1:k−1 ) (3.3)
k=1
Applyingthechainruletowords,weget
P(w ) = P(w )P(w |w )P(w |w )...P(w |w )
1:n 1 2 1 3 1:2 n 1:n−1
n
(cid:89)
= P(w k |w 1:k−1 ) (3.4)
k=1
Thechainruleshowsthelinkbetweencomputingthejointprobabilityofasequence
and computing the conditional probability of a word given previous words. Equa-
tion3.4suggeststhatwecouldestimatethejointprobabilityofanentiresequenceof
wo

\textbf{Page 4:}

4 CHAPTER3 • N-GRAMLANGUAGEMODELS
size,soN=2meansbigramsandN=3meanstrigrams. Thenweapproximatethe
probabilityofawordgivenitsentirecontextasfollows:
P(w n |w 1:n−1 )≈P(w n |w n−N+1:n−1 ) (3.8)
Giventhebigramassumptionfortheprobabilityofanindividualword,wecancom-
putetheprobabilityofacompletewordsequencebysubstitutingEq.3.7intoEq.3.4:
n
(cid:89)
P(w )≈ P(w |w ) (3.9)
1:n k k−1
k=1
3.1.2 Howtoestimateprobabilities
How do we estimate these bigram or n-gram probabilities? An intuitive way to
maximum


\textbf{Page 5:}

3.1 • N-GRAMS 5
Equation 3.12 (like Eq. 3.11) estimates the n-gram probability by dividing the
observedfrequencyofaparticularsequencebytheobservedfrequencyofaprefix.
relative This ratio is called a relative frequency. We said above that this use of relative
frequency
frequenciesasawaytoestimateprobabilitiesisanexampleofmaximumlikelihood
estimationorMLE.InMLE,theresultingparametersetmaximizesthelikelihoodof
thetrainingsetT giventhemodelM(i.e.,P(T|M)). Forexample,supposetheword
Chinese occurs 400 

\textbf{Page 6:}

6 CHAPTER3 • N-GRAMLANGUAGEMODELS
i want to eat chinese food lunch spend
i 0.002 0.33 0 0.0036 0 0 0 0.00079
want 0.0022 0 0.66 0.0011 0.0065 0.0065 0.0054 0.0011
to 0.00083 0 0.0017 0.28 0.00083 0 0.0025 0.087
eat 0 0 0.0027 0 0.021 0.0027 0.056 0
chinese 0.0063 0 0 0 0 0.52 0.0063 0
food 0.014 0 0.014 0 0.00092 0.0037 0 0
lunch 0.0059 0 0 0 0 0.0029 0 0
spend 0.0036 0 0.0036 0 0 0 0 0
Figure3.2 BigramprobabilitiesforeightwordsintheBerkeleyRestaurant Projectcorpus
of9332sentences.Zeroprobabilit

\textbf{Page 7:}

3.2 • EVALUATINGLANGUAGEMODELS: TRAININGANDTESTSETS 7
Longer context Although for pedagogical purposes we have only described bi-
trigram grammodels,whenthereissufficienttrainingdataweusetrigrammodels,which
4-gram conditionontheprevioustwowords,or4-gramor5-grammodels. Fortheselarger
5-gram n-grams,we’llneedtoassumeextracontextstotheleftandrightofthesentenceend.
Forexample,tocomputetrigramprobabilitiesattheverybeginningofthesentence,
weusetwopseudo-wordsforthefirsttrigram(i.e.,P(I|<s><s>).
Some l

\textbf{Page 8:}

8 CHAPTER3 • N-GRAMLANGUAGEMODELS
The training set is the data we use to learn the parameters of our model; for
simple n-gram language models it’s the corpus from which we get the counts that
wenormalizeintotheprobabilitiesofthen-gramlanguagemodel.
Thetestsetisadifferent,held-outsetofdata,notoverlappingwiththetraining
set, that we use to evaluate the model. We need a separate test set to give us an
unbiasedestimateofhowwellthemodelwetrainedcangeneralizewhenweapply
ittosomenewunknowndataset. Amac

\textbf{Page 9:}

3.3 • EVALUATINGLANGUAGEMODELS: PERPLEXITY 9
to measure a statistically significant difference between two potential models. It’s
importantthatthedevsetbedrawnfromthesamekindoftextasthetestset, since
itsgoalistomeasurehowwewoulddoonthetestset.
3.3 Evaluating Language Models: Perplexity
We said above that we evaluate language models based on which one assigns a
higher probability to the test set. A better model is better at predicting upcoming
words, andso itwill be lesssurprised by(i.e., assign 

\textbf{Page 10:}

10 CHAPTER3 • N-GRAMLANGUAGEMODELS
The perplexity ofW computed with a bigram language model is still a geometric
mean,butnowoftheinverseofthebigramprobabilities:
(cid:118)
(cid:117) N
(cid:117)(cid:89) 1
perplexity(W) = (cid:116)N (3.17)
P(w|w )
i i−1
i=1
What we generally use for word sequence in Eq. 3.15 or Eq. 3.17 is the entire
sequence of words in some test set. Since this sequence will cross many sentence
boundaries,ifourvocabularyincludesabetween-sentencetoken<EOS>orseparate
begin- and en

\textbf{Page 11:}

3.4 • SAMPLINGSENTENCESFROMALANGUAGEMODEL 11
languagethatisdeterministic(noprobabilities),anywordcanfollowanyword,and
whosevocabularyconsistsofonlythreecolors:
L={red,blue,green} (3.18)
Thebranchingfactorofthislanguageis3.
Nowlet’smakeaprobabilisticversionofthesameLM,let’scallitA,whereeach
wordfollowseachotherwithequalprobability 1 (itwastrainedonatrainingsetwith
3
equalcountsforthe3colors),andatestsetT =“red red red red blue”.
Let’sfirstconvinceourselvesthatifwecomputetheperplexityofthisartific

\textbf{Page 12:}

12 CHAPTER3 • N-GRAMLANGUAGEMODELS
polyphonic
however p=0.0000018
the of a to in …(p=0.0003)
0.06 0.03 0.02 0.020.02
… …
.06 .09.11.13.15 .66 .99
0 1
Figure3.3 Avisualizationofthesamplingdistributionforsamplingsentencesbyrepeat-
edlysamplingunigrams. Thebluebarrepresentstherelativefrequencyofeachword(we’ve
orderedthemfrommostfrequenttoleastfrequent,butthechoiceoforderisarbitrary). The
numberlineshowsthecumulativeprobabilities. Ifwechoosearandomnumberbetween0
and1,itwillfallinanintervalcorrespond

\textbf{Page 13:}

3.5 • GENERALIZINGVS. OVERFITTINGTHETRAININGSET 13
–To him swallowed confess hear both. Which. Of save on trail for are ay device and
1
rotelifehave
gram –Hillhelatespeaks;or! amoretoleglessfirstyouenter
–Whydoststandforththycanopy,forsooth;heisthispalpablehittheKingHenry. Live
2
king. Follow.
gram –Whatmeans,sir. Iconfessshe? thenallsorts,heistrim,captain.
–Fly,andwillridmethesenewsofprice. Thereforethesadnessofparting,astheysay,
3
’tisdone.
gram –Thisshallforbiditshouldbebranded,ifrenownmadeit

\textbf{Page 14:}

14 CHAPTER3 • N-GRAMLANGUAGEMODELS
(3.22) Boredafdenmyphonefinnadie!!!
whiletweetsfromEnglish-basedlanguageslikeNigerianPidginhavemarkedlydif-
ferentvocabularyandn-grampatternsfromAmericanEnglish(Jurgensetal.,2017):
(3.23) @usernameRuawizardorwatgansef: indmornin-utweet,afternoon-u
tweet,nytganudeytweet. betageturITplacementwivtwitter
Isitpossibleforthetestsetnonethelesstohaveawordwehaveneverseenbe-
fore?WhathappensifthewordJurafskyneveroccursinourtrainingset,butpopsup
inthetestset? Theanswerist

\textbf{Page 15:}

3.6 • SMOOTHING,INTERPOLATION,ANDBACKOFF 15
ofthewordw isitscountc normalizedbythetotalnumberofwordtokensN:
i i
c
i
P(w)=
i
N
Laplace smoothing merely adds one to each count (hence its alternate name add-
add-one one smoothing). Since there areV words in the vocabulary and each one was in-
cremented,wealsoneedtoadjustthedenominatortotakeintoaccounttheextraV
observations. (WhathappenstoourPvaluesifwedon’tincreasethedenominator?)
c +1
i
P Laplace (w i )= (3.24)
N+V
Nowthatwehavetheintuitionfortheu

\textbf{Page 16:}

16 CHAPTER3 • N-GRAMLANGUAGEMODELS
i want to eat chinese food lunch spend
i 0.0015 0.21 0.00025 0.0025 0.00025 0.00025 0.00025 0.00075
want 0.0013 0.00042 0.26 0.00084 0.0029 0.0029 0.0025 0.00084
to 0.00078 0.00026 0.0013 0.18 0.00078 0.00026 0.0018 0.055
eat 0.00046 0.00046 0.0014 0.00046 0.0078 0.0014 0.02 0.00046
chinese 0.0012 0.00062 0.00062 0.00062 0.00062 0.052 0.0012 0.00062
food 0.0063 0.00039 0.0063 0.00039 0.00079 0.002 0.00039 0.00039
lunch 0.0017 0.00056 0.00056 0.00056 0.00056 0.0

\textbf{Page 17:}

3.6 • SMOOTHING,INTERPOLATION,ANDBACKOFF 17
language modeling, generating counts with poor variances and often inappropriate
discounts(GaleandChurch,1994).
3.6.3 LanguageModelInterpolation
There is an alternative source of knowledge we can draw on to solve the problem
of zero frequency n-grams. If we are trying to compute P(w |w w ) but we
n n−2 n−1
have no examples of a particular trigram w w w , we can instead estimate its
n−2 n−1 n
probabilitybyusingthebigramprobabilityP(w |w ). Similarly,ifw

\textbf{Page 18:}

18 CHAPTER3 • N-GRAMLANGUAGEMODELS
haszero counts, weapproximate itbybacking offtothe (n-1)-gram. Wecontinue
backing off until we reach a history that has some counts. For a backoff model to
discount giveacorrectprobabilitydistribution,wehavetodiscountthehigher-ordern-grams
to save some probability mass for the lower order n-grams. In practice, instead of
discounting, it’scommontouseamuchsimplernon-discountedbackoffalgorithm
stupidbackoff calledstupidbackoff(Brantsetal.,2007).
Stupid backoff giv

\textbf{Page 19:}

3.7 • ADVANCED: PERPLEXITY’SRELATIONTOENTROPY 19
Horse1 1 Horse5 1
2 64
Horse2 1 Horse6 1
4 64
Horse3 1 Horse7 1
8 64
Horse4 1 Horse8 1
16 64
The entropy of the random variable X that ranges over horses gives us a lower
boundonthenumberofbitsandis
i=8
(cid:88)
H(X) = − p(i)log p(i)
2
i=1
= −1
2
log22 1−1
4
log2 1
4
−1
8
log2 1
8
−
1
1
6
log21 1
6
−4(
6
1
4
log26 1
4
)
= 2bits (3.33)
A code that averages 2 bits per race can be built with short encodings for more
probable horses, and longer encodi

\textbf{Page 20:}

20 CHAPTER3 • N-GRAMLANGUAGEMODELS
TheShannon-McMillan-Breimantheorem(AlgoetandCover1988,CoverandThomas
1991)statesthatifthelanguageisregularincertainways(tobeexact,ifitisboth
stationaryandergodic),
1
H(L)= lim− logp(w 1:n ) (3.38)
n→∞ n
Thatis,wecantakeasinglesequencethatislongenoughinsteadofsummingover
all possible sequences. The intuition of the Shannon-McMillan-Breiman theorem
is that a long-enough sequence of words will contain in it many other shorter se-
quencesandthateachoftheseshorterse

\textbf{Page 21:}

3.8 • SUMMARY 21
lowercross-entropy. (Thecross-entropycanneverbelowerthanthetrueentropy,so
amodelcannoterrbyunderestimatingthetrueentropy.)
We are finally ready to see the relation between perplexity and cross-entropy
as we saw it in Eq. 3.40. Cross-entropy is defined in the limit as the length of the
observed word sequence goes to infinity. We approximate this cross-entropy by
relyingona(sufficientlylong)sequenceoffixedlength. Thisapproximationtothe
cross-entropyofamodelM=P(w|w )onasequenceofwo

\textbf{Page 22:}

22 CHAPTER3 • N-GRAMLANGUAGEMODELS
trigram probability that a given letter would be a vowel given the previous one or
twoletters. Shannon(1948)appliedn-gramstocomputeapproximationstoEnglish
wordsequences.BasedonShannon’swork,Markovmodelswerecommonlyusedin
engineering,linguistic,andpsychologicalworkonmodelingwordsequencesbythe
1950s. InaseriesofextremelyinfluentialpapersstartingwithChomsky(1956)and
includingChomsky(1957)andMillerandChomsky(1963),NoamChomskyargued
that “finite-state Markov process

\textbf{Page 23:}

EXERCISES 23
Exercises
3.1 Writeouttheequationfortrigramprobabilityestimation(modifyingEq.3.11).
Nowwriteoutallthenon-zerotrigramprobabilitiesfortheI am Samcorpus
onpage4.
3.2 Calculatetheprobabilityofthesentencei want chinese food. Givetwo
probabilities,oneusingFig.3.2andthe‘usefulprobabilities’justbelowiton
page6, andanotherusingtheadd-1smoothedtableinFig.3.7. Assumethe
additionaladd-1smoothedprobabilitiesP(i|<s>)=0.19andP(</s>|food)=
0.40.
3.3 Whichofthetwoprobabilitiesyoucomputedinthepreviou

\textbf{Page 24:}

24 CHAPTER3 • N-GRAMLANGUAGEMODELS
3.9 Runyourn-gramprogramontwodifferentsmallcorporaofyourchoice(you
might use email text or newsgroups). Now compare the statistics of the two
corpora. Whatarethedifferencesinthemostcommonunigramsbetweenthe
two? Howaboutinterestingdifferencesinbigrams?
3.10 Addanoptiontoyourprogramtogeneraterandomsentences.
3.11 Addanoptiontoyourprogramtocomputetheperplexityofatestset.
3.12 You are given a training set of 100 numbers that consists of 91 zeros and 1
eachoftheothe

\textbf{Page 25:}

Exercises 25
Algoet,P.H.andT.M.Cover.1988. Asandwichproofof Jelinek,F.andR.L.Mercer.1980. Interpolatedestimation
theShannon-McMillan-Breimantheorem. TheAnnalsof ofMarkovsourceparametersfromsparsedata. InE.S.
Probability,16(2):899–909. GelsemaandL.N.Kanal, eds, Proceedings, Workshop
Bahl, L.R., F.Jelinek, andR.L.Mercer.1983. Amaxi- onPatternRecognitioninPractice,381–397.NorthHol-
mumlikelihoodapproachtocontinuousspeechrecogni- land.
tion.IEEETransactionsonPatternAnalysisandMachine Jelinek,F.,R.L.

\textbf{Page 26:}

26 Chapter3 • N-gramLanguageModels
Stolcke,A.2002.SRILM–anextensiblelanguagemodeling
toolkit.ICSLP.
Trnka, K., D. Yarrington, J. McCaw, K. F. McCoy, and
C.Pennington.2007. Theeffectsofwordpredictionon
communicationrateforAAC.NAACL-HLT.
Witten,I.H.andT.C.Bell.1991.Thezero-frequencyprob-
lem:Estimatingtheprobabilitiesofnoveleventsinadap-
tivetextcompression.IEEETransactionsonInformation
Theory,37(4):1085–1094.


\subsection{Module 1 Notes}

\begin{itemize}
    \item Introduction: 
    \item Generic Natural Language Processing (NLP) system [CO1]
    \item NLP is the field that enables computers to understand and respond to human language. It is crucial for applications like email filtering (e.g., spam detection). NLP has evolved from traditional rule-based systems to using machine learning and deep learning techniques, especially from the 1980s onward. Natural language processing (NLP) is a branch of artificial intelligence that helps computers understand, interpret and manipulate human language. ​
    \item ​
    \item Natural Language Processing, usually shortened as NLP, is a branch of artificial intelligence that deals with the interaction between computers and humans using the natural language. ​
    \item ​
    \item The ultimate objective of NLP is to read, decipher, understand, and make sense of the human languages in a manner that is valuable. It is an interdisciplinary subfield of linguistics, computer science, and artificial intelligence concerned with the interactions between computers and human language​.
    \item  ​
    \item Text Analytics and NLP 
    \item Text analytics focuses on extracting meaningful insights from text data, which may not always be human language. In contrast, NLP includes both text and speech recognition. The two branches of NLP are: - 
    \item Natural Language Understanding (NLU): Comprehension of spoken or written language by machines. 
    \item Natural Language Generation (NLG): Machines producing human-understandable text.
    \item Stages in NLP, [CO1]
    \item To converse with human a program must understand the syntax(grammar), semantics (word meaning) and morphology (word level analysis), and pragmatics(conversation). ​
    \item There are certain phases in which Natural language processing is performed.
    \item  ​
    \item 1. Morphological Analysis
    \item It deals with morphemes which are the smallest units of meaning in a word. It is important for understanding structure of words and their parts by identifying free morphemes (independent words like "cat") and bound morphemes (like prefixes or suffixes e.g. "un-" or "-ing").
    \item Key tasks in morphological analysis:
    \item Stemming: Reducing words to their root form like "running" to "run".
\end{itemize}


\section{Module 2: Semantic Analysis and Word Embeddings}


\subsection{1 Vector Semantics _Final}

\textbf{Slide 1:}

Computational Semantics and Semantic Parsing

\textbf{Slide 2:}

What Is Semantic Analysis in NLP?

To understand the meaning of Natural Language

Captures the meaning of the given text while considering context, logical structuring of sentences and grammar roles

Due to the vast complexity and subjectivity involved in human language, interpreting it is quite a complicated task for machines

\textbf{Slide 3:}

What is Semantics ? Word

A single distinct meaningful element of speech or writing, used with others to form a sentence.

Words that occur in similar contexts tend to have similar meanings.

Words which are synonyms (like oculist and eye-doctor) tended to occur in the same environment.

\textbf{Slide 8:}

\textbf{Vector Semantics:} 

It represent a word as a point in some multidimensional semantic space.

Vectors for representing words are generally called embeddings, because the word is embedded in a particular vector space.

Vector semantic models can be learned automatically from text without any complex labeling or supervision.

As a result of these advantages, vector models is now a standard way to represent the meaning of words in NLP.

\textbf{Slide 10:}

Vector Space

\textbf{Slide 11:}

Word Space

\textbf{Target words:}  automobile, car, soccer, football

\textbf{Term vocabulary:}  wheel, transport, passenger, tournament, London, goal, match

\textbf{Slide 12:}

Constructing Word spaces

Informal algorithm for constructing word spaces

\textbf{Pick the words you are interested in:}   target words

Define a context window, number of words surrounding target word

The context can in general be defined in terms of documents, paragraphs or sentences.

Count number of times the target word co-occurs with the context words: co-occurrence matrix

Build vectors out of (a function of) these co-occurrence counts

\textbf{Slide 13:}

\textbf{Constructing Word spaces:}  distributional vectors

\textbf{Slide 14:}

Vector Space

\textbf{Slide 15:}

Computing similarity

\textbf{Slide 16:}

Vector Space Model without distributional similarity

Words are treated as atomic symbols

One-hot representation

\textbf{Slide 17:}

Building a DSM step-by-step

\textbf{Slide 18:}

Building a DSM step-by-step

\textbf{Slide 19:}

Many design choices

\textbf{Slide 20:}

The parameter space

\textbf{Slide 21:}

\textbf{Documents as context:}  Word × document

\textbf{Slide 22:}

\textbf{Words as context:}  Word × Word

\textbf{Slide 24:}

fool

\textbf{Slide 25:}

\textbf{Context weighting:}  documents as context

\textbf{Slide 27:}

Bag of Words

\textbf{Slide 28:}

Bag of words Model

Text modelling technique .

Convert a text into its equivalent vector of numbers.

It is a method of feature extraction with text data.

\textbf{Slide 29:}

Understanding Bag of Words with an example

\textbf{Example(1) without preprocessing:} 

\textbf{Sentence 1:}   ”Welcome to Great Learning, Now start learning”

\textbf{Sentence 2:}  “Learning is a good practice”

\textbf{Slide 30:}

\textbf{Solution:} 

\textbf{Example:} 

\textbf{Sentence 1:}   ”Welcome to Great Learning, Now start learning”​

\textbf{Sentence 2:}  “Learning is a good practice”

Bag of Words with an example

\textbf{Slide 31:}

Step 1: Go through all the words in the above text and make a list of all the words in our model vocabulary.

Welcome

To

Great

Learning

,

Now

start

learning

is

a

good

practice

\textbf{Slide 32:}

Frequency Table for both sentences

writing the above frequencies in the vector form

Sentence 1 ➝ [ 1,1,1,1,1,1,1,1,0,0,0 ]

Sentence 2 ➝ [ 0,0,0,0,0,0,0,1,1,1,1,1 ]

\textbf{Slide 33:}

\textbf{Example(2) with preprocessing:} 

\textbf{Example(2) with preprocessing:} 

\textbf{Sentence 1:}  ”Welcome to Great Learning, Now start learning”

\textbf{Sentence 2:}  “Learning is a good practice”

Step 1: Convert the above sentences in lower case as the case of the word does not hold any information.

Step 2: Remove special characters and stopwords from the text. Stopwords are the words that do not contain much information about text like ‘is’, ‘a’,’the and many more’.

\textbf{Slide 34:}

Occurrence matrix

Writing the above frequencies in the vector

Sentence 1 ➝ [ 1,1,2,1,1,0,0 ]

Sentence 2 ➝ [ 0,0,1,0,0,1,1 ]

\textbf{Slide 35:}

Disadvantages of bag of Words

Treats all words equally → Common words (e.g., the, is, and) get same importance as meaningful words.

No importance weighting → Cannot highlight rare but informative words.

High frequency bias → Words appearing many times dominate even if not meaningful.

No document-level context → Doesn’t consider how unique a word is across documents.

Sparse and noisy features → Large vocabulary with many low-value words.

\textbf{Slide 36:}

Disadvantages of bag of Words

\textbf{Slide 40:}

Practice Problem

\textbf{Compute TFIDF value for word Cat given in a small corpus with three documents:} 

\textbf{Document 1:}  "The cat sat on the mat"

\textbf{Document 2:}  "The dog sat on the log"

\textbf{Document 3:}  "The cat and the dog played"

\textbf{Slide 41:}

Compute Term Frequency (TF)

\textbf{The Term Frequency (TF) is calculated as:} 

TF = Number of times term appears in document / Total terms in document

\textbf{Let’s compute TF for the word "cat" in each document:} 

Document 1: "The cat sat on the mat"Total words = 6, Occurrences of "cat" = 1TF(cat, D1) = 1/6 = 0.1667

Document 2: "The dog sat on the log"Total words = 6, Occurrences of "cat" = 0TF(cat, D2) = 0/6 = 0

Document 3: "The cat and the dog played"Total words = 7, Occurrences of "cat" = 1TF(cat, D3) = 1/7 = 0.1429

\textbf{Slide 42:}

\textbf{Step 3:}  Compute Inverse Document Frequency (IDF)

\textbf{The Inverse Document Frequency (IDF) is calculated as:} 

IDF = log ( Total number of documents / Number of documents containing the term )

\textbf{For the term "cat":} 

The total number of documents = 3

Documents containing "cat" = 2 (D1 and D3)

IDF(cat) = log (3/2) = log(1.5) ≈ 0.1761

\textbf{Slide 43:}

Compute TF-IDF Score

\textbf{The TF-IDF score is obtained by multiplying TF and IDF:} 

TFIDF = TF × IDF

\textbf{For "cat":} 

\textbf{Document 1:} TFIDF(cat, D1) = 0.1667 × 0.1761 = 0.0294

\textbf{Document 2:} TFIDF(cat, D2) = 0 × 0.1761 = 0

\textbf{Document 3:} TFIDF(cat, D3) = 0.1429 × 0.1761 = 0.0251

Final TF-IDF Scores for "cat"

\textbf{Slide 44:}

\textbf{Example :} TFIDF

News Article Classification Problem

\textbf{Slide 45:}

\textbf{Example :} TFIDF

News Article Classification Problem

\textbf{Slide 46:}

Consider The Document Frequency

\textbf{Inverse Document Frequency :} -

\textbf{Slide 47:}

TF-IDF

\textbf{Slide 48:}

Why Log of IDF

\textbf{Slide 49:}

Numerical TF-IDF

\textbf{Consider 3 documents:} 

\textbf{D1:}  "data science is fun"

\textbf{D2:}  "machine learning is fun"

\textbf{D3:}  "data mining and machine learning"

\textbf{Step 1:}  Vocabulary (Unique Terms)

data, science, is, fun, machine, learning, mining, and

\textbf{Slide 50:}

\textbf{Step 2:}  Term Frequency (TF)

TF = (Number of times term appears in document) / (Total terms in that document)

\textbf{Document D1:}  "data science is fun"

\textbf{Document D2:}  "machine learning is fun"

\textbf{Document D3:}  "data mining and machine learning"

\textbf{Slide 51:}

\textbf{Step 3:}  Document Frequency (DF)

\textbf{Slide 53:}

TF–IDF for D2

TF–IDF for D3

\textbf{Slide 54:}

Final TF–IDF Vector Form

\textbf{Slide 55:}

Point Mutual Information

An alternative weighting function to tf-idf

Ask whether a context word is particularly informative about the target word.

Pointwise mutual information is one of the most important concepts in NLP.

It is a measure of how often two events x and y occur, compared with what we would expect if they were independent

\textbf{Slide 56:}

Do events x and y co-occur more than if they were independent?

PMI(X, Y) = log₂ ( P(x, y) / (P(x) P(y)) )

\textbf{PMI between two words:}  (Church \& Hanks 1989)

Do words word₁ and word₂ co-occur more than if they were independent?

PMI(word₁, word₂) = log₂ ( P(word₁, word₂) / (P(word₁) P(word₂)) )

\textbf{Pointwise mutual information:} 

\textbf{Slide 57:}

Computing PPMI on a term-context matrix

Matrix F with W rows (words) and C columns (contexts)

fij is \# of times wi occurs in context cj

57

\textbf{Slide 58:}

p(w=information , c=data) =

p(w=information) =

p(c=data) =

58

= .3399

3982/111716

7703/11716

= .6575

5673/11716

= .4842

\textbf{Slide 59:}

59

PMI(information,data) =

Resulting PPMI matrix (negatives replaced by 0)

\textbf{Slide 60:}

Weighting PMI

PMI is biased toward infrequent events

Very rare words have very high PMI values

\textbf{Two solutions:} 

Give rare words slightly higher probabilities

Use add-one smoothing (which has a similar effect)

60

\textbf{Slide 61:}

\textbf{Weighting PMI:}  Giving rare context words slightly higher probability

61

\textbf{Slide 62:}

p(w=information,c=data) =  6/19 p(w=information) =  11/19  = .58 p(c=data) = 7/19  = .37

p(w,context)

p(w)

= .32

ij

p	

fij

W	C

 fij

i1 j1

i

p(w ) 

C

 fij

j1

N

p(cj )   i1

W

 fij

N

\textbf{Slide 63:}

ij

pmi	 log

2

pij

p	p

i*	* j

pmi(information,data) = log2 ( .32 /	(.37*.58) ) = .57

PPMI(w,context)

p(w,context)

p(w)

\textbf{Slide 64:}

Add-2 Smoothed Count

(w,context

p(w,context) [add02]

p(w)

\textbf{Slide 65:}

PPMI versus add-2 smoothed PPMI

PPMI(w,context)

\textbf{Slide 66:}

Numerical Example: Consider the following word-context frequencies from a corpus:

\textbf{Slide 68:}

PPMI and Rare Words

PPMI tends to favor rare co-occurrences, which means it gives higher scores to words that rarely occur but co-occur with certain context words more often than expected by chance.

\textbf{Example:} 

\textbf{Let’s say we have the following corpus:} 

markdown

1. The dog barked at the cat.

2. The cat sat on the mat.

3. The dog and the cat watched the **unicorn**.

\textbf{"dog", "cat":}  These words are common and occur in several sentences.

\textbf{"unicorn":}  This word is rare—it appears only once, in a single unique sentence.

\textbf{Slide 69:}

\textbf{PPMI Calculation:} ​

Even though "unicorn" is rare, it gets a high PPMI score with the context "cat" because it occurs with "cat" more frequently than expected (given the general rarity of "unicorn").​

This is useful in some NLP tasks (e.g., word similarity) because rare but important relationships are highlighted.​

\textbf{PPMI Example Value:} ​

PPMI("unicorn", "cat")=3.5(High value)

PPMI("dog", "cat")=0.2(Low value due to frequent co-occurrence)

\textbf{Slide 70:}

Weighting PMI

PMI is biased toward infrequent events

Very rare words have very high PMI values

\textbf{Two solutions:} 

Give rare words slightly higher probabilities

Use add-one smoothing (which has a similar effect)

\textbf{Slide 71:}

How PPMI is Different from TF-IDF

\textbf{Slide 72:}

\textbf{Computing word similarity:}  Dot product and cosine

\textbf{The dot product between two vectors is a scalar:} 

The dot product tends to be high when the two vectors have large values in the same dimensions

Dot product can thus be a useful similarity metric between vectors

\textbf{Slide 73:}

Problem with raw dot-product

Dot product favors long vectors

Dot product is higher if a vector is longer (has higher values in many dimension)

\textbf{Vector length:} 

Frequent words (of, the, you) have long vectors (since they occur many times with other words).

So dot product overly favors frequent words

\textbf{Slide 74:}

\textbf{Alternative:}  cosine for computing word similarity

Based on the definition of the dot product between two vectors a and b

\textbf{Slide 75:}

Cosine as a similarity metric

\textbf{-1:}  vectors point in opposite directions

\textbf{+1:}   vectors point in same directions

\textbf{0:}  vectors are orthogonal

But since raw frequency values are non-negative, the cosine for term-term matrix vectors ranges from 0–1

75

\textbf{Slide 76:}

Cosine examples

76

\textbf{Slide 77:}

Visualizing cosines (well, angles)

\textbf{Slide 78:}

Word2Vec

Word Embedding is a language modeling technique used for mapping words to vectors of real numbers.

It represents words or phrases in vector space with several dimensions.

Word embeddings can be generated using various methods like neural networks, co-occurrence matrix, probabilistic models, etc.

\textbf{Slide 79:}

Word2Vec

Word2Vec creates vectors of the words that are distributed numerical representations of word features

these word features could comprise of words that represent the context of the individual words present in our vocabulary.

Word embeddings eventually help in establishing the association of a word with another similar meaning word through the created vectors.

\textbf{Slide 80:}

Word2Vec

Word2Vec consists of models for generating word embedding.

These models are shallow two-layer neural networks having one input layer, one hidden layer, and one output layer.

\textbf{Word2Vec utilizes two architectures :} 

CBOW (Continuous Bag of Words)

Skip Gram

\textbf{Slide 82:}

Issues in BOW /TF\_IDF

Semantics meaning is not captured

Similar words should have similar kind of vectors

Sparse Matrix is generated

Increase no of dimensions

Vectors created in word2Vec

Limited dimension

Sparsity is reduce ( Create dense vector )

Semantic Meaning is maintained

\textbf{Slide 84:}

How does CBOW work?

Word2Vec is an unsupervised model

you can give a corpus without any label information and the model can create dense word embeddings,

Word2Vec internally leverages a supervised classification model to get these embeddings from the corpus.

\textbf{Slide 85:}

How does CBOW work

The CBOW architecture comprises a deep learning classification model in which we take in context words as input, X, and try to predict our target word, Y.

\textbf{Slide 87:}

Architecture of CBOW

\textbf{Slide 88:}

Math Behind The Above Calculation

\textbf{Slide 89:}

Skip-Gram Model

The skip-gram model is a simple neural network with one hidden layer trained in order to predict the probability of a given word being present when an input word is present.

Intuitively, you can imagine the skip-gram model being the opposite of the CBOW model.

In this architecture, it takes the current word as an input and tries to accurately predict the words before and after this current word.

This model essentially tries to learn and predict the context words around the specified input word.


\subsection{2 SemanticAnalysisPart2}

\textbf{Slide 1:}

Semantic Analysis

---

Lexical Semantics

\textbf{Slide 2:}

\textbf{WORDNET:} 

WordNet is a database of lexical relations for English (Fellbaum, 1998).

\textbf{WordNet consists of Three databases:}  Noun, Verb, Adjective \& Adverb.

\textbf{Slide 3:}

\textbf{WORDNET:} 

A hierarchically organized lexical database

• On-line thesaurus + aspects of a dictionary

• Word senses and sense relations

• Some other languages available or under development(Arabic, Finnish, German, Portuguese…)

\textbf{Slide 6:}

SYNSETS

Different words with similar meaning is a set of synonms representing a sense

\textbf{Example:} 

Chump🡪 person who is gullible & easy to take advantage of

{Chump1 , fool 2, gull 1, guy 1, mark 9}

9th sense of mark is this meaning

1st sense of guy is same as 2nd sense of fool

\textbf{Slide 7:}

2

Relationships between word  meanings

Homonymy

Polysemy

Synonymy

Antonymy

Hypernomy

Hyponomy

Meronomy

\textbf{Slide 8:}

lemma and wordform

Semantic Analysis

8

\textbf{Slide 9:}

Lemmas have senses

\textbf{One lemma “bank” can have many meanings:} 

\textbf{Sense 1:} 

• ...a bank can hold the investments in a custodial 1

Account...

\textbf{Sense 2:} 

“...as agriculture burgeons on the east bank the river will shrink even more” 2

• Sense (or word sense)

• A discrete representationof an aspect of a word’s meaning.

• The lemma bank here has two senses 14

Semantic Analysis

9

\textbf{Slide 10:}

Lemma Pepper

\textbf{Sense 1:}  spice from pepper plant

\textbf{Sense 2:}  the pepper plant itself

\textbf{Sense 3:}  another similar plant (Jamaican pepper)

\textbf{Sense 4:}  another plant with peppercorns  (California pepper)

\textbf{Sense 5:}  capsicum (i.e. chili, paprika, bell pepper,  etc)

A sense or word sense is discrete representation of one aspect of  meaning of word.

\textbf{Slide 11:}

Word Senses

A sense is discrete representation of one aspect of meaning of  word.

A sense or “concept” is the meaning component of a word

There are relations between senses

Semantic Analysis

9

\textbf{Slide 12:}

Although a form can have more than one meaning,  there is always a primary or original meaning that  is expresses.

The original meaning of a form in a language is normally

called “denotation”.

With respect to words (lexical terms), the primary  (original) meaning is the meaning that we can find in a  dictionary.

\textbf{Example:}  “Rose” “A garden plant with thorns on its stem

and pleasant-smelling flowers, or a flower from this plant”.

⚫

Semantic Analysis

12

“Denotation” (Primary Dictionary Meaning)

\textbf{Slide 13:}

Besides, there are additional meaning, which is known as “connotation”.

With respect to words, connotation or additional meanings of a  word are not listed in a dictionary; and therefore, we cannot find  them in a dictionary

Normally, a person expresses a connotative meaning through a  word, phrase, clause or sentence based on certain  characteristics of the entity or event that he/she is referring to.

\textbf{Example:}  one may call a beautiful girl a “rose” or a “lily”.

13

Connotation

\textbf{Slide 14:}

HOMONYMY

Sense relation in which one form has different meaning; Same Spelling has different Meaning (treated as  such in dictionaries)

no relatedness in meaning

e.g. bank1-side of a river

bank2-financial institution

Semantic Analysis

14

\textbf{Slide 15:}

Semantic Analysis

15

\textbf{Slide 16:}

Semantic Analysis

16

\textbf{Slide 17:}

POLYSEMY

One word (Same Spelling)  having several closely related senses

Native speaker has clear intuition that the  different senses are related to each other

Polysemy-close relatedness in m. which is usually  connected to metaphorical extension

meaning relatedness has to be synchronic

Semantic Analysis

17

\textbf{Slide 18:}

The following sentences contain some examples of polysemy.

He drank a glass of milk.

He forgot to milk the cow.

His cottage is near a small wood.

The statue was made out of a block of wood.

He fixed his hair.

They fixed a date for the wedding.

Semantic Analysis

18

\textbf{Slide 19:}

25

Polysemy

The bank is constructed from red brick  I withdrew the money from the bank

Are those the same sense?

Which sense of bank is this?

Is it distinct from (homonymous with) the river bank sense?

How about the savings bank sense?

A single lexeme with multiple related meanings (bank the building,  bank the finantial institution)

\textbf{Slide 20:}

Semantic Analysis

20

\textbf{Slide 21:}

Polysemy is the state or phenomenon in which the words  that have more than one meaning.

In other words, it can  be described as multiple meanings of words.The words  are considered to be related etymologically.

The concrete form of polysemy is called “polyseme”.

\textbf{Examples:}  Simple (e.g. English is extremely plain  subject).

Plain With nothing added/not decorated in any way  (e.g.This blouse is too plain)

Semantic Analysis

21

\textbf{Slide 22:}

HOMONYMY-POLYSEMY DISTINCTION

Homonymy-polysemy distinction: closeness or  relatedness of senses of ambiguous w.

Homonymy- different senses of an ambiguous

are far apart from each other and not related via

speaker’s intuition;

matter of accident or coincidence e.g. bank, mug; kit, bar

Semantic Analysis

22

\textbf{Slide 23:}

Semantic Analysis

23

\textbf{Slide 24:}

\textbf{Relation:}  Synonymy

Synonyms have the same meaning in some or all  contexts.

filbert / hazelnut

couch / sofa

big / large

automobile / car

vomit / throw up

Water / H20

\textbf{Slide 25:}

\textbf{Relation:}  Synonymity

Note that there are probably no examples of  perfect synonymy.

Even if many aspects of meaning are identical

Still may not preserve the acceptability based on  notions of politeness, slang, register, genre, etc.

\textbf{The Linguistic Principle of Contrast:} 

Difference in form -> difference in meaning

\textbf{Slide 26:}

Synonymy is a relation

between senses rather than words

Consider the words big and large

•  Are they synonyms?

How big is that plane?

Would I be flying on a large or small plane?

\textbf{•	How about here:} 

Miss Nelson became a kind of big sister to Benjamin.

?Miss Nelson became a kind of large sister to Benjamin.

• Why?

big has a sense that means being older, or grown up • large

lacks this sense

Semantic Analysis

26

\textbf{Slide 27:}

Hyponymy and Hypernymy

One sense is a hyponym of another if the first sense is more  specific, deno(ng a subclass of the other

• car is a hyponym of vehicle

• mango is a hyponym of fruit

Conversely hypernym/superordinate (“hyper is super”)

vehicle is a hypernym of car

fruit is a hypernym of mango

Semantic Analysis

27

\textbf{Slide 28:}

Hyponymy (i.e. category membership)

\textbf{It may be problematic to identify the superordinate terms:} 

brother & sister < sibling (formal)

uncle & aunt < ?

cow & bull < cow/cattle (collective)/bovine (technical)

human being & animal < animal (vs. vegetable, mineral)

fish

snapper

trout	bass [bæs]	sole

salmon [ˈsæm ən]

chinook [(ˌ)tʃɪ ˈnuːk]

spring  coho [ˈkəʊ həʊ]

king

sockeye [ˈsɒk aɪ]

hypernym

(co)hyponyms

\textbf{Slide 29:}

Hyponymy more formally

\textbf{Extensional:} 

The class denoted by the superordinate extensionally

includes the class denoted by the hyponym

\textbf{Entailment:} 

A sense A is a hyponym of sense B if being an A entails being a B •

Hyponymy is usually transitive

(A hypo B and B hypo C entails A hypo C)

\textbf{• Another name:}  the IS-A hierarchy

A IS-A B	(or A ISA B)

B subsumes A

Semantic Analysis

43

\textbf{Slide 30:}

Antonyms

Senses that are opposites with respect to one feature of  meaning

Otherwise, they are very similar!

dark/light short/long fast/slow rise/fall hot/cold up/down  in/out

\textbf{• More formally:}  antonyms can • define a binary opposition or

be at opposite ends of a scale

• long/short, fast/slow

\textbf{• Be reversives:} 

rise/fall,up/down

Semantic Analysis

44

\textbf{Slide 31:}

Metaphor and Metonymy

31

Specific types of polysemy

\textbf{Metaphor:} 

Germany will pull Slovenia out of its economic slump.(compares Slovenia's economic difficulties to a physical "slump" or pit,)

I spent 2 hours on that homework.(time being treated as money or a resource)

Metonymy	-the act of referring to something by the name of  something else that is closely connected with it, for example using  the White House for the US President

The White House announced yesterday.

This chapter talks about part-of-speech tagging

Bank (building) and bank (financial institution)

\textbf{Slide 32:}

How do we know when a word has  more than one sense?

32

ATIS examples

Which flights serve breakfast?

Does America West serve Philadelphia?

\textbf{The “zeugma” test:} 

?Does United serve breakfast and San Jose?

\textbf{Slide 33:}

Synonyms

33

Word that have the same meaning in some or all contexts.

hazelnut  adolescent  large

filbert

youth

big

automobile car

Two lexemes are synonyms if they can be successfully substituted for  each other in all situations

\textbf{Slide 34:}

Synonyms

34

But there are few (or no) examples of perfect synonymy.

Why should that be?

Even if many aspects of meaning are identical

Still may not preserve the acceptability based on notions of politeness, slang,  register, genre, etc.

\textbf{Example:} 

Big and large?

That’s my big sister

That’s my large sister

\textbf{Slide 35:}

Antonyms

35

light

Words that are opposites with respect to one feature of their  meaning

Otherwise, they are very similar!

Dark

Boy	girl

Hot cold

Up

In

down  out

\textbf{Slide 36:}

Word Similarity Computation

36

For various computational applications it’s useful to find words which

are similar to another word.

Machine translation (to find near-synonyms)

Information retrieval (to do “query expansion”)

\textbf{Two ways to do this:} 

Automatic computation based on distributional similarity

Lsa.colorado.edu

Use a thesaurus which lists similar words.

WordNet (we’ll return to this in a few slides)


\subsection{3 Word Sense Disambiguation}

\textbf{Slide 1:}

Word Sense Disambiguation

\textbf{Slide 3:}

Word Sense Disambiguation Problem

WSD is the problem of identifying which “sense” (meaning) of a word is activated by the use of the word in a particular context or scenario.

deals with determining the intended meaning of a word in a given context.

It is the process of identifying the correct sense of a word from a set of possible senses, based on the context in which the word appears.

\textbf{Slide 4:}

Sense ambiguity

Many words have several meanings or senses

The meaning of bass depends on the context

Are we talking about music, or fish?

An electric guitar and bass player stand off to one side, not really part of the scene, just as a sort of nod to gringo expectations perhaps.

And it all started when fishermen decided the striped bass in Lake Mead were too skinny.

Disambiguation

The task of disambiguation is to determine which of the senses of an ambiguous word is invoked in a particular use of the word.

\textbf{Slide 5:}

Application

Machine Translation

Information Retrieval

Question Answering

Text Summarization

Sentiment Analysis

Chatbots & Virtual Assistants

Search Engines

\textbf{Example:} 

Apple released a new product.

Is Apple → fruit 🍎 or company 🏢?

\textbf{Slide 6:}

Algorithms

Knowledge Based Approaches

Overlap Based Approaches

Machine Learning Based Approaches

Supervised Approaches

Semi-supervised Algorithms

Unsupervised Algorithms

Hybrid Approaches

\textbf{Slide 7:}

\textbf{Approaches:} 

• Dictionary- and knowledge-based methods: These rely primarily on dictionaries,

thesauri, and lexical knowledge bases, without using any corpus evidence.

\textbf{• Supervised methods:}  These make use of sense-annotated corpora to train from.

\textbf{• Semi-supervised or minimally-supervised methods:}  These make use of a

secondary source of knowledge such as a small annotated corpus as seed data in a

bootstrapping process, or a word-aligned bilingual corpus

\textbf{Slide 8:}

\textbf{WORDNET:} 

WordNet is a database of lexical relations for English (Fellbaum, 1998).

\textbf{WordNet consists of Three databases:}  Noun, Verb, Adjective \& Adverb.

\textbf{Slide 9:}

\textbf{WORDNET:} 

A hierarchically organized lexical database

• On-line thesaurus + aspects of a dictionary

• Word senses and sense relations

• Some other languages available or under development(Arabic, Finnish, German, Portuguese…)

\textbf{Slide 12:}

SYNSETS

Different words with similar meaning is a set of synonms representing a sense

\textbf{Example:} 

Chump🡪 person who is gullible & easy to take advantage of

{Chump1 , fool 2, gull 1, guy 1, mark 9}

9th sense of mark is this meaning

1st sense of guy is same as 2nd sense of fool

\textbf{Slide 13:}

\textbf{HIERARCHICAL SYNSET RELATIONS:}  NOUNS

Hypernym/hyponym (between concepts)

The more general ‘meal’ is a hypernym of the more specific ‘breakfast’

Instance hypernym/hyponym (between concepts and instances)

Austen is an instance hyponym of author

Member holonym/meronym (groups and members)

professor is a member meronym of (a university’s) faculty

Part holonym/meronym (wholes and parts)

wheel is a part meronym of (is a part of) car.

Substance meronym/holonym (substances and components)

flour is a substance meronym of (is made of) bread

\textbf{Slide 14:}

\textbf{HIERARCHICAL SYNSET RELATIONS:}  VERBS

\textbf{Hypernym/troponym (between events):} 

travel/fly, walk/stroll

\textbf{Flying is a troponym of traveling:} 

it denotes a specific manner of traveling

\textbf{Entailment (between events):} 

snore/sleep

Snoring entails (presupposes) sleeping

\textbf{Slide 17:}

\textbf{Dictionary-based WSD:}  Lesk algorithm

\textbf{Slide 19:}

Dictionary / Knowledge /Thesaurus Based Approaches

Lesk algorithm

Walkers algorithm

\textbf{We have a dictionary/thesaurus that contains glosses and examples:} 

bank1

Gloss: a financial institution that accepts deposits and channels the money into lending activities

Examples: “he cashed the check at the bank”, “that bank holds the mortgage on my home”

bank2

\textbf{Gloss:}  sloping land (especially the slope beside a body of water)

Examples: “they pulled the canoe up on the bank”, “he sat on the bank of the river and watched the current”

We could use anything in the thesaurus

§ Example sentences

§ Derivational relations and sentence frames

\textbf{Slide 20:}

Knowledge Based Approaches

Overlap Based Approaches

Require a Machine Readable Dictionary (MRD).

For Example WordNet (all the information should in form that machine can understand)

\textbf{Steps:} 

Take the ambiguous word.

List all possible senses from dictionary.

Compare definition words with surrounding context.

Select sense with maximum word overlap.

Find the overlap between the features of different senses of an ambiguous word (sense bag) and the features of the words in its context (context bag).

The features could be sense definitions, example sentences, hypernyms etc.

The features could also be given weights.

The sense which has the maximum overlap is selected as the contextually appropriate sense.

\textbf{Slide 21:}

Lesk’s Algorithm

Sense Bag: contains the words in the definition of a candidate sense of the ambiguous word.

Context Bag: contains the words in the definition of each sense of each context word.

\textbf{Example :} -On burning coal we get ash.

\textbf{Slide 22:}

She drank Some Milk

She drank Some books

\textbf{Slide 23:}

Walker’s Algorithm

A Thesaurus Based approach

What is Thesaurus ?

A thesaurus is a book, software program, or online service that provides alternative or similar words to a word.

For example, searching for "hope" may return synonyms like "achievement," "faith," "ambition," and "optimism.“

Thesaurus vs. dictionary

A thesaurus is used to group different words with the same meaning (synonyms) and similar words.

On the other hand, a dictionary explains the definition of a word.

For example, looking up the word "computer" in a thesaurus may give words like PC, CPU, calculator, abacus, and laptop that could be used in place of the word computer.

Looking up the word "computer" in a dictionary would define the word like what is found on our computer definition.

\textbf{Slide 24:}

Walker’s Algorithm -A Thesaurus Based approach

Steps

Identify the ambiguous word.

Extract all possible senses of the word from a dictionary / WordNet.

Identify context words surrounding the ambiguous word.

Measure semantic similarity between each sense of the ambiguous word and each context word.

Select the sense with maximum total similarity score.

\textbf{semantic similarity Calculation:} 

Step 1: For each sense of the target word find the thesaurus category to which that sense belongs

\textbf{Step 2:}  Calculate the score for each sense by using the context words.

A context word will add 1 to the score of the sense if the thesaurus

category of the word matches that of the sense.

E.g. The money in this bank fetches an interest of 8% per annum

\textbf{Target word:}  bank

\textbf{Clue words from the context:}  money, interest, annum, fetch

\textbf{Slide 25:}

\textbf{Key Differences:} 

\textbf{Slide 26:}

MRD – A Resource for Knowledge-based WSD

\textbf{For each word in the language vocabulary, an MRD provides:} 

A list of meanings

Definitions (for all word meanings)

Typical usage examples (for most word meanings)

WordNet definitions/examples for the noun plant

buildings for carrying on industrial labor; "they built a large plant to  manufacture automobiles“

a living organism lacking the power of locomotion

something planted secretly for discovery by another; "the police used a plant to

trick the thieves"; "he claimed that the evidence against him was a plant"

an actor situated in the audience whose acting is rehearsed but seems  spontaneous to the audience

\textbf{Slide 30:}

\textbf{1.Constructing the Graph:} 

Identify ambiguous words in the sentence.

Extract all their possible senses using WordNet / BabelNet.

Build a semantic graph connecting all senses.

Add weighted edges based on semantic relationships.

Run graph algorithms to compute node importance.

Choose highest-ranked sense.

E.g.

\textbf{The sentence "She drank some milk" contains two key words:}  "drink" and "milk".

These words are inserted into a graph.

Each word has multiple senses (meanings), which are represented as nodes (e.g., "drink₁", "drink₂", etc., and "milk₁", "milk₂", etc.).

Directed edges link the words to their possible senses.

Additional edges connect the senses based on their relationships, such as synonyms, hypernyms, or other lexical relations.

Edge Weights can be assigned based on Semantic similarity, Relation type, Path length

\textbf{2. Choosing the Most Central Sense:} 

The algorithm assigns some initial probability to each word ("drink" and "milk").

Using PageRank, a centrality algorithm, the most important or frequently visited sense in the graph is identified.

The sense with the highest PageRank score is chosen as the correct meaning.   (good page rank is given based on Connectivity , Edge weights, Importance of neighboring nodes)

\textbf{Slide 31:}

WSD Using Random Walk Algorithm

\textbf{Slide 34:}

\textbf{Page rank Algorithm :} ->

\textbf{Slide 36:}

KB Approaches –Conclusions

36

CF

Drawbacks of WSD using Selectional Restrictions

Needs exhaustive Knowledge Base.

Drawbacks of Overlap based approaches

Dictionary definitions are generally very small.

Dictionary entries rarely take into account the distributional constraints  of different word senses (e.g. selectional preferences, kinds of  prepositions, etc.  cigarette and ash never co-occur in a dictionary).

Suffer from the problem of sparse match.

Proper nouns are not present in a MRD. Hence these approaches fail to  capture the strong clues provided by proper nouns.

E.g. “Sachin Tendulkar” will be a strong indicator of the category “sports”.

Sachin Tendulkar plays cricket.

\textbf{Slide 37:}

Supervised methods

Supervised methods assumes that the context can provide enough evidence on its own to disambiguate words

Probably every machine learning algorithm  has been applied to WSD, including associated techniques such as feature selection, parameter optimization, and ensemble learning.

Support vector machines and memory-based learning have been shown to be the most successful approaches, to date, probably because they can cope with the high dimensionality of the feature space.

However, these supervised methods are subject to a new knowledge acquisition bottleneck since they rely on substantial amounts of manually sense tagged corpora for training, which are laborious and expensive to create.

\textbf{Slide 39:}

Word Sense  Disambiguation

Classification

\textbf{Slide 40:}

Dan  Jurafsky

\textbf{Classification:}   definition

\textbf{Input:} 

a word w and some features f

a fixed set of classes C = {c1, c2,…, cJ}

\textbf{Output:}  a predicted class c∈C

\textbf{Slide 41:}

Dan  Jurafsky

\textbf{Classification Methods:}   Supervised Machine

Learning

\textbf{Input:} 

a word w in a text window d (which we’ll call a  “document”)

a fixed set of classes	C = {c1, c2,…, cJ}

A training set of m hand-‐labeled text windows again  called	“documents” (d1,c1),....,(dm,cm)

\textbf{Output:} 

\textbf{a learned classifier γ:} d  c

2

2

\textbf{Slide 42:}

Dan  Jurafsky

\textbf{Classification Methods:}   Supervised Machine Learning

Any kind of classifier

Naive Bayes

Logistic regression

Neural Networks

Support-‐vector  machines

k-‐Nearest Neighbors

…

\textbf{Slide 43:}

Applying Naive Bayes to WSD

P(c) is the prior probability of that sense

Counting in a labeled training set.

P(w|c)	conditional probability of a word given a  particular sense

P(w|c) = count(w,c)/count(c)

We get both of these from a tagged corpus like SemCor

Can also generalize to look at other features besides  words.

Then it would be P(f|c)

Conditional probability of a feature given a sense


\subsection{PMI}

\textbf{Page 1:}

CHAPTER
J Pointwise Mutual Information
(PMI)
Analternativeweightingfunctiontotf-idf,PPMI(positivepointwisemutualinfor-
mation),isusedforterm-term-matrices,whenthevectordimensionscorrespondto
wordsratherthandocuments.PPMIdrawsontheintuitionthatthebestwaytoweigh
theassociationbetweentwowordsistoaskhowmuchmorethetwowordsco-occur
inourcorpusthanwewouldhaveaprioriexpectedthemtoappearbychance.
pointwise
mutual
Pointwisemutualinformation(Fano,1961)1isoneofthemostimportantcon-
information
ceptsinNLP.Iti

\textbf{Page 2:}

2 APPENDIXJ • POINTWISEMUTUALINFORMATION(PMI)
Moreformally,let’sassumewehaveaco-occurrencematrixFwithWrows(words)
andCcolumns(contexts),where f givesthenumberoftimeswordw occurswith
ij i
context c . This can be turned into a PPMI matrix where PPMI gives the PPMI
j ij
value of word w with context c (which we can also express as PPMI(w,c ) or
i j i j
PPMI(w=i,c= j))asfollows:
p ij = (cid:80)W i(cid:48)=1 (cid:80) f i C j j(cid:48)=1 f i(cid:48)j(cid:48) , p i∗ = (cid:80)W i(cid:48)= (cid:80) 1 (ci

\textbf{Page 3:}

3
p(w,context) p(w)
computer data result pie sugar p(w)
cherry 0.0002 0.0007 0.0008 0.0377 0.0021 0.0415
strawberry 0.0000 0.0000 0.0001 0.0051 0.0016 0.0068
digital 0.1425 0.1436 0.0073 0.0004 0.0003 0.2942
information 0.2838 0.3399 0.0323 0.0004 0.0011 0.6575
p(context) 0.4265 0.4842 0.0404 0.0437 0.0052
FigureJ.3 ReplacingthecountsinFig.J.1withjointprobabilities,showingthemarginals
intherightcolumnandthebottomrow.
computer data result pie sugar
cherry 0 0 0 4.38 3.30
strawberry 0 0 0 4.10 5.5

\textbf{Page 4:}

4 AppendixJ • PointwiseMutualInformation(PMI)
Church,K.W.andP.Hanks.1989.Wordassociationnorms,
mutualinformation,andlexicography.ACL.
Church,K.W.andP.Hanks.1990.Wordassociationnorms,
mutual information, and lexicography. Computational
Linguistics,16(1):22–29.
Dagan,I.,S.Marcus,andS.Markovitch.1993. Contextual
wordsimilarityandestimationfromsparsedata.ACL.
Fano,R.M.1961.TransmissionofInformation:AStatistical
TheoryofCommunications.MITPress.
Levy,O.,Y.Goldberg,andI.Dagan.2015. Improvingdis-
tribut


\subsection{WordSenses}

\textbf{Page 1:}

Speech and Language Processing. Daniel Jurafsky & James H. Martin. Copyright © 2026. All
rights reserved. Draft of January 25, 2026.
CHAPTER
I Word Senses and WordNet
LadyBracknell. Areyourparentsliving?
Jack. Ihavelostbothmyparents.
Lady Bracknell. To lose one parent, Mr. Worthing, may be regarded as a
misfortune;tolosebothlookslikecarelessness.
OscarWilde,TheImportanceofBeingEarnest
ambiguous Wordsareambiguous: thesamewordcanbeusedtomeandifferentthings. In
Chapter 5 we saw that the word “mouse

\textbf{Page 2:}

2 APPENDIXI • WORDSENSESANDWORDNET
analyticdirection.
I.1 Word Senses
wordsense Asense(orwordsense)isadiscreterepresentationofoneaspectofthemeaningof
aword. Looselyfollowinglexicographictradition, werepresenteachsensewitha
superscript: bank1 andbank2,mouse1 andmouse2. Incontext,it’seasytoseethe
differentmeanings:
mouse1: .... amousecontrollingacomputersystemin1968.
mouse2: .... aquietanimallikeamouse
bank1: ...abankcanholdtheinvestmentsinacustodialaccount...
bank2: ...asagricultureburgeonsonthee

\textbf{Page 3:}

I.1 • WORDSENSES 3
Yet despite their circularity and lack of formal representation, glosses can still
beusefulforcomputationalmodelingofsenses.Thisisbecauseaglossisjustasen-
tence,andfromsentenceswecancomputesentenceembeddingsthattellussome-
thing about the meaning of the sense. Dictionaries often give example sentences
alongwithglosses,andthesecanagainbeusedtohelpbuildasenserepresentation.
Thesecondwaythatthesaurusesofferfordefiningasenseis—likethedictionary
definitions—definingasensethroughits

\textbf{Page 4:}

4 APPENDIXI • WORDSENSESANDWORDNET
I.2 Relations Between Senses
Thissectionexplorestherelationsbetweenwordsenses,especiallythosethathave
receivedsignificantcomputationalinvestigationlikesynonymy,antonymy,andhy-
pernymy.
Synonymy
We introduced in Chapter 5 the idea that when two senses of two different words
synonym (lemmas) are identical, or nearly identical, we say the two senses are synonyms.
Synonymsincludesuchpairsas
couch/sofa vomit/throwup filbert/hazelnut car/automobile
And we mentioned t

\textbf{Page 5:}

I.2 • RELATIONSBETWEENSENSES 5
andhyponym)areverysimilarandhenceeasilyconfused;forthisreason,theword
superordinate superordinateisoftenusedinsteadofhypernym.
Superordinate vehicle fruit furniture mammal
Subordinate car mango chair dog
We can define hypernymy more formally by saying that the class denoted by
the superordinate extensionally includes the class denoted by the hyponym. Thus,
theclassofanimalsincludesasmembersalldogs,andtheclassofmovingactions
includes all walking actions. Hypernymy c

\textbf{Page 6:}

6 APPENDIXI • WORDSENSESANDWORDNET
I.3 WordNet: A Database of Lexical Relations
The most commonly used resource for sense relations in English and many other
WordNet languages is the WordNet lexical database (Fellbaum, 1998). English WordNet
consists of three separate databases, one each for nouns and verbs and a third for
adjectivesandadverbs;closedclasswordsarenotincluded. Eachdatabasecontains
asetoflemmas,eachoneannotatedwithasetofsenses. TheWordNet3.0release
has117,798nouns, 11,529verbs, 22,

\textbf{Page 7:}

I.3 • WORDNET: ADATABASEOFLEXICALRELATIONS 7
supersense calledsupersenses,becausetheyactascoarsesemanticcategoriesorgroupingsof
senses which can be useful when word senses are too fine-grained (Ciaramita and
Johnson2003,CiaramitaandAltun2006). Supersenseshavealsobeendefinedfor
adjectives(Tsvetkovetal.,2014)andprepositions(Schneideretal.,2018).
Category Example Category Example Category Example
ACT service GROUP place PLANT tree
ANIMAL dog LOCATION area POSSESSION price
ARTIFACT car MOTIVE reason

\textbf{Page 8:}

8 APPENDIXI • WORDSENSESANDWORDNET
bass3, basso (an adult male singer with the lowest voice)
=> singer, vocalist, vocalizer, vocaliser
=> musician, instrumentalist, player
=> performer, performing artist
=> entertainer
=> person, individual, someone...
=> organism, being
=> living thing, animate thing,
=> whole, unit
=> object, physical object
=> physical entity
=> entity
bass7 (member with the lowest range of a family of instruments)
=> musical instrument, instrument
=> device
=> instrumentalit

\textbf{Page 9:}

I.4 • WORDSENSEDISAMBIGUATION 9
I.4.1 WSD:TheTaskandDatasets
In this section we introduce the task setup for WSD, and then turn to algorithms.
The inventory of sense tags depends on the task. For sense tagging in the context
oftranslationfromEnglishtoSpanish,thesensetaginventoryforanEnglishword
might be the set of different Spanish translations. For automatic indexing of med-
ical articles, the sense-tag inventory might be the set of MeSH (Medical Subject
Headings) thesaurus entries. Or we can u

\textbf{Page 10:}

10 APPENDIXI • WORDSENSESANDWORDNET
y y
5 6
y
3
stand1: side1:
y y
1 bass1: 4 upright relative
low range … region
electric1: … player1: stand5: …
using bass4: in game bear side3:
electricity sea fish player2: … of body
electric2: y … musician stand10: …
tense 2 bass7: player3: put side11:
electric3: instrument actor upright slope
thrilling guitar1 … … … …
x x x x x x
1 2 3 4 5 6
an electric guitar and bass player stand off to one side
FigureI.8 The all-words WSD task, mapping from input words (x

\textbf{Page 11:}

I.4 • WORDSENSEDISAMBIGUATION 11
sensewhosesenseembeddinghasthehighestcosinewitht:
sense(t)= argmax cosine(t,v s ) (I.14)
s∈senses(t)
Fig.I.9illustratesthemodel.
5
4 find v
find
v 1
find
v
9
find
v
c c c c c
I found the jar empty
ENCODER
I found the jar empty
FigureI.9 Thenearest-neighboralgorithmforWSD.Ingreenarethecontextualembed-
dingsprecomputedforeachsenseofeachword; herewejustshowafewofthesensesfor
find. A contextual embedding is computed for the target word found, and then the nearest
nei

\textbf{Page 12:}

12 APPENDIXI • WORDSENSESANDWORDNET
I.5 Alternate WSD algorithms and Tasks
I.5.1 Feature-BasedWSD
Feature-based algorithms for WSD are extremely simple and function almost as
wellascontextuallanguagemodelalgorithms. ThebestperformingIMSalgorithm
(ZhongandNg,2010),augmentedbyembeddings(Iacobaccietal.2016,Raganato
etal.2017b), usesanSVMclassifiertochoosethesenseforeachinputwordwith
thefollowingsimplefeaturesofthesurroundingwords:
• part-of-speech tags (for a window of 3 words on each side, stoppin

\textbf{Page 13:}

I.5 • ALTERNATEWSDALGORITHMSANDTASKS 13
functionSIMPLIFIEDLESK(word,sentence)returnsbestsenseofword
best-sense←mostfrequentsenseforword
max-overlap←0
context←setofwordsinsentence
foreachsenseinsensesofworddo
signature←setofwordsintheglossandexamplesofsense
overlap←COMPUTEOVERLAP(signature,context)
ifoverlap>max-overlapthen
max-overlap←overlap
best-sense←sense
end
return(best-sense)
FigureI.10 TheSimplifiedLeskalgorithm.TheCOMPUTEOVERLAPfunctionreturnsthe
numberofwordsincommonbetweentwosets,ignor

\textbf{Page 14:}

14 APPENDIXI • WORDSENSESANDWORDNET
F There’salotoftrashonthebedoftheriver—
IkeepaglassofwaternexttomybedwhenIsleep
F Justifythemargins—Theendjustifiesthemeans
T Airpollution—Openawindowandletinsomeair
T Theexpandedwindowwillgiveustimetocatchthethieves—
Youhaveatwo-hourwindowofclearweathertofinishworkingonthelawn
FigureI.11 Positive (T) and negative (F) pairs from the WiC dataset (Pilehvar and
Camacho-Collados,2019).
word-in-context first cluster the word senses into coarser clusters, so that th

\textbf{Page 15:}

I.6 • USINGTHESAURUSESTOIMPROVEEMBEDDINGS 15
I.6 Using Thesauruses to Improve Embeddings
Thesauruses have also been used to improve both static and contextual word em-
beddings. For example, static word embeddings have a problem with antonyms.
A word like expensive is often very similar in embedding cosine to its antonym
like cheap. Antonymy information from thesauruses can help solve this problem;
Fig.I.12showsnearestneighborstosometargetwordsinGloVe,andtheimprove-
mentafteronesuchmethod.
Befor

\textbf{Page 16:}

16 APPENDIXI • WORDSENSESANDWORDNET
1. Computeacontextvectorcfort.
2. Retrieveallsensevectorss forw.
j
3. Assignt tothesenserepresentedbythesensevectors thatisclosesttot.
j
All we need is a clustering algorithm and a distance metric between vectors.
Clusteringisawell-studiedproblemwithawidenumberofstandardalgorithmsthat
canbeappliedtoinputsstructuredasvectorsofnumericalvalues(DudaandHart,
1973). Afrequentlyusedtechniqueinlanguageapplicationsisknownasagglom-
agglomerative erative clustering. In t

\textbf{Page 17:}

HISTORICALNOTES 17
• An important baseline for WSD is the most frequent sense, equivalent, in
WordNet,totakethefirstsense.
• Another baseline is a knowledge-based WSD algorithm called the Lesk al-
gorithmwhichchoosesthesensewhosedictionarydefinitionsharesthemost
wordswiththetargetword’sneighborhood.
• Wordsenseinductionisthetaskoflearningwordsensesunsupervised.
Historical Notes
Word sense disambiguation traces its roots to some of the earliest applications of
digitalcomputers. Theinsightthatunde

\textbf{Page 18:}

18 APPENDIXI • WORDSENSESANDWORDNET
resourcesincludeAmsler’s1981useoftheMerriamWebsterdictionaryandLong-
man’sDictionaryofContemporaryEnglish(BoguraevandBriscoe,1989).
Supervised approaches to disambiguation began with the use of decision trees
byBlack(1988). InadditiontotheIMSandcontextual-embeddingbasedmethods
forsupervisedWSD,recentsupervisedalgorithmsincludesencoder-decodermodels
(Raganatoetal.,2017a).
The need for large amounts of annotated text in supervised methods led early
ontoinvestiga

\textbf{Page 19:}

Exercises 19
Agirre,E.andO.L.deLacalle.2003. ClusteringWordNet Hirst,G.1987. SemanticInterpretationandtheResolution
wordsenses.RANLP2003. ofAmbiguity.CambridgeUniversityPress.
Agirre,E.andP.Edmonds,eds.2006. WordSenseDisam- Hirst,G.1988.Resolvinglexicalambiguitycomputationally
biguation:AlgorithmsandApplications.Kluwer. withspreadingactivationandpolaroidwords. InS.L.
Small,G.W.Cottrell,andM.K.Tanenhaus,eds,Lexical
Amsler, R. A. 1981. A taxonomy for English nouns and
AmbiguityResolution,73–108.Mo

\textbf{Page 20:}

20 AppendixI • WordSensesandWordNet
Manandhar,S.,I.P.Klapaftis,D.Dligach,andS.Pradhan. Ponzetto,S.P.andR.Navigli.2010. Knowledge-richword
2010. SemEval-2010task14: Wordsenseinduction& sensedisambiguationrivalingsupervisedsystems.ACL.
disambiguation.SemEval.
Pu, X., N. Pappas, J. Henderson, and A. Popescu-Belis.
Martin,J.H.1986.Theacquisitionofpolysemy.ICML. 2018. Integratingweaklysupervisedwordsensedisam-
Masterman,M.1957.Thethesaurusinsyntaxandsemantics. biguationintoneuralmachinetranslation. T

\textbf{Page 21:}

Exercises 21
Weaver, W.1949/1955. Translation. InW.N.Lockeand
A.D.Boothe, eds, MachineTranslationofLanguages,
15–23.MITPress. Reprintedfromamemorandumwrit-
tenbyWeaverin1949.
Wilks,Y.1975a.Anintelligentanalyzerandunderstanderof
English.CACM,18(5):264–274.
Wilks, Y.1975b. Preferencesemantics. InE.L.Keenan,
ed., The Formal Semantics of Natural Language, 329–
350.CambridgeUniv.Press.
Wilks, Y. 1975c. A preferential, pattern-seeking, seman-
ticsfornaturallanguageinference.ArtificialIntelligence,
6(1


\section{Module 3: Text Classification with Naive Bayes}


\subsection{Text Classification using Naive Bayes  (3)}

\textbf{Slide 1:}

Chapter 2Text Classification

\textbf{Slide 2:}

What is Text Classification?

Have you performed The Task of Text Classification ?

\textbf{Slide 3:}

Some Examples of Text Classification?

\textbf{Slide 4:}

Is this spam?  How will we decide ?

\textbf{Slide 5:}

What is the subject of this medical article?

Antogonists and Inhibitors

Blood Supply

Chemistry

Drug Therapy

Embryology

Epidemiology

…

5

MeSH Subject Category Hierarchy

?

MEDLINE Article

\textbf{Slide 6:}

Positive or negative movie review?

...zany characters and richly applied satire, and some great plot twists

It was pathetic. The worst part about it was the boxing scenes...

...awesome caramel sauce and sweet toasty almonds. I love this place!

...awful pizza and ridiculously overpriced...

6

+

+

−

−

\textbf{Slide 7:}

Positive or negative movie review?

...zany characters and richly applied satire, and some great plot twists

It was pathetic. The worst part about it was the boxing scenes...

...awesome caramel sauce and sweet toasty almonds. I love this place!

...awful pizza and ridiculously overpriced...

7

+

+

−

−

\textbf{Slide 8:}

Why sentiment analysis?

\textbf{Movie:}   is this review positive or negative?

\textbf{Products:}  what do people think about the new iPhone?

\textbf{Public sentiment:}  how is consumer confidence?

\textbf{Politics:}  what do people think about this candidate or issue?

\textbf{Prediction:}  predict election outcomes or market trends from sentiment

8

\textbf{Slide 9:}

Scherer Typology of Affective States

\textbf{Emotion:}  brief organically synchronized … evaluation of a major event

angry, sad, joyful, fearful, ashamed, proud, elated

Mood: diffuse non-caused low-intensity long-duration change in subjective feeling

cheerful, gloomy, irritable, listless, depressed, buoyant

Interpersonal stances: affective stance toward another person in a specific interaction

friendly, flirtatious, distant, cold, warm, supportive, contemptuous

Attitudes: enduring, affectively colored beliefs, dispositions towards objects or persons

liking, loving, hating, valuing, desiring

Personality traits: stable personality dispositions and typical behavior tendencies

nervous, anxious, reckless, morose, hostile, jealous

\textbf{Slide 10:}

Basic Sentiment Classification

Sentiment analysis is the detection of attitudes

Simple task we focus on in this chapter

Is the attitude of this text positive or negative?

We return to affect classification in later chapters

\textbf{Slide 11:}

Text Classification

Sentiment analysis

Spam detection

Authorship identification

Language Identification

Assigning subject categories, topics, or genres

…

\textbf{Slide 12:}

\textbf{Text Classification:}  definition

\textbf{Input:} 

a document d

a fixed set of classes  C = {c1, c2,…, cJ}

\textbf{Output:}  a predicted class c  C

\textbf{Slide 13:}

\textbf{Classification Methods:}   Hand-coded rules

Rules based on combinations of words or other features

\textbf{spam:}  black-list-address OR (“dollars” AND “you have been selected”)

Accuracy can be high

If rules carefully refined by expert

But building and maintaining these rules is expensive

\textbf{Slide 14:}

\textbf{Classification Methods:}  Supervised Machine Learning

\textbf{Input:} 

a document d

a fixed set of classes  C = {c1, c2,…, cJ}

A training set of m hand-labeled documents (d1,c1),....,(dm,cm)

\textbf{Output:} 

\textbf{a learned classifier γ:} d  c

14

\textbf{Slide 15:}

\textbf{Classification Methods:} Supervised Machine Learning Types

Any kind of classifier

Naïve Bayes

Logistic regression

Neural networks

k-Nearest Neighbors

SVM…

\textbf{Slide 16:}

How to use naïve bayes classifier for text classification?

\textbf{Slide 17:}

Generative and Discriminative Classifiers

Naive Bayes is a generative classifier

\textbf{by contrast:} 

Logistic regression/SVM is a discriminative classifier

\textbf{Slide 19:}

Generative and Discriminative Classifiers

Suppose we're distinguishing cat from dog images

imagenet

imagenet

\textbf{Slide 20:}

\textbf{Generative Classifier:} 

Build a model of what's in a cat image

Knows about whiskers, ears, eyes

\textbf{Assigns a probability to any image:} 

how cat-y is this image?

Also build a model for dog images

\textbf{Now given a new image:} 

Run both models and see which one fits better

\textbf{Slide 21:}

Discriminative Classifier

Just try to distinguish dogs from cats

Oh look, dogs have collars!

Let's ignore everything else

\textbf{Slide 22:}

Finding the correct class c from a document d inGenerative vs Discriminative Classifiers

Naive Bayes

22

\textbf{Slide 23:}

How SVM Makes Decision

\textbf{y^=sign(w⋅x+b)Where:} 

w = weight vector

b = bias

sign() → decides class (+1 or −1)

\textbf{Slide 24:}

Why is Naïve Bayes generative?

NB assumes that features are conditionally independent given the class.

It first estimates P(X∣Y)(the likelihood of the features given the class) and P(Y) (the prior probability of each class).

Then, it uses Bayes' theorem to compute

Since it explicitly models how the data is generated, it is a generative model.

\textbf{Slide 25:}

Why is Logistic Regression discriminative?​

It does not model how the data is generated; instead, it directly finds the best decision boundary.​

\textbf{It learns a function of the form:} 

It maximizes the likelihood of the observed data, adjusting weights to separate classes optimally.​

Since it directly maps input X to output Y without modeling P(X) or P(X∣Y), it is a discriminative model.​

\textbf{Slide 26:}

Naive Bayes Intuition

Simple ("naive") classification method based on Bayes rule

Relies on very simple representation of document

Bag of words

\textbf{Slide 27:}

The Bag of Words Representation

27

\textbf{Slide 28:}

The bag of words representation

γ(

)=c

\textbf{Slide 29:}

Supervised Learning

\textbf{Slide 30:}

Bayes’ Rule Applied to Documents and Classes

For a document d and a class c

\textbf{Slide 31:}

Naive Bayes Classifier (I)

MAP is “maximum a posteriori”  = most likely class

Bayes Rule

Dropping the denominator

\textbf{Slide 32:}

Naive Bayes Classifier (II)

Document d represented as features x1..xn

"Likelihood"

"Prior"

\textbf{Slide 33:}

Naïve Bayes Classifier (IV)

How often does this class occur?

O(|X|n•|C|) parameters

We can just count the relative frequencies in a corpus

Could only be estimated if a very, very large number of training examples was available.

\textbf{Slide 34:}

Multinomial Naive Bayes Independence Assumptions

\textbf{Bag of Words assumption:}  Assume position doesn’t matter

Conditional Independence: Assume the feature probabilities P(xi|cj) are independent of others, given the class.

\textbf{Slide 35:}

Multinomial Naive Bayes Classifier

\textbf{Slide 36:}

Applying Multinomial Naive Bayes Classifiers to Text Classification

positions  all word positions in test document

\textbf{Slide 37:}

Problems with multiplying lots of probs

\textbf{There's a problem with this:} 

Multiplying lots of probabilities can result in floating-point underflow!

.0006 * .0007 * .0009 * .01 * .5 * .000008….

\textbf{Idea:}    Use logs, because  log(ab) = log(a) + log(b)

We'll sum logs of probabilities instead of multiplying probabilities!

\textbf{Slide 38:}

We actually do everything in log space

\textbf{Instead of this:} 

\textbf{This:} 

\textbf{Notes:} 

1) Taking log doesn't change the ranking of classes!

The class with highest probability also has highest log probability!

\textbf{2) It's a linear model:} 

\textbf{Just a max of a sum of weights:}  a linear function of the inputs

So naive bayes is a linear classifier

\textbf{Slide 39:}

Parameter estimation

Create mega-document for topic j by concatenating all docs in this topic

Use frequency of w in mega-document

fraction of times word wi appears

among all words in documents of topic cj

\textbf{Slide 40:}

Learning the Multinomial Naive Bayes Model

\textbf{First attempt:}  maximum likelihood estimates

simply use the frequencies in the data

Sec.13.3

\textbf{Slide 41:}

Problem with Maximum Likelihood

What if we have seen no training documents with the word fantastic  and classified in the topic positive (thumbs-up)?

Zero probabilities cannot be conditioned away, no matter the other evidence!

Sec.13.3

\textbf{Slide 42:}

Laplace (add-1) smoothing for Naïve Bayes

\textbf{Slide 43:}

\textbf{Multinomial Naïve Bayes:}  Learning

Calculate P(cj) terms

For each cj in C do

docsj  all docs with  class =cj

Calculate P(wk | cj) terms

Textj  single doc containing all docsj

For each word wk in Vocabulary

nk  \# of occurrences of wk in Textj

From training corpus, extract Vocabulary

\textbf{Slide 44:}

Unknown words

What about unknown words

that appear in our test data

but not in our training data or vocabulary?

We ignore them

Remove them from the test document!

Pretend they weren't there!

Don't include any probability for them at all!

Why don't we build an unknown word model?

It doesn't help: knowing which class has more unknown words is not generally helpful!

\textbf{Slide 45:}

Stop words

Some systems ignore stop words

\textbf{Stop words:}  very frequent words like the and a.

Sort the vocabulary by word frequency in training set

Call the top 10 or 50 words the stopword list.

Remove all stop words from both training and test sets

As if they were never there!

But removing stop words doesn't usually help

So in practice most NB algorithms use all words and don't use stopword lists

\textbf{Slide 46:}

Text Classification and Naive Bayes

Sentiment and Binary Naive Bayes

\textbf{Slide 47:}

Let's do a worked sentiment example!

\textbf{Slide 48:}

\textbf{Step1:} 

\textbf{1. Prior from training:} 

P(-) = 3/5

P(+) = 2/5

\textbf{Slide 49:}

Step 2 and 3

2. Drop "with"

\textbf{3. Likelihoods from training:} 

\textbf{Slide 50:}

\textbf{4. Scoring the test set:} 

\textbf{Slide 51:}

Solve the following

\textbf{Slide 52:}

Training Data

\textbf{Slide 55:}

\textbf{Step 3:}  Word Frequency Table

Spam Class

D1 + D2 → buy cheap meds cheap meds online

Total words in Spam = 6

Ham Class

D3 + D4 → meeting schedule today project meeting today

Total words in Ham = 6

\textbf{Slide 56:}

Optimizing for sentiment analysis

For tasks like sentiment, word occurrence seems to be more important than word frequency.

The occurrence of the word fantastic tells us a lot

The fact that it occurs 5 times may not tell us much more.

Binary multinominal naive bayes, or binary NB

Clip our word counts at 1

\textbf{Slide 57:}

\textbf{Binary Multinomial Naïve Bayes:}  Learning

Calculate P(cj) terms

For each cj in C do

docsj  all docs with  class =cj

From training corpus, extract Vocabulary

Calculate P(wk | cj) terms

\textbf{Remove duplicates in each doc:} 

For each word type w in docj

Retain only a single instance of w

\textbf{Slide 58:}

Binary Multinomial Naive Bayes on a test document d

58

First remove all duplicate words from d

\textbf{Then compute NB using the same equation:} 

\textbf{Slide 59:}

Binary multinominal naive Bayes

\textbf{Slide 60:}

\textbf{Sentiment Classification:}  Dealing with Negation

I really like this movie

I really don't like this movie

Negation changes the meaning of "like" to negative.

Negation can also change negative to positive-ish

Don't dismiss this film

Doesn't let us get bored

\textbf{Slide 64:}

\textbf{Sentiment Classification:}  Dealing with Negation

\textbf{Simple baseline method:} 

\textbf{Add NOT\_ to every word between negation and following punctuation:} 

didn’t like this movie , but I

didn’t NOT\_like NOT\_this NOT\_movie but I

Das, Sanjiv and Mike Chen. 2001. Yahoo! for Amazon: Extracting market sentiment from stock message boards. In Proceedings of the Asia Pacific Finance Association Annual Conference (APFA).

Bo Pang, Lillian Lee, and Shivakumar Vaithyanathan.  2002.  Thumbs up? Sentiment Classification using Machine Learning Techniques. EMNLP-2002, 79—86.


\subsection{Naive Bayes Classifiers}

\textbf{Page 1:}

Speech and Language Processing. Daniel Jurafsky & James H. Martin. Copyright © 2026. All
rights reserved. Draft of January 6, 2026.
CHAPTER
B Naive Bayes, Text Classification,
and Sentiment
Classification lies at the heart of both human and machine intelligence. Deciding
what letter, word, or image has been presented to our senses, recognizing faces
or voices, sorting mail, assigning grades to homeworks; these are all examples of
assigningacategorytoaninput.Thepotentialchallengesofthistaskarehig

\textbf{Page 2:}

2 APPENDIXB • NAIVEBAYES,TEXTCLASSIFICATION,ANDSENTIMENT
Finally, one of the oldest tasks in text classification is assigning a library sub-
ject category or topic label to a text. Deciding whether a research paper concerns
epidemiologyorinstead,perhaps,embryology,isanimportantcomponentofinfor-
mationretrieval.Varioussetsofsubjectcategoriesexist,suchastheMeSH(Medical
SubjectHeadings) thesaurus. Infact,aswewillsee,subjectcategoryclassification
isthetaskforwhichthenaiveBayesalgorithmwasinventedin1

\textbf{Page 3:}

B.1 • NAIVEBAYESCLASSIFIERS 3
howthefeaturesinteract.
TheintuitionoftheclassifierisshowninFig.B.1. Werepresentatextdocument
bagofwords as if it were a bag of words, that is, an unordered set of words with their position
ignored,keepingonlytheirfrequencyinthedocument. Intheexampleinthefigure,
insteadofrepresentingthewordorderinallthephraseslike“Ilovethismovie”and
“I would recommend it”, we simply note that the word I occurred 5 times in the
entireexcerpt,thewordit6times,thewordslove,recommend,and

\textbf{Page 4:}

4 APPENDIXB • NAIVEBAYES,TEXTCLASSIFICATION,ANDSENTIMENT
WecanconvenientlysimplifyEq.B.3bydroppingthedenominatorP(d). This
ispossiblebecausewewillbecomputing P(d|c)P(c) foreachpossibleclass.ButP(d)
P(d)
doesn’tchangeforeachclass;wearealwaysaskingaboutthemostlikelyclassfor
the same document d, which must have the same probability P(d). Thus, we can
choosetheclassthatmaximizesthissimplerformula:
cˆ=argmaxP(c|d)=argmaxP(d|c)P(c) (B.4)
c∈C c∈C
WecallNaiveBayesagenerativemodelbecausewecanreadEq.B.4as

\textbf{Page 5:}

B.2 • TRAININGTHENAIVEBAYESCLASSIFIER 5
bywalkinganindexthrougheverywordpositioninthedocument:
positions ← allwordpositionsintestdocument
(cid:89)
c NB = argmaxP(c) P(w i |c) (B.9)
c∈C
i∈positions
NaiveBayescalculations, likecalculationsforlanguagemodeling, aredoneinlog
space, to avoid underflow and increase speed. Thus Eq. B.9 is generally instead
expressed1as
(cid:88)
c NB = argmaxlogP(c)+ logP(w i |c) (B.10)
c∈C
i∈positions
Byconsideringfeaturesinlogspace,Eq.B.10computesthepredictedclassasali

\textbf{Page 6:}

6 APPENDIXB • NAIVEBAYES,TEXTCLASSIFICATION,ANDSENTIMENT
But since naive Bayes naively multiplies all the feature likelihoods together, zero
probabilities in the likelihood term for any class will cause the probability of the
classtobezero,nomattertheotherevidence!
The simplest solution is the add-one (Laplace) smoothing introduced in Chap-
ter3.WhileLaplacesmoothingisusuallyreplacedbymoresophisticatedsmoothing
algorithms in language modeling, it is commonly used in naive Bayes text catego-
riza

\textbf{Page 7:}

B.4 • OPTIMIZINGFORSENTIMENTANALYSIS 7
functionTRAINNAIVEBAYES(D,C)returnsV, log P(c),log P(w|c)
foreachclassc∈C \#CalculateP(c)terms
N =numberofdocumentsinD
doc
Nc=numberofdocumentsfromDinclassc
Nc
logprior[c]← log
N
doc
V←vocabularyofD
bigdoc[c]←append(d)ford∈Dwithclassc
foreachwordwinV \# CalculateP(w|c)terms
count(w,c)←\#ofoccurrencesofwinbigdoc[c]
count(w,c) + 1
loglikelihood[w,c]← log (cid:80) (count(w(cid:48),c) + 1)
w(cid:48)inV
returnlogprior,loglikelihood,V
functionTESTNAIVEBAYES(testdoc,

\textbf{Page 8:}

8 APPENDIXB • NAIVEBAYES,TEXTCLASSIFICATION,ANDSENTIMENT
First,forsentimentclassificationandanumberofothertextclassificationtasks,
whetherawordoccursornotseemstomattermorethanitsfrequency. Thusitoften
improves performance to clip the word counts in each document at 1 (see the end
ofthechapterforpointerstotheseresults). Thisvariantiscalledbinarymultino-
binarynaive mialnaiveBayesorbinarynaiveBayes. Thevariantusesthesamealgorithmasin
Bayes
Fig.B.2exceptthatforeachdocumentweremoveallduplicatewordsb

\textbf{Page 9:}

B.5 • NAIVEBAYESFOROTHERTEXTCLASSIFICATIONTASKS 9
thesenegationwordsandthepredicatestheymodify,butthissimplebaselineworks
quitewellinpractice.
Finally, in some situations we might have insufficient labeled training data to
trainaccuratenaiveBayesclassifiersusingallwordsinthetrainingsettoestimate
positive and negative sentiment. In such cases we can instead derive the positive
sentiment and negative word features from sentiment lexicons, lists of words that are pre-
lexicons
annotatedwithpositive

\textbf{Page 10:}

10 APPENDIXB • NAIVEBAYES,TEXTCLASSIFICATION,ANDSENTIMENT
• Emailsubjectlinecontains“onlinepharmaceutical”
• HTMLhasunbalanced“head”tags
• Claimsyoucanberemovedfromthelist
languageid For other tasks, like language id—determining what language a given piece
of text is written in—the most effective naive Bayes features are not words at all,
but charactern-grams, 2-grams (‘zw’)3-grams (‘nya’, ‘Vo’), or4-grams (‘iez’,
‘thei’),or,evensimplerbyten-grams,whereinsteadofusingthemultibyteUnicode
character

\textbf{Page 11:}

B.7 • EVALUATION: PRECISION,RECALL,F-MEASURE 11
P(“Ilovethisfunfilm”|+) = 0.1×0.1×0.01×0.05×0.1=5×10−7
P(“Ilovethisfunfilm”|−) = 0.2×0.001×0.01×0.005×0.1=1.0×10−9
As it happens, the positive model assigns a higher probability to the sentence:
P(s|pos)>P(s|neg). Note that this is just the likelihood part of the naive Bayes
model; once we multiply in the prior a full naive Bayes model might well make a
differentclassificationdecision.
B.7 Evaluation: Precision, Recall, F-measure
Tointroducethemeth

\textbf{Page 12:}

12 APPENDIXB • NAIVEBAYES,TEXTCLASSIFICATION,ANDSENTIMENT
gold standard labels
gold positive gold negative
system p s o y s s i t t e i m ve true positive false positive precision = tp t + p fp
output
system
labels negative false negative true negative
tp tp+tn
recall = accuracy =
tp+fn tp+fp+tn+fn
FigureB.4 Aconfusionmatrixforvisualizinghowwellabinaryclassificationsystemper-
formsagainstgoldstandardlabels.
That’swhyinsteadofaccuracywegenerallyturntotwoothermetricsshownin
precision Fig.B.4: prec

\textbf{Page 13:}

B.7 • EVALUATION: PRECISION,RECALL,F-MEASURE 13
andhenceF-measureis
1 (cid:18) 1−α (cid:19) (β2+1)PR
F = or withβ2= F = (B.18)
α1+(1−α)1 α β2P+R
P R
Harmonic mean is used because the harmonic mean of two values is closer to the
minimumofthetwovaluesthanthearithmeticmeanis. Thusitweighsthelowerof
thetwonumbersmoreheavily,whichismoreconservativeinthissituation.
B.7.1 Evaluatingwithmorethantwoclasses
Uptonowwehavebeendescribingtextclassificationtaskswithonlytwoclasses.
But lots of classification ta

\textbf{Page 14:}

14 APPENDIXB • NAIVEBAYES,TEXTCLASSIFICATION,ANDSENTIMENT
Class 1: Urgent Class 2: Normal Class 3: Spam Pooled
true true true true true true true true
urgent not normal not spam not yes no
system system system system
8 11 60 55 200 33 268 99
urgent normal spam yes
system 8 340 system 40 212 system 51 83 system 99 635
not not not no
8 60 200 microaverage 268
precision = = .42 precision = = .52 precision = = .86 = = .73
8+11 60+55 200+33 precision 268+99
macroaverage .42+.52+.86
= = .60
precision 

\textbf{Page 15:}

B.9 • STATISTICALSIGNIFICANCETESTING 15
Training Iterations Testing
1 Dev Training
2 Dev Training
3 Dev Training
4 Dev Training
Test
5 Training Dev Training
Set
6 Training Dev
7 Training Dev
8 Training Dev
9 Training Dev
10 Training Dev
FigureB.7 10-foldcross-validation
B.9 Statistical Significance Testing
Inbuildingsystemsweoftenneedtocomparetheperformanceoftwosystems.How
can we know if the new system we just built is better than our old one? Or better
thansomeothersystemdescribedintheliteratur

\textbf{Page 16:}

16 APPENDIXB • NAIVEBAYES,TEXTCLASSIFICATION,ANDSENTIMENT
wewouldencounterthevalueofδ(x)thatwefound,ifwerepeatedtheexperiment
p-value a great many times. We formalize this likelihood as the p-value: the probability,
assumingthenullhypothesisH istrue,ofseeingtheδ(x)thatwesaworoneeven
0
greater
P(δ(X)≥δ(x)|H 0 istrue) (B.21)
Soinourexample,thisp-valueistheprobabilitythatwewouldseeδ(x)assuming
A is not better than B. If δ(x) is huge (let’s say A has a very respectable F of .9
1
andBhasaterribleF of

\textbf{Page 17:}

B.9 • STATISTICALSIGNIFICANCETESTING 17
Consideratinytextclassificationexamplewithatestsetxof10documents.The
first row of Fig. B.8 shows the results of two classifiers (A and B) on this test set.
Eachdocumentislabeledbyoneofthefourpossibilities(AandBbothright,both
wrong, A right and B wrong, A wrong and B right). A slash through a letter ((cid:19)B)
meansthatthatclassifiergottheanswerwrong. OnthefirstdocumentbothAand
Bgetthecorrectclass(AB),whileontheseconddocumentAgotitrightbutBgot
it wrong (A(

\textbf{Page 18:}

18 APPENDIXB • NAIVEBAYES,TEXTCLASSIFICATION,ANDSENTIMENT
δ(x(i))exceedstheexpectedvalueofδ(x)byδ(x)ormore:
b
1(cid:88) (cid:16) (cid:17)
p-value(x) = 1 δ(x(i))−δ(x)≥δ(x)
b
i=1
b
1(cid:88) (cid:16) (cid:17)
= 1 δ(x(i))≥2δ(x) (B.22)
b
i=1
Soifforexamplewehave10,000testsetsx(i)andathresholdof.01,andinonly47
ofthetestsetsdowefindthatAisaccidentallybetterδ(x(i))≥2δ(x),theresulting
p-valueof.0047issmallerthan.01,indicatingthatthedeltawefound,δ(x)isindeed
sufficientlysurprisingandunlikelytohavehappene

\textbf{Page 19:}

B.11 • SUMMARY 19
toxicity icity detection is the task of detecting hate speech, abuse, harassment, or other
detection
kinds of toxic language. While the goal of such classifiers is to help reduce soci-
etalharm,toxicityclassifierscanthemselvescauseharms. Forexample,researchers
haveshownthatsomewidelyusedtoxicityclassifiersincorrectlyflagasbeingtoxic
sentencesthatarenon-toxicbutsimplymentionidentitieslikewomen(Parketal.,
2018), blind people (Hutchinson et al., 2020) or gay people (Dixon et al., 

\textbf{Page 20:}

20 APPENDIXB • NAIVEBAYES,TEXTCLASSIFICATION,ANDSENTIMENT
• Statistical significance tests should be used to determine whether we can be
confidentthatoneversionofaclassifierisbetterthananother.
• Designers of classifiers should carefully consider harms that may be caused
by the model, including its training data and other components, and report
modelcharacteristicsinamodelcard.
Historical Notes
Multinomial naive Bayes text classification was proposed by Maron (1961) at the
RANDCorporationforthet

\textbf{Page 21:}

EXERCISES 21
2014,Droretal.2017).
Featureselectionisamethodofremovingfeaturesthatareunlikelytogeneralize
well.Featuresaregenerallyrankedbyhowinformativetheyareabouttheclassifica-
information tiondecision. Averycommonmetric,informationgain,tellsushowmanybitsof
gain
informationthepresenceofthewordgivesusforguessingtheclass. Otherfeature
selection metrics include χ2, pointwise mutual information, and GINI index; see
YangandPedersen(1997)foracomparisonandGuyonandElisseeff(2003)foran
introductiontofe

\textbf{Page 22:}

22 AppendixB • NaiveBayes,TextClassification,andSentiment
Aggarwal,C.C.andC.Zhai.2012. Asurveyoftextclassi- Heckerman,D.,E.Horvitz,M.Sahami,andS.T.Dumais.
ficationalgorithms. InC.C.AggarwalandC.Zhai,eds, 1998.Abayesianapproachtofilteringjunke-mail.AAAI-
Miningtextdata,163–222.Springer. 98WorkshoponLearningforTextCategorization.
Bayes,T.1763. AnEssayTowardSolvingaProbleminthe Hu,M.andB.Liu.2004.Miningandsummarizingcustomer
DoctrineofChances,volume53.ReprintedinFacsimiles reviews.KDD.
ofTwoPapersb

\textbf{Page 23:}

Exercises 23
Pang, B., L. Lee, and S. Vaithyanathan. 2002. Thumbs
up?Sentimentclassificationusingmachinelearningtech-
niques.EMNLP.
Park,J.H.,J.Shin,andP.Fung.2018.Reducinggenderbias
inabusivelanguagedetection.EMNLP.
Pennebaker, J. W., R. J. Booth, and M. E. Francis. 2007.
LinguisticInquiryandWordCount:LIWC2007.Austin,
TX.
Popp,D.,R.A.Donovan,M.Crawford,K.L.Marsh,and
M.Peele.2003. Gender,race,andspeechstylestereo-
types.SexRoles,48(7-8):317–325.
Sahami,M.,S.T.Dumais,D.Heckerman,andE.Horvitz.
199


\section{Module 4: Recurrent Neural Networks}


\subsection{1. Chapter 4}

\textbf{Slide 1:}

Deep Learning Architectures

for Sequence Processing

\textbf{Slide 2:}

Recap of Feed Forward Neural Networks

Inputs → pass through hidden layers → output

Each layer has weights + activation functions

\textbf{Limitation:}  FFNNs assume inputs are independent (no memory of past  inputs).

Works well for images, tabular data, but not for sequences.

\textbf{Example to give students:} 

Predicting tomorrow’s weather using only today’s temperature (FFNN).

But in reality, weather depends on past few days → we need memory.

\textbf{Slide 3:}

Introduce the Idea of Sequences

\textbf{Some problems need context:} 

Sentences in language (words depend on previous words).

Stock price prediction (depends on past days).

Speech recognition (next sound depends on earlier sounds).

“If I say ‘The cat sat on the \_\_\_’, what word do you expect?” They’ll naturally answer ‘mat’. That shows context matters.

\textbf{Slide 4:}

Recurrent Neural Networks

Works well with sequence of data for example Natural language data, Time series data, sales forecasting, stock forecasting .

Problems in NLP using Machine Learning

For example, SPAM Classifier

NLP Text Preprocessing  technique used to convert text data to vector are Bag of Words ,TFIDF,wprd2Vec

Steps followed for Classification

Step 1   Covert the text into vector

Step 2   Apply machine learning algorithm(naïve Bayes)

\textbf{Problem :} -While converting text to vector sequence information is discarded

If the sequence information is decarded Accuracy of the result may get affected.

In many NLP applications sequence information is very important..

And for such applications RNN is important concept

\textbf{Slide 5:}

RNN

where the output from the previous step is fed as input to the current step

when in language models, it is required to predict the next word of a sentence, the previous words are required and hence there is a need to remember the previous words.

main and most important feature of RNN is its Hidden state which remembers some information about a sequence.

The state is also referred to as Memory State since it remembers the previous input to the network.

It uses the same parameters for each input as it performs the same task on all the inputs or hidden layers to produce the output.

\textbf{Slide 6:}

Architecture of RNN

\textbf{where:} 

ht -> current state

ht-1 -> previous state

xt -> input state

\textbf{Slide 7:}

Simple recurrent neural network illustrated as a feed-forward network

\textbf{Most significant change:}  new set of weights, U

connect the hidden layer from the previous time step to the current hidden layer.

determine how the network should make use of past context in calculating the output for the current input.

INPUT

OUTPUT

WEIGHT MATRIX

WEIGHT MATRIX

HIDDEN LAYER

PREVIOUS TIME STEP

weights from the hiddenlayer to the output layer

\textbf{Slide 8:}

\textbf{RNNs:}  Passing hidden state to next time step

While processing, it passes the previous hidden state to the next step of the sequence.

The hidden state acts as the memory.

It holds information on previous data the network has seen before.

\textbf{Slide 9:}

Simple-RNN abstraction

y2

y1

y3

\textbf{Slide 10:}

\textbf{RNNs:}  Understanding a single cell

\textbf{Slide 11:}

Unrolled RNN

\textbf{Slide 12:}

Training

we’ll use a training set, a loss function, and back-propagation

to obtain the gradients needed to adjust the weights in these recurrent networks.

\textbf{3 sets of weights to update while training :} 

W , the weights from the input layer to the hidden layer,

U, the weights from the previous hidden layer to the current hidden layer, and finally

V , the weights from the hidden layer to the output layer.

\textbf{Slide 13:}

Steps in Training

to compute the loss function for the output at time t we need the hidden layer from time t − 1.

the hidden layer at time t influences both the output at time t and the hidden layer at time t + 1 (and hence the output and loss at t + 1).

to assess the error accruing to ht , we’ll need to know its influence on both the current output as well as the ones that follow.

\textbf{Slide 14:}

Backpropagation algorithm – two pass algorithm for training the weights in RNNs.

Step 1:- we perform forward inference, computing ht , yt , accumulating the loss at each step-in time, saving thevalue of the hidden layer at each step for use at the next time step.

Step 2:- we process the sequence in reverse, computing the required gradients as we go, computing and saving the error term for use in the hidden layer for each step backward in time.

This general approach is commonly referred to as Backpropagation Through Time

\textbf{Slide 15:}

Language Modeling

Language Modelling (LM) is the task of predicting what word comes next or more generally, a system that assigns probability to a piece of a text sequence.

N-gram is the simplest language model, and its performance is limited by its lack of complexity.

Simplistic models like this one cannot achieve fluency, enough language variation and correct writing style for long texts.

For these reasons, neural networks (NN) are explored as the new gold standard despite their complexity and Recurrent Neural Networks (RNN) became a fundamental architecture for sequences of any kind.

\textbf{Slide 16:}

RNNs as Language Models

Processes sequences one word at a time

Predicts the next word in a sequence

\textbf{Uses:}  Current word \& Previous hidden state

Hidden state captures information from all previous words

Overcomes limited context problem of N-gram models & maintains long-term context from the beginning of the

\textbf{Slide 17:}

RNN as language Model

Input sequence x consists of word embeddings (one-hot vectors) of size |V| × 1

Output y is a probability distribution over the vocabulary

\textbf{At each step:} 

Retrieve word embedding using embedding matrix (E)

Combine with previous hidden state

Compute new hidden state

Hidden state is used to generate the output layer

Output is passed through Softmax function

Produces probabilities for all words in the vocabulary

\textbf{Slide 18:}

Training RNNs as language models.

\textbf{Slide 19:}

Generation with RNN-Based Language Models

Generate random sentences using a trained language model to understand its behavior

Start with beginning-of-sentence token (<s>)Sample first word from Softmax output distribution

Feed the sampled word’s embedding as input for the next step

\textbf{Repeat the process:} 

Generate next word Sample from probability distribution

\textbf{Continue until:} 

End-of-sentence token (</s>) is generated

A fixed length limit is reached

\textbf{Slide 20:}

RNN as a Language Model

Steps for using RNN as Language Model

Consider the following sentences

"The cat sat on the mat."

\textbf{Let’s assume our vocabulary consists of 5 unique words:} 

The, cat, sat, on, mat

\textbf{We will manually compute:} 

Word Representation

Forward Pass in RNN

Softmax Output Calculation

Loss Calculation (Cross-Entropy)

Backpropagation (BPTT) Calculation

\textbf{Slide 21:}

\textbf{Step 1:}  Assign Numerical Representation

\textbf{We convert words into one-hot vectors:} 

\textbf{Slide 22:}

\textbf{Our training data consists of input-output pairs:} 

\textbf{Slide 23:}

\textbf{Step 2:}  Forward Pass in RNN

\textbf{An RNN updates its hidden state ht  using the formula:} 

\textbf{where:} 

ht  = hidden state at time t

ht−1  = previous hidden state

xt  = input one-hot vector

Wh = hidden state weight matrix (random initialization)

Wx = input weight matrix (random initialization)

bh  = bias term

tanh⁡(⋅) = activation function

\textbf{Let's assume:} 

Hidden state size = 3

\textbf{Initial hidden state:}  h0=[0,0,0]

\textbf{Random weight matrices:} 

\textbf{Slide 24:}

\textbf{Step 2.1:}  Compute h1  for Input "The"

\textbf{Using the one-hot vector for "The":} 

x1=[1,0,0,0,0]

\textbf{Hidden state calculation:} 

h1=tanh⁡(Wh*h0+Wx*x1+bh)

\textbf{Slide 25:}

\textbf{Step 3:}  Compute Softmax Output

\textbf{The output layer computes:} 

yt=softmax(Wo*ht+bo)

\textbf{where:} 

Wo = output weight matrix

bo  = bias

\textbf{Slide 26:}

\textbf{Step 4:}  Compute Cross-Entropy Loss

\textbf{Given correct output "cat" (index = 1):} 

L=−log P("cat")

L= −log(0.206)=1.58L

\textbf{Step 5:}  Backpropagation Through Time (BPTT)

Gradient calculations are computed for Wx,Wh,Wo using the chain rule.

\textbf{Slide 27:}

\textbf{Step 1:}  Assign Word Embeddings

\textbf{Let’s assume we use a 3-dimensional embedding for each word:} 

So instead of using one-hot vectors, we directly pass the embedding vector to the RNN.

\textbf{Slide 35:}

Autoregressive Generation

The technique of the  generation since the word generated at each  autoregressive generation time step is conditioned on the word selected by the network from the previous step.

\textbf{Slide 36:}

What is Autoregressive RNN?

Autoregressive RNN is a sequence generation model where the next word is predicted based on the previously generated words.

The key characteristic is that the model uses its own generated outputs as inputs for the next step.

It is commonly used in text generation, language modeling, and machine translation.

in actual generation (inference, real time use), the model predicts words one by one without knowing the ground truth at each step, so per-step loss is not useful.

\textbf{Slide 37:}

Step-by-Step Numerical Example

Problem Statement

\textbf{Let's say we train an RNN language model on the sentence:} 

"I love deep learning"

We want to use autoregressive RNN to generate text starting from the seed word "I".

\textbf{Slide 38:}

\textbf{Step 2:}  Initialize RNN Parameters

We assume a simple Recurrent Neural Network (RNN) with the

\textbf{following parameters:} 

Hidden state size = 3

\textbf{Weight matrices:} 

\textbf{Input-to-hidden weights:}  Wx​ (3×3 matrix)

\textbf{Hidden-to-hidden weights:}  Wh​ (3×3 matrix)

\textbf{Hidden-to-output weights:}  Wy​ (3×3 matrix)

\textbf{Biases:}  bh,by​

\textbf{Slide 39:}

\textbf{Step 3:}  Compute Hidden State for First Word ("I")

Hidden State Equation

ht=tanh⁡(Wx*xt+Wh*ht−1+bh)

Initial hidden state (h0​) = [0,0,0]

For word "I" (x1=[0.1,0.3,0.5])

\textbf{Slide 41:}

\textbf{Step 5:}  Continue the Process

The model feeds "love" as input into the RNN.

It generates "deep" next.

Then "learning".

The process continues until a stopping condition is met.

\textbf{Slide 42:}

Conclusion

Autoregressive RNNs generate text one word at a time.

The hidden state captures context from previous words.

The softmax layer selects the most probable next word.

The model recursively feeds predictions back as input.

\textbf{Slide 43:}

Sequence Labeling

In sequence labeling, the network’s task is to assign a label chosen from a small fixed set of labels to each element of a sequence.

Canonical examples of sequence labeling include part-of-speech tagging and named entity recognition

In an RNN approach to sequence labeling, inputs are wordembeddings and the outputs are tag probabilities generated by a softmax layer over the given tag set,

\textbf{Slide 44:}

Part-of-speech tagging as sequence labeling with a simple RNN.

The inputs at each time step are pre-trained word embeddings corresponding to the input tokens.

The RNN block is an abstraction that represents an unrolled simple recurrent network consisting of an input layer, hidden layer, and output layer at each time step, as well as the shared U, V and W weight matrices that comprise the network.

The outputs of the network at each time step represent the distribution over the POS tagset generated by a softmax layer.

\textbf{Slide 45:}

RNNs for Sequence Classification

RNNs is to classify entire sequences rather than the tokens within them

document-level topic classification, spam detection, message routing for customer service applications, and deception detection

To apply RNNs in this setting, the text to be classified is passed through the RNN a word at a time generating a new hidden layer at each time step.

The hidden layer for the final element of the text, hn, is taken to constitute a compressed representation of the entire sequence. In the simplest approach to classification, hn, serves as theinput to a subsequent feedforward network that chooses a class via a SoftMax over the possible classes

\textbf{Slide 46:}

there are no intermediate outputs for the words in the sequence preceding the last element.

there are no loss terms associated with those elements.

Instead, the loss function used to train the weights in the network is based entirely on the final text classification task.

Specifically, the output from the SoftMax output from the feedforward classifier together with a cross-entropy loss drives the training.

The error signal from the classification is backpropagated all the way through the weights in the feedforward classifier through, to its input, and then through to the three sets of weights in the RNN

\textbf{Slide 47:}

Performing Sentiment analysis using RNN

Sentiment analysis using Recurrent Neural Networks (RNNs) involves processing textual data to classify whether the sentiment of a sentence is positive, negative, or neutral.

\textbf{Slide 48:}

\textbf{Step 1:}  Understanding the Sentiment Analysis Task

\textbf{Objective:} 

To classify given text (e.g., movie reviews, tweets) into sentiment categories such as:

Positive sentiment (e.g., "I love this movie")

Negative sentiment (e.g., "I hate this film")

\textbf{Slide 49:}

\textbf{Step 2:}  Preprocessing the Text Data

Before feeding text into an RNN model, we must convert it into numerical form.

Tokenization

\textbf{Definition:}  Splitting a sentence into words (tokens).

\textbf{Example:} 

\textbf{Input:}  "This movie is great"

\textbf{Output tokens:}  ["this", "movie", "is", "great"]

\textbf{Slide 50:}

Convert Words into Word Embeddings

Instead of one-hot encoding, we use word embeddings (like Word2Vec, GloVe, or trainable embeddings) that convert words into dense vector representations.

\textbf{Example:} 

So, the sentence "This movie is great" becomes a matrix of word embeddings.​

\textbf{Slide 51:}

\textbf{Step 3:}  Define RNN Model

RNN processes sequences one word at a time and maintains a hidden state ht​.

3.1 Forward Pass in RNN

\textbf{For each word, the hidden state is updated using:} 

ht=tanh(Wh*ht−1+Wx*xt+bh)

\textbf{where:} 

xt​ = Word embedding of the current word.

ht​ = Hidden state at time t.

Wx​, Wh​, bh​ = Learnable parameters.

\textbf{Slide 52:}

\textbf{Each word in a sentence updates the hidden state as:} 

\textbf{Time Step 1:}  "This"

h1=tanh(Wx*x1+bh)

\textbf{Time Step 2:}  "movie"

h2=tanh(Wh*h1+Wx*x2+bh)

\textbf{Time Step 3:}  "is"

h3=tanh(Wh*h2+Wx*x3+bh)

\textbf{Time Step 4:}  "great"

h4=tanh(Wh*h3+Wx*x4+bh)

At the final time step, the hidden state hT​ represents the entire sentence.

\textbf{Slide 53:}

\textbf{Step 4:}  Sentiment Prediction

The final hidden state hT​ is passed through a fully connected layer and a sigmoid activation function to get the final sentiment score:

y=σ(Wy*hT + by)

\textbf{where:} 

Wy​, by​ = Learnable parameters.

(Sigmoid function converts the output to a probability between 0 and 1).

If y>0.5→ Positive Sentiment

If y<0.5→ Negative Sentiment

\textbf{Slide 54:}

\textbf{Step 5:}  Loss Function \& Backpropagation

\textbf{To train the model, we use Binary Cross-Entropy Loss:} 

\textbf{where:} 

Ytrue = Actual sentiment label (1 for positive, 0 for negative).

y = Predicted probability.

Gradients are computed using Backpropagation Through Time (BPTT), and parameters are updated using Gradient Descent:

where α  (learning rate) is a small value like 0.01.

\textbf{Slide 55:}

\textbf{Step 6:}  Model Training

Train on a dataset of labeled reviews (e.g., IMDB reviews dataset).

Adjust hyperparameters like learning rate, embedding size, and number of RNN units.

Use dropout to prevent overfitting.

\textbf{Slide 56:}

\textbf{Step 7:}  Model Evaluation

Evaluate performance using accuracy, precision, recall, and F1-score.

\textbf{Test with new sentences like:} 

\textbf{"I love this film" → Predicted Sentiment:}  Positive (1)

\textbf{"This movie is bad" → Predicted Sentiment:}  Negative (0)

\textbf{Slide 57:}

\textbf{Step 1:}  Understanding Sentiment Analysis with RNN

\textbf{Objective:}  Predict sentiment (e.g., positive/negative) from a given text.

\textbf{Input:}  Text data (reviews, tweets, etc.).

Output: Sentiment score (e.g., binary classification: positive = 1, negative = 0).

\textbf{Example:} 

\textbf{Input:}  "I love this movie"

\textbf{Output:}  Positive sentiment (1)

\textbf{Slide 58:}

\textbf{Step 2:}  Preprocessing the Data

Before feeding text into an RNN model, we must convert it into numerical form.

2.1 Tokenization

Tokenization splits the text into individual words:Sentence: "This movie is great"Tokens: ["this", "movie", "is", "great"]

2.2 Word Embeddings (Instead of One-Hot Encoding)

Each word is mapped to a fixed-size vector using pretrained embeddings (like Word2Vec, GloVe) or a trainable embedding layer.

\textbf{Slide 60:}

\textbf{Step 3:}  Forward Pass in RNN\

An RNN processes each word sequentially and updates a hidden state.

ht=tanh⁡(Wh*ht−1+Wx*xt+bh)

\textbf{where:} 

xt​ = Word embedding at time t.

ht​ = Hidden state at time t.

Wx​, Wh​, bh​ = Learnable parameters.

\textbf{Let's initialize:} 

h0=[0,0,0]

\textbf{Slide 65:}

Stacked RNN

Stacked recurrent networks.

The output of a lower level serves as the input to higher levels with the output of the last network serving as the final output

The optimal number of stacked RNNs is specific to each application and to each training set.

However, as the number of stacks is increased the training costs rise quickly.

\textbf{Slide 66:}

\textbf{RNNs:}  Vanishing/exploding gradients

Product of so many derivatives of tanh functions (<1) will lead to exponentially small values (~0). Vanishing Gradient problem

Product of so many derivatives of tanh functions (>1) will lead to exponentially large values. Exploding Gradient problem

Exploding Gradient problem can be simply solved by clipping the gradients

However, vanishing gradients is a complex issue and disallow capturing long-range dependencies

Proper initialization of the W matrix can reduce the effect of vanishing gradients

Use ReLU (derivative either 0 or 1) instead of sigmoid or tanh activation functions

Use Long Short-Term Memory (LSTM) or Gated Recurrent Unit (GRU) architectures

\textbf{Slide 67:}

\textbf{Managing Context in RNNs:}  LSTMs and GRUs

to make use of information distant from the current point of processing is difficult while training RNN.

the information encoded in hidden states tends to be local, more relevant to the most recent parts of the input sequence and recent decisions.

\textbf{Slide 68:}

Long short-term memory (LSTM)

A type of RNN

The differences are the operations within the LSTM’s cells

What to keep and what to forget?

\textbf{Slide 69:}

\textbf{LSTM:}  The main idea

The cell state and various gates

The cell state act as a transport highway that transfers relative information all the way down the sequence chain.

The gates decide which information is allowed on the cell state

Gates contains sigmoid activations

\textbf{Forget info:}  if sigmoid activation ~0

\textbf{Keep info:}  if sigmoid activation ~1

\textbf{Slide 70:}

Consider the following example

step by step and calculation of  all intermediate values for an LSTM cell when processing the input sequence "Movie is good" using given weights and biases.

\textbf{Slide 71:}

Consider the following example

step by step and calculation of  all intermediate values for an LSTM cell when processing the input sequence "Movie is good" using given weights and biases.

\textbf{Slide 72:}

\textbf{Step 2:}  LSTM Forward Pass for "Movie"

We compute all gate activations for time step t = 1.

\textbf{Slide 73:}

\textbf{LSTM:}  Forget gate

This gate decides what information should be thrown away or kept.

Information from previous hidden state and the current input is passed through the sigmoid function.

Values come out between 0 and 1. The closer to 0 means to forget, and the closer to 1 means to keep.

\textbf{Slide 74:}

\textbf{LSTM:}  Input gate

First, the previous hidden state and current input is passed into a sigmoid function

Also pass the hidden state and current input into the tanh function to squish values between -1 and 1.

Multiply the tanh output with the sigmoid output

Sigmoid output will decide which information is important to keep from the tanh output

\textbf{Slide 76:}

\textbf{LSTM:}  Cell state

The cell state gets pointwise multiplied by the forget vector.

Take the output from the input gate and do a pointwise addition which updates the cell state to new values

\textbf{Slide 77:}

\textbf{LSTM:}  Output gate

The output gate decides what the next hidden state should be

First, we pass the previous hidden state and the current input into a sigmoid function.

Then we pass the newly modified cell state to the tanh function.

Multiply the tanh output with the sigmoid output

\textbf{Slide 78:}

Example

Traditional RNN Fails to predict because of short term memory

\textbf{Slide 79:}

Example

\textbf{Slide 80:}

Example

\textbf{Slide 81:}

\textbf{LSTM:}  The operations

Forget gate

\textbf{Slide 82:}

\textbf{LSTM:}  The operations

Input gate

\textbf{Slide 83:}

\textbf{LSTM:}  Updating weights

\textbf{Slide 84:}

\textbf{LSTM:}  The operations

Output gate

\textbf{Slide 85:}

\textbf{LSTM:}  Updating weights

\textbf{Slide 86:}

\textbf{LSTM:}  Updating weights

\textbf{Slide 87:}

\textbf{LSTM:}  Updating weights

Gradient behaves similar to the forget gate

\textbf{Slide 88:}

\textbf{LSTM:}  Updating weights

This can be controlled initially using +ve bias term!

\textbf{Slide 89:}

Gated Recurrent Unit (GRU)

New generation RNNs

GRU’s got rid of the cell state and used the hidden state to transfer information

Little speedier to train than LSTMs

But not a clear winner

Update Gate: decides what information to throw away and what new information to add

\textbf{Reset Gate:}  decides how much past information to forget.

\textbf{Slide 90:}

GRU Mathematical Explanation to the Diagram

\textbf{A GRU cell consists of:} 

Reset Gate (rt , red sigmoid in the left part of the diagram)

Update Gate (zt , red sigmoid in the center)

Candidate Hidden State (h~t)

New Hidden State (ht , final output)

\textbf{Slide 92:}

Purpose of the Reset Gate

The reset gate (rt​) determines how much of the past hidden state (ht−1​) should be forgotten before computing the new candidate hidden state (ht~​).

If rt is close to 0, the past hidden state is ignored, meaning the GRU relies mostly on the current input.

If rt is close to 1, the past hidden state is preserved and influences the new candidate hidden state.

\textbf{Slide 93:}

For example, if rt=

The first value 0.42 means that for the first hidden unit, about 42% of the previous hidden state is used in computing h~t​, while 58% is ignored.

The second value 0.55 means that for the second hidden unit, 55% of the past state is preserved, and 45% is forgotten.

\textbf{Slide 95:}

Purpose of the Candidate Hidden State

The candidate hidden state (h~t​) represents the new memory that could be used to update the final hidden state.

It is computed based on the current input xt and a modified version of the previous hidden state ht−1​, which is influenced by the reset gate rt​.

If rt​ is close to 0, the past hidden state is largely ignored. If rt​ is close to 1, the past hidden state has more influence.

\textbf{Slide 98:}

Understanding 1−z  and zt

The update gate zt​ plays a critical role in determining how much of the previous hidden state ht−1​ should be kept and how much of the newly computed candidate hidden state h~t​ should be used.

\textbf{The final hidden state ht  is computed as:} 

\textbf{This equation is a weighted sum where:} 

zt​ controls how much of h~t (the new memory) is retained.

1−zt controls how much of ht−1​ (the old memory) is kept.

\textbf{Slide 100:}

Interpretation of 1−zt​ and zt

1−ztacts as a forgetting factor for the old state ht−1​.

If ztis small, then 1−zt is large, meaning more of the old state is preserved.

If zt​ is large, then 1−zt=​ is small, meaning more of the old state is discarded.

zt​ determines how much of h~t​ (new memory) should be incorporated into the hidden state.

If zt​ is close to 1, the new memory is fully used, making the network update its knowledge rapidly.

If zt​ is close to 0, the network relies mostly on past memory, preventing unnecessary updates.

This mechanism makes the GRU effective at handling long-term dependencies in sequential data.

\textbf{Slide 101:}

Example

\textbf{Slide 102:}

GRU

\textbf{Slide 103:}

Example 2

\textbf{Slide 104:}

Different types of RNNs

\textbf{Slide 105:}

\textbf{one-to-one:}  Binary Classification [single images/words → single output]

\textbf{one-to-many:}  Image Captioning [image → Sequence of words](variability sized)

many-to-one: Sentiment Classification [sequence of words (variable sized) → single sentiment]

many-to-many (non-matching): Machine Translation [sequence of words → sequence of words], where both sequences have variable length and these lenghts may be difference

\textbf{many-to-many (matching):}  Video classification on a…


\subsection{2. Bidirectional RNN}

\textbf{Slide 1:}

Bi-directional Recurrent Neural Network

In order for the network to use information from both the past and future context in its predictions, BRNNs process input sequences in both the forward and backward directions.

This is the main distinction between BRNNs and conventional recurrent neural networks.

\textbf{Slide 2:}

BRNN

A BRNN has two distinct recurrent hidden layers, one of which processes the input sequence forward and the other of which processes it backward.

After that, the results from these hidden layers are collected and input into a prediction-making final layer.

Any recurrent neural network cell, such as Long Short-Term Memory (LSTM) or Gated Recurrent Unit, can be used to create the recurrent hidden layers.

\textbf{Slide 3:}

BRNN functions

The BRNN functions similarly to conventional recurrent neural networks in the forward direction, updating the hidden state depending on the current input and the prior hidden state at each time step.

The backward hidden layer, on the other hand, analyses the input sequence in the opposite manner, updating the hidden state based on the current input and the hidden state of the next time step.

Compared to conventional unidirectional recurrent neural networks, the accuracy of the BRNN is improved since it can process information in both directions and account for both past and future contexts.

Because the two hidden layers can complement one another and give the final prediction layer more data, using two distinct hidden layers also offers a type of model regularization.

\textbf{Slide 4:}

Gradient Updating

In order to update the model parameters, the gradients are computed for both the forward and backward passes of the backpropagation through the time technique that is typically used to train BRNNs.

The input sequence is processed by the BRNN in a single forward pass at inference time, and predictions are made based on the combined outputs of the two hidden layers.

\textbf{Slide 6:}

Working of Bidirectional Recurrent Neural Network

Inputting a sequence: A sequence of data points, each represented as a vector with the same dimensionality, are fed into a BRNN. The sequence might have different lengths.

Dual Processing:  Both the forward and backward directions are used to process the data. Based on the input at that step and the hidden state at step t-1, the hidden state at time step t is determined in the forward direction. The input at step t and the hidden state at step t+1 are used to calculate the hidden state at step t in a reverse way.

Computing the hidden state: A non-linear activation function on the weighted sum of the input and previous hidden state is used to calculate the hidden state at each step. This creates a memory mechanism that enables the network to remember data from earlier steps in the process.

Determining the output: A non-linear activation function is used to determine the output at each step from the weighted sum of the hidden state and a few output weights. This output has two options: it can be the final output or input for another layer in the network.

Training: The network is trained through a supervised learning approach where the goal is to minimize the discrepancy between the predicted output and the actual output. The network adjusts its weights in the input-to-hidden and hidden-to-output connections during training through backpropagation.

\textbf{Slide 7:}

\textbf{To calculate the output from an RNN unit, we use the following formula:} 

\textbf{Slide 8:}

The training of a BRNN is like backpropagation through a time algorithm. BPTT algorithm works as follows:

\textbf{Slide 9:}

Applications of Bidirectional Recurrent Neural Network

Bi-RNNs have been applied to various natural language processing (NLP) tasks, including:

Sentiment Analysis: By considering both the prior and subsequent context, BRNNs can be utilized to categorize the sentiment of a particular sentence.

Named Entity Recognition: By considering the context both before and after the stated thing, BRNNs can be utilized to identify those entities in a sentence.

Part-of-Speech Tagging: The classification of words in a phrase into their corresponding parts of speech, such as nouns, verbs, adjectives, etc., can be done using BRNNs.

Machine Translation: BRNNs can be used in encoder-decoder models for machine translation, where the decoder creates the target sentence, and the encoder analyses the source sentence in both directions to capture its context.

Speech Recognition: When the input voice signal is processed in both directions to capture the contextual information, BRNNs can be used in automatic speech recognition systems.

\textbf{Slide 10:}

Advantages of Bidirectional RNN

Context from both past and future: With the ability to process sequential input both forward and backward, BRNNs provide a thorough grasp of the full context of a sequence. Because of this, BRNNs are effective at tasks like sentiment analysis and speech recognition.

Enhanced accuracy: BRNNs frequently yield more precise answers since they take both historical and upcoming data into account.

Efficient handling of variable-length sequences: When compared to conventional RNNs, which require padding to have a constant length, BRNNs are better equipped to handle variable-length sequences.

Resilience to noise and irrelevant information: BRNNs may be resistant to noise and irrelevant data that are present in the data. This is so because both the forward and backward paths offer useful information that supports the predictions made by the network.

Ability to handle sequential dependencies: BRNNs can capture long-term links between sequence pieces, making them extremely adept at handling complicated sequential dependencies.

\textbf{Slide 11:}

Disadvantages of Bidirectional RNN

Computational complexity: Given that they analyze data both forward and backward, BRNNs can be computationally expensive due to the increased amount of calculations needed.

Long training time: BRNNs can also take a while to train because there are many parameters to optimize, especially when using huge datasets.

Difficulty in parallelization: Due to the requirement for sequential processing in both the forward and backward directions, BRNNs can be challenging to parallelize.

Overfitting: BRNNs are prone to overfitting since they include many parameters that might result in too complicated models, especially when trained on short datasets.

Interpretability: Due to the processing of data in both forward and backward directions, BRNNs can be tricky to interpret since it can be difficult to comprehend what the model is doing and how it is producing predictions.


\subsection{Gated Recurrent Unit (3)}

\textbf{Slide 1:}

Gated Recurrent Unit

\textbf{Slide 8:}

Element Wise Multiplication

\textbf{Slide 10:}

The Basic Idea

It uses gating mechanisms to selectively update the hidden state of the network at each time step.

The gating mechanisms are used to control the flow of information in and out of the network.

The GRU has two gating mechanisms, called the reset gate and the update gate.

The reset gate determines how much of the previous hidden state should be forgotten, while the update gate determines how much of the new input should be used to update the hidden state.

The output of the GRU is calculated based on the updated hidden state.

\textbf{Slide 11:}

The equations used to calculate the reset gate, update gate, and hidden state of a GRU are as follows:

where W\_r, W\_z, and W\_h are learnable weight matrices,

x\_t is the input at time step t,

h\_{t-1} is the previous hidden state,

and h\_t is the current hidden state.

\textbf{Slide 12:}

Collective information passed onto the next Gated Recurrent Unit

\textbf{The different gates of a GRU are as described below:} -

\textbf{Update Gate(z):} 

It determines how much of the past knowledge needs to be passed along into the future.

It is analogous to the Output Gate in an LSTM recurrent unit.

\textbf{Reset Gate(r):} 

It determines how much of the past knowledge to forget.

It is analogous to the combination of the Input Gate and the Forget Gate in an LSTM recurrent unit.

\textbf{Current Memory Gate(    ):} 

It is often overlooked during a typical discussion on Gated Recurrent Unit Network. It is incorporated into the Reset Gate just like the Input Modulation Gate is a sub-part of the Input Gate and is used to introduce some non-linearity into the input and to also make the input Zero-mean. Another reason to make it a sub-part of the Reset gate is to reduce the effect that previous information has on the current information that is being passed into the future.

\textbf{Slide 13:}

\textbf{Working of a Gated Recurrent Unit:} 

Take input the current input and the previous hidden state as vectors.

Calculate the values of the three different gates by following the steps given below:-

For each gate, calculate the parameterized current input and previously hidden state vectors by performing element-wise multiplication (Hadamard Product) between the concerned vector and the respective weights for each gate.

Apply the respective activation function for each gate element-wise on the parameterized vectors. Below given is the list of the gates with the activation function to be applied for the gate.

\textbf{Slide 14:}

Step 3

The process of calculating the Current Memory Gate is a little different.

First, the Hadamard product of the Reset Gate and the previously hidden state vector is calculated. Then this vector is parameterized and then added to the parameterized current input vector.

\textbf{Slide 15:}

Step 4

To calculate the current hidden state, first, a vector of ones and the same dimensions as that of the input is defined.

This vector will be called ones and mathematically be denoted by 1.

First, calculate the Hadamard Product of the update gate and the previously hidden state vector.

Then generate a new vector by subtracting the update gate from ones and then calculate the Hadamard Product of the newly generated vector with the current memory gate.

Finally, add the two vectors to get the currently hidden state vector.

\textbf{Slide 16:}

Training

the training process for a GRU network is also diagrammatically like that of a basic Recurrent Neural Network and differs only in the internal working of each recurrent unit.

The Back-Propagation Through Time Algorithm for a Gated Recurrent Unit Network is like that of a Long Short Term Memory Network and differs only in the differential chain formation.


\section{Module 5: Morphological and Syntactic Analysis}


\subsection{1_Morphological analysis (1) (2)}

\textbf{Slide 1:}

Morphology

Morphology is the study of words

Study's internal structure of words

How words change their form to generate new word

Different role they play in sentence ,strictly following linguistic rule.

\textbf{Slide 2:}

Morphemes

words are built from smaller meaningful grammatical units called morphemes.

dogs

2 morphemes, ‘dog’ and ‘s’

‘s’ is a plural marker on nouns

\textbf{Slide 3:}

Morpheme

A Smallest unit carries meaning

While doing analysis of a word, we split the given word into parts.

Each part is called a morpheme.

\textbf{Eg.:} 

uneducated –>

un + educate + d.

it has three morphemes.

\textbf{Slide 4:}

Free Morpheme

Can appear as a word by itself; often can combine with other morphemes too.

house (house-s), walk (walk-ed), of, the, or

Independent/can stand by themselves as single words.

When combined with bound morphemes

free morphemes are called stems/root.

\textbf{Slide 5:}

Lexical  Morpheme

Visualize/Content word

Nouns –Glass

Adjectives -Black,

Verbs-Dance,

Adverb - beautifully.

‘open' class of morphemes because we can add new words to the language easily

Eg.(girl, play, google, e-mail, blog).

\textbf{Slide 6:}

Grammatical Morpheme

Also called as Functional morphemes

consist of functional/Grammer words

Conjunctions- And/Or

Prepositions-On/Under/At

Articles-An/The

Pronouns- She/He.

Limited words

This is a ‘closed' class of morphemes

\textbf{Slide 7:}

Bound  Morpheme

Cannot appear as a word by itself.

-s (dog-s), -ly (quick-ly), -ed (walk-ed)

cannot stand alone

Not complete in themselves

Depends on free Morpheme

Affixes are bound morpheme

typically attached to free morpheme.

They can prefixes, Infix and suffixes (re-, un-, dis-, pre-, -ness, -less, -ly).

\textbf{Slide 8:}

Prefix

Attached before lexical morpheme

Be-come

Un-happy

\textbf{Slide 9:}

Infix

unladylike

3 morphemes

un- ‘not’

lady ‘well-behaved woman’

-like ‘having the characteristic of’

\textbf{Slide 10:}

Suffix

Words attached after root

Teach-er

Use-less

\textbf{Slide 11:}

Suffix

8 Suffixes

Noun

Plural –   Book->Books

Possessional- Ram’s

Adjective

er –   Tall->Taller

est-   Tall->Tallest

Verb

S=     Play->Plays

ed=   Play->Played

Ing=  Play->Playing

en=   Brake->Broken

\textbf{Slide 12:}

Derivational  Morpheme

make new words of a different grammatical category from a stem

Class Changing

Care	         careful

(Verb)            (Adjective)

Care	   	  carelessly

(Verb)                 (Adverb)

Class Maintaining

Child		Childhood

(Noun)                 (Noun)

\textbf{Slide 13:}

Inflectional  Morpheme

Morpheme when attached to root doesn’t change its class

There are eight inflectional morphemes in English.

They are all suffixes.

\textbf{Slide 14:}

Regular Inflectional  Morpheme

Refers to those inflections which follows a standard pattern

Eg.

Dance+ed->Danced

Walk+ed->Walked

Boy+s->Boys

\textbf{Slide 15:}

Irregular Inflectional  Morpheme

Refers to those inflections which does not follow a standard pattern

Eg.

Wife+s - >  Wives

Child+s  ->  Children

Mouse+s   ->   Mice

\textbf{Slide 16:}

Completely change the morpheme

Occurrence of phonemically unrelated morpheme

Eg.

Go+ed->Went

Good+er->Better

am+ed->was

Suppletion Inflectional  Morpheme

\textbf{Slide 17:}

Identify Inflectional or Derivational Morphology

I jumped into the puddle this morning

This is unbelievable

I have john’s umbrella

Emma goes to school

She is working carelessly

\textbf{Slide 18:}

Identify Inflectional or Derivational Morphology

I jumped into the puddle this morning

This is unbelievable

I have john’s umbrella

Emma goes to school

She is working carelessly

Inflectional Morphology

Derivational Morphology

Inflectional Morphology

Inflectional Morphology

Derivational Morphology

\textbf{Slide 20:}

Regular Expressions

\textbf{Slide 21:}

Some Questions on Regular Expression.....

\textbf{Slide 22:}

What is Regular Expression?

What does the command ab+c search for?

ac,abc,abbc, and so on

ab,abc,abcc and so on

abc,abbc,abbbc and so on

None of the above

The number of matches does the command a{1,3} give with the string aabbaaaa?

3

2

1

4

\textbf{Slide 23:}

What is the output of the below code?

print(re.match(‘On’,”one”))

'e'

Match Object

‘n’

None

What does the sequence \D finds the match?

Decimal digits

Non-decimal digits

Division

None of the above

Which of the following command is used to search a match for 1,2,3,4?

[1-4]

(1-3)

[1234]

Both a and c

\textbf{Slide 24:}

What is the output of the below code?

re.sub(‘a’,’u’,’aeiou!’)

'ueiou!'

'eiou!'

'eio!'

None of the above

\textbf{Slide 26:}

What is Regular  Expression?

A  Formal Language for specifying text search strings

A regular expression is an algebraic notation for characterizing a set  of strings.

Regular expression search function will search through the corpus, returning all texts that match the pattern or the first match.

The corpus can be a single document or a collection

\textbf{Slide 27:}

Regular expressions

The simplest kind of regular expression is sequence of simple characters

A formal language for specifying text strings

How can we search for any of these?

woodchuck

woodchucks

Woodchuck

Woodchucks

[Ww]oodchucks?

Regular expressions are case sensitive

All the above are different regular expression

\textbf{Slide 28:}

Regular Expression

\textbf{Slide 29:}

\textbf{Regular Expressions:}  Disjunctions

Letters inside square brackets []

Ranges [A-Z]

\textbf{Slide 30:}

\textbf{Regular Expressions:}  Negation in Disjunction

Negations [^Ss]

Carat means negation only when first in []

\textbf{Slide 31:}

\textbf{Regular Expressions:}  More Disjunction

Woodchuck is another name for groundhog!

The pipe | for disjunction

\textbf{Slide 32:}

\textbf{Regular Expressions:}  ? *+.

Stephen C Kleene

Kleene *,   Kleene +

\textbf{Slide 33:}

\textbf{Regular Expressions:}  Anchors  ^   \$

993, 99

If your pattern is 99 it will be matched to both above string if you want to match only 99 you can use \b99\b

\textbf{Slide 36:}

Example

Find me all instances of the word “the” in a text.

the

Misses capitalized examples

[tT]he

Incorrectly returns other or theology

[^a-zA-Z][tT]he[^a-zA-Z]

\textbf{Slide 37:}

Errors

\textbf{The process we just went through was based on fixing two kinds of errors:} 

Matching strings that we should not have matched (there, then, other)

False positives (Type I errors)

Not matching things that we should have matched (The)

False negatives (Type II errors)

\textbf{Slide 38:}

Errors cont.

In NLP we are always dealing with these kinds of errors.

Reducing the error rate for an application often involves two antagonistic efforts:

Increasing accuracy or precision (minimizing false positives)

Increasing coverage or recall (minimizing false negatives).

\textbf{Slide 41:}

Summary

Regular expressions play a surprisingly large role

Sophisticated sequences of regular expressions are often the first model for any text processing text

For hard tasks, we use machine learning classifiers

But regular expressions are still used for pre-processing, or as features in the classifiers

Can be very useful in capturing generalizations

41

\textbf{Slide 42:}

\textbf{More Regular Expressions:} 

Substitutions and ELIZA

\textbf{Slide 43:}

An important use of regular expressions is in substitutions.

For example, the substitution operator s/regexp1/pattern/ used in Python and in Unix commands like vim or sed allows a string characterized by a regular expression to be replaced by another string:

\textbf{Slide 44:}

Substitutions

\textbf{Substitution in Python and UNIX commands:} 

s/regexp1/pattern/

\textbf{e.g.:} 

s/colour/color/

\textbf{Slide 45:}

Capture Groups

\textbf{Say we want to put angles around all numbers:} 

the 35 boxes  the <35> boxes

Use parens () to "capture" a pattern into a numbered register (1, 2, 3…)

Use \1  to refer to the contents of the register

s/([0-9]+)/<\1>/

\textbf{Slide 46:}

\textbf{Capture groups:}  multiple registers

/the (.*)er they (.*), the \1er we \2/

Matches

the faster they ran, the faster we ran

But not

the faster they ran, the faster we ate

Sham went to Paris with Ram, Ram loves Paris."

/sham went to (Paris) with (Ram),\1 loves \2

\textbf{Slide 47:}

The harder he worked, the harder she worked.

the faster they ran, the faster we ran

The sooner they arrived, the sooner we arrived.

\textbf{Slide 48:}

But suppose we don't want to capture?

\textbf{Parentheses have a double function:}  grouping terms, and capturing

\textbf{Non-capturing groups:}  add a ?: after paren:

\textbf{/(?:} some|a few) (people|cats) like some \1/

matches

some cats like some cats

but not

some cats like some some

\textbf{Slide 49:}

\textbf{Simple Application:}  ELIZA

Early NLP system that imitated a Rogerian psychotherapist

Joseph Weizenbaum, 1966.

\textbf{Uses pattern matching to match, e.g.,:} 

“I need X”

and translates them into, e.g.

“What would it mean to you if you got X?

\textbf{Slide 50:}

\textbf{Simple Application:}  ELIZA

Men are all alike.IN WHAT WAY

They're always bugging us about something or other. CAN YOU THINK OF A SPECIFIC EXAMPLE

Well, my boyfriend made me come here.YOUR BOYFRIEND MADE YOU COME HERE

He says I'm depressed much of the time.I AM SORRY TO HEAR YOU ARE DEPRESSED

\textbf{Slide 51:}

How ELIZA works

s/.* I’M (depressed|sad) .*/I AM SORRY TO HEAR YOU ARE \1/

s/.* I AM (depressed|sad) .*/WHY DO YOU THINK YOU ARE \1/

s/.* all .*/IN WHAT WAY?/

s/.* always .*/CAN YOU THINK OF A SPECIFIC EXAMPLE?/

\textbf{Slide 52:}

Write the regular expression for the following

Write the regular expression that matches a string that has an a followed by zero or more b's

Write the regular expression for given language  L = {E, a, aa, b,bb,ab,ba, aba, bab,.....},any combination of a and b.

Write the regular expression for the language  L = {a, aba, aab, aba, aaa, abab, .....}

Write the regular expression for the language L = {a, aa,aaa, ....}

Write the regular expression for the language L = {E, 0, 1,00, 11,10,100,.....}

Write the regular expression for the language accepting all the string which are starting with 1 and ending with 0, over ∑ = {0, 1}.

\textbf{Slide 53:}

Write the regular expression for the following

Write the regular expression that matches a string that has an a followed by zero or more b's R=ab*

Write the regular expression for given language  L = {E, a, aa, b,bb,ab,ba, aba, bab,.....},any combination of a and b.  Answer:-Solution

The regular expression will be −(a + b)*

Write the regular expression for the language  L = {a, aba, aab, aba, aaa, abab, .....} R = {a + ab}*

Write the regular expression for the language L = {a, aa,aaa, ....} R = a+

Write the regular expression for the language L = {E, 0, 1,00, 11,10,100,.....}

R = (1* O*)

Write the regular expression for the language accepting all the string which are starting with 1 and ending with 0, over ∑ = {0, 1}. R = 1 (0+1)* 0

\textbf{Slide 54:}

Finite State Automata

Finite Automata(FA) is the simplest machine to recognize patterns.

The finite automata or finite state machine is an abstract machine that has five elements or tuples.

It has a set of states and rules for moving from one state to another, but it depends upon the applied input symbol.

Basically, it is an abstract model of a digital computer.

Finite automata have two states, Accept state or Reject state.

When the input string is processed successfully, and the automata reached its final state, then it will accept.

\textbf{Slide 55:}

\textbf{The above figure shows the following features of automata:} 

Input

Output

States of automata

State relation

Output relation

\textbf{Slide 56:}

\textbf{A Finite Automata consists of the following:} 

\textbf{Q :}  Finite set of states.

\textbf{Σ :}  set of Input Symbols.

\textbf{q :}  Initial state.

\textbf{F :}  set of Final States.

\textbf{δ :}  Transition Function.

\textbf{Slide 57:}

\textbf{Deterministic Finite Automata (DFA):} 

DFA consists of 5 tuples {Q, Σ, q, F, δ}.

\textbf{Q :}  set of all states.

\textbf{Σ :}  set of input symbols. ( Symbols which machine takes as input )

\textbf{q :}  Initial state. ( Starting state of a machine )

\textbf{F :}  set of final state.

\textbf{δ :}  Transition Function, defined as δ : Q X Σ --> Q.

In a DFA, for a particular input character, the machine goes to one state only.

A transition function is defined on every state for every input symbol.

Also in DFA null (or ε) move is not allowed, i.e., DFA cannot change state without any input character.

\textbf{Slide 58:}

For example, below DFA with Σ = {0, 1} accepts all strings ending with 0.

\textbf{Slide 64:}

State Transition Diagram

\textbf{Slide 71:}

Draw a deterministic  finite automate which accept 00 and 11 at the end of a string containing 0, 1 in it, e.g., 01010100 but not 000111010.

\textbf{Slide 73:}

FSA as language Recognizer

Recognition is the process of determining if a string should be accepted by a machine •

Or… it’s the process of determining if a string is in the language, we’re defining with the machine

Or… it’s the process of determining if a regular expression matches a string

• Those all amount the same thing in the end

\textbf{Slide 74:}

Recognition

Simply a process of starting in the start state

Examining the current input

Consulting the table

Going to a new state and updating the input pointer.

Until you run out of input

\textbf{Slide 81:}

\textbf{Finite-state can capture the generalization here:} 

Eg.

I eat Sushi

Ram like Mango

Noun+ Verb Noun+

\textbf{Slide 82:}

Language is recursive

§ the ball

§ the ball in the garden

§ the ball in the garden behind the house

§ the ball in the garden behind the house     next to school

§ the ball

§ the big ball

§ the big, red ball

§ the big, red, heavy ball

\textbf{Slide 84:}

Morphological parsing

Morphological parsing is to find the lexical form of a word from its surface form.

Given the input, for example, cats, we would like to produce cat +N +PL.

Stem-----Prefix+Stem+Suffix

\textbf{Slide 85:}

Morphological parsing

Two-level morphology,

– Representing a word as a correspondence between

a lexical level

• Representing a simple concatenation of morphemes making up a word, and

– The surface level

• Representing the actual spelling of the final word.

• Morphological parsing is implemented by building mapping rules that maps letter sequences like cats on the surface level  into morpheme and features sequence like cat +N +PL on the lexical level

\textbf{Slide 86:}

Morphological Parsing

–cats ------------	cat +N +PLU

–cat -------------	cat +N +SG

–goose ------------goose +N +SG

–geese ------------goose +N +PLU

–catch ------------catch +V

–caught -----------catch +V +PAST

Surface Form

Lexical Form

\textbf{Slide 87:}

Parts of A Morphological Processor

\textbf{•Lexicon:}  The list of stems

with basic information about

categories (noun, verb, adjective, …)

sub-categories (regular noun, irregular noun, …)

•Morphotactics: Explains ordering -which classes of morphemes can follow other classes of morphemes inside a word.

•Orthographic Rules (Spelling Rules): These spelling rules are used to model changes that occur in a word (normally when two morphemes combine).

\textbf{Slide 88:}

Lexicon

•A lexicon is a repository for words (stems).

•They are grouped according to their main categories.

–noun, verb, adjective, adverb, …

•They may be also divided into sub-categories.

–regular-nouns, irregular-singular nouns, irregular-plural nouns

•The simplest way to create a morphological parser, put all possible words (together with its inflections) into a lexicon.

\textbf{Slide 89:}

Morphotactic

In a natural language, a word may have many morphemes attached to it as affixes (prefix, suffix, infix).

Morphotactic is about placing morphemes with stem to form a meaningful word.

It gives the details about which classes of morpheme can follow which other classes of morpheme.

For example, in English language the plural morpheme (s or es) follows the stem.

\textbf{Slide 90:}

Morphotactics

\textbf{Slide 91:}

Orthographic Rules

It refers to the spelling rules used in a particular language to model the spelling changes that occur in a word.

For example, when a stem ‘cook’ is attached with the morpheme ‘ing’, it becomes a new valid word with the spelling ‘cooking’.

What happened here was the concatenation of two strings.

This is not true for all words. If we attach the morpheme ‘ing’ with the word ‘care’, the spelling changes to ‘caring’ rather than ‘careing’.

In this case, it was not just concatenation, but it also involved removal of the letter ‘e’.

\textbf{Slide 94:}

Finite state transducers

A finite state transducer essentially is a finite state automaton that works on two (or more) tapes. The most common way to think about transducers is as a kind of ``translating machine''.

They read from one of the tapes and write onto the other. This, for instance, is a transducer that translates as into bs:

\textbf{Slide 95:}

Finite state transducers

A finite state transducer (FST) is a finite state machine where transitions are conditioned on a pair of symbols

The machine moves between the states based on input symbol, while it outputs the corresponding output symbol

An FST encodes a relation, a mapping from a set to another . The relation defined by an FST is called a regular (or rational) relation

\textbf{Slide 99:}

Two-Level Morphology

Two-level morphology represents the correspondence between lexical and surface levels.

•We use a finite-state transducer to find mapping between these two levels.

\textbf{•A FST is a two-tape automaton:} 

–Reads from one tape and writes to other one.

•For morphological processing, one tape holds lexical representation, the second one holds the surface form of a word.

\textbf{Slide 100:}

Morphological segmentation (or Stemming)

Taking a surface input and breaking it down into its morphemes

cat- 		cat+N+SL

cats- 		cat+N+PL

Plays- 	Play+V+Singular

Played- 	Play+V+ PastParticipat

foxes – 	fox+N+PL

\textbf{Slide 102:}

\#

+Sg

\textbf{Slide 103:}

\#

+Sg

\textbf{Slide 104:}

For each spelling rule we will have a FST, and these FSTs run parallel.

\textbf{Slide 105:}

Orthographic Rules

•For each spelling rule we will have a FST, and these FSTs run parallel.

\textbf{•We represent these rules using two-level morphology rules:} 

a => b / c \_\_ d

rewrite a as b when it occurs between c and d.

\textbf{English Spelling Rules:} 

––E insertion --e added after s, z, x, ch, sh  before s --watch/watches

\textbf{Slide 106:}

Stemming

Stemming is suffix stripping operation

Process of reducing word into its base form (Root form/stem form)

This is achieved by cutting of begging or end of the word

\textbf{Eg:} 

Playing-    Play

Boys-         Boy

Popular Algorithm is Porter Stemmer

\textbf{Slide 107:}

Porter Stemmer

– Five sets of rules, applied in order

– Within each set, if more than one of the rules can apply, only the one with the longest matching suffix (S1) is followed

\textbf{Advantage:}   easy to see understand, easy to implement.

\textbf{Slide 108:}

Convention

Consonant( C ): other than A, E, I, O or U, and other than    Y preceded by a consonant.

So in TOY the consonants are T and Y

\textbf{Vowel(V) :}  Any other letter.

\textbf{Slide 109:}

\textbf{Any word in English has this forms:} 

[C](VCm)[V]

[]denotes arbitrary presence of content

C-Consonant

V-Vowel

m will be called the measure of any word or word part

m=0 TR, EE, TREE, Y, BY.

m=1 TROUBLE, OATS, TREES, IVY.

m=2 TROUBLES, PRIVATE, OATEN, ORRERY.

\textbf{Slide 110:}

RULES

The rules for removing a suffix will be given in the form		(condition) S1 -> S2if a word ends with the suffix S1 will be replaced by S2, if it satisfies the given condition.

e.g.   (m > 1) EMENT -> ЄHere S1 is 'EMENT' and S2 is null.

REPLACEMENT to REPLAC,

since REPLAC is a word part for which m = 2.

\textbf{Slide 111:}

E.g. (m>1 or *S)

\textbf{Slide 112:}

\textbf{Step 1a :} 

\textbf{sses -> ss (Example :}  caresses -> caress)

\textbf{ies -> i (Example :}  ponies -> poni ; ties -> ti)

\textbf{ss -> ss (Example :}  caress -> caress)

\textbf{s ->Є (Example :}  cats -> cat)

PORTER STEMMER

\textbf{Slide 113:}

\textbf{Step 1b :} 

\textbf{(m>0) eed -> ee (Example :}  agreed -> agree; feed -> feed )

\textbf{(*v*) ed -> є (Example :}  plastered -> plaster ; bled -> bled)

\textbf{(*v*) ing -> є (Example :}  motoring -> motor ; sing -> sing)

\textbf{s -> є (Example :}  cats -> cat)

\textbf{Slide 114:}

If the second or third of the rules in Step 1b is successful, the following is done: Cleaning Step

\textbf{at -> ate (Example :}  conflat(ed) -> conflate)

\textbf{bl -> ble (Example :}  troubl(ed) -> trouble)

\textbf{iz -> ize (Example :}  siz(ed) -> size)

\textbf{s -> є (Example :}  cats -> cat)

(*d &! (*l or *s or *z)) -> single letter

(Example : hopp(ing) -> hop ; tann(ed) -> tan ; fall(ing) -> fall ; hiss(ing) -> hiss ; fizz(ed) -> fizz)

(m=1 and *o) -> e

\textbf{(Example :}  fil(ing) -> file); fail(ing) -> fail

\textbf{Slide 115:}

\textbf{Step 1c :}  Y Elimination

\textbf{( \*v\*) y -> i (Example :}  happy -> happi ; sky -> sky)

Step 1 deals with plurals and past participles. The subsequent steps are much more straightforward.

\textbf{Step 2 :} Derivational Morphology -1

\textbf{(m>0) ational -> ate (Example :}  relational -> relate

(m>0) ization -> ize (generalization -> generalize)

(m>0) biliti -> ble (sensibiliti -> sensible)

\textbf{Slide 116:}

\textbf{Step 3 :} 

\textbf{(m>0) icate -> ic (Example :}  triplicate -> triplic)

\textbf{(m>0) ful -> є (Example :}  hopeful -> hope)

\textbf{(m>0) ness -> є (Example :}  goodness -> good)

\textbf{Slide 117:}

\textbf{Step 4 :} Derivational Morphology-II

\textbf{( (m>1) ance -> є (Example :}  allowance -> allow)

\textbf{(m>1) ment -> є(Example :}  adjustment -> adjust)

\textbf{(m>1) ent -> є (Example :}  dependent -> depend)

\textbf{(m>1) ive -> є (Example :}  effective -> effect)

\textbf{Slide 118:}

The suffixes are now removed. All that remains is a little tidying up.

\textbf{Step 5a :} 

\textbf{(m>1) e -> є (Example :}  probate -> probat ; rate -> rate)

\textbf{(m=1 and not *o) ness -> є (Example :}  goodness -> good)

Step 5b

(m > 1 and *d &*l) -> single letter

\textbf{(Example :}  controll -> control ; roll -> roll)

\textbf{Slide 119:}

Disadvantage

This indiscriminate cutting will be successful at some occasion and fail in some

Eg.

Studies- 	 Studi

Giving-   	 Giv

Intelligence-Intelligen

Produced immediate representation of the word  may not have any meaning

\textbf{Slide 120:}

Stemming					Lemmatization

Cutting of Suffixes to extract stem

Eg.

Studies- 	 Studi

Giving-   	 Giv

Caring-         Car

Produced immediate representation of the word  may not have any meaning

Cutting of Suffixes to extract Lemma

Eg.

Studies- 	 Study

Giving-   	 Give

Caring-         Care

Produced immediate representation of the word  having meaning

\textbf{Slide 121:}

Lemmatization

Same as Stemming but immediate representation have some meaning

Stemming and lemmatization both of these concepts are used to normalized the given word by removing affixes

\textbf{Slide 122:}

N-Gram Language Model

Predicting Nth    word from  N-1 words

Predicting 3rd word from previous 2 words – Model is called Trigram

Predicting 2nd word from previous 1 words – Model is called Bigram

He is going to \_\_\_\_\_      Predicting 5th word from last four words is 5 gram Language Model

going to \_\_\_\_\_     Predicting 3rd  word from last two words is trigram

\textbf{Slide 123:}

Application of N-Gram

Optical Character Recognition

Image   Text

If words are missing or not clear can be predicted

Grammar correction

Spelling is not correct then based on context can be suggested

If spelling is right

E.g. Deer sir instead of Dear sir can be correctly as, Sir is normally followed by Dear not Deer

Speech to Text

Words with different pronunciation. e.g. Eye am Fine & I am Fine

Translation

Multiple Synonym of words

E.g., He is biggest minister of Pakistan  Prime

Suggestion while typing in mobile text prediction mode.

\textbf{Slide 124:}

He is going to school Meaningful

He school is to going Doesn’t Make sense

For N-gram models

P(w1,w2, ...,wn)

\textbf{By the Chain Rule we can decompose a joint probability, as follows:} 

P(w1,w2, ...,wn) = P(w1) P(w2|w1) P(w3|w1w2 ) P(wn|w1..., Wn-2 wn-1)

Join Probability of sentence

P(He is going to school)= P(He)P(is|He)P(going|He is)P(to|He is going)P(school|He is going to)

Looking back so much


\subsection{2_Porter Stemmer Algorithm (3)}

\textbf{Slide 1:}

Porter Stemmer Algorithm

\textbf{Slide 2:}

Stemming

Stemming is suffix stripping operation

Process of reducing word into its base form (Root form/stem form)

This is achieved by cutting of begging or end of the word

\textbf{Eg:} 

Playing-    Play

Boys-         Boy

Popular Algorithm is Porter Stemmer

\textbf{Slide 3:}

Porter Stemmer

– Five sets of rules, applied in order

– Within each set, if more than one of the rules can apply, only the one with the longest matching suffix (S1) is followed

\textbf{Advantage:}   easy to see understand, easy to implement.

\textbf{Slide 4:}

Algorithm

Few definitions.

Consonant-

it  is a letter other than A, E, I, O or U, and other than Y preceded by a consonant.

So, in TOY the consonants are T and Y, and in SYZYGY they are S, Z and G.

If a letter is not a consonant it is a vowel.

\textbf{Slide 5:}

Algorithm

A consonant will be denoted by c, a vowel by v.

A list ccc... of length greater than 0 will be denoted by C,

and a list vvv... of length greater than 0 will be denoted by V.

\textbf{Any word, or part of a word, therefore, has one of the four forms:} 

CVCV ... C

CVCV ... V

VCVC ... C

VCVC ... V

\textbf{Slide 6:}

Algorithm

These may all be represented by the single form[C]VCVC ... [V]where the square brackets denote arbitrary presence of their contents.

Using (\$VC^m\$) to denote VC repeated m times, this may again be written as[C](\$ VC^m \$)[V] m will be called the measure of any word or word part when represented in this form.

The case m = 0 covers the null word. Here are some examples:(m is no of VC pairs)

m=0 TR, EE, TREE, Y, BY.

m=1 TROUBLE, OATS, TREES, IVY.

m=2 TROUBLES, PRIVATE, OATEN, ORRERY.

\textbf{Slide 8:}

The rules for removing a suffix will be given in the form(condition) S1 -> S2This means that if a word ends with the suffix S1, and the stem before S1 satisfies the given condition, S1 is replaced by S2. The condition is usually given in terms of m,

e.g.(m > 1) EMENT ->Here S1 is 'EMENT' and S2 is null.

This would map REPLACEMENT to REPLAC, since REPLAC is a word part for which m = 2.

\textbf{Slide 9:}

\textbf{The 'condition' part may also contain the following:} 

*S - the stem ends with S (and similarly for the other letters).

*v* - the stem contains a vowel.

m=2 TROUBLES, PRIVATE, OATEN, ORRERY.

*d - the stem ends with a double consonant (e.g. -TT, -SS).

*o - the stem ends cvc, where the second c is not W, X or Y (e.g. -WIL, -HOP).

\textbf{Slide 10:}

And the condition part may also contain expressions with and, or and not, so that:

\textbf{(m>1 and (*S or *T)) :}  tests for a stem with m>1 ending in S or T, while

(*d and not (*L or *S or *Z)) : tests for a stem ending with a double consonant other than L, S or Z.

Elaborate conditions like this are required only rarely.

\textbf{Slide 11:}

In a set of rules written beneath each other, only one is obeyed, and this will be the one with the longest matching S1 for the given word.

For example, with

SSES -> SS

IES -> I

SS -> SS

S ->

(here the conditions are all null) CARESSES maps to CARESS since SSES is the longest match for S1. Equally CARESS maps to CARESS (S1=SS) and CARES to CARE (S1=S).

In the rules below, examples of their application, successful or otherwise, are given on the right in lower case. The algorithm now follows:

\textbf{Slide 12:}

\textbf{Step 1a :} 

\textbf{SSES -> SS (Example :}  caresses -> caress)

\textbf{IES -> I (Example :}  ponies -> poni ; ties -> ti)

\textbf{SS -> SS (Example :}  caress -> caress)

\textbf{S -> (Example :}  cats -> cat)

\textbf{Slide 13:}

\textbf{Step 1b :} 

\textbf{(m>0) EED -> EE (Example :}  feed -> feed ; agreed -> agree)

\textbf{(v) ED -> (Example :}  plastered -> plaster ; bled -> bled)

\textbf{(v) ING -> (Example :}  motoring -> motor ; sing -> sing)

\textbf{S -> (Example :}  cats -> cat)

\textbf{Slide 14:}

If the second or third of the rules in Step 1b is successful, the following is done:

\textbf{AT -> ATE (Example :}  conflat(ed) -> conflate)

\textbf{BL -> BLE (Example :}  troubl(ed) -> trouble)

\textbf{IZ -> IZE (Example :}  siz(ed) -> size)

\textbf{S -> (Example :}  cats -> cat)

(*d and not (*L a q a or *S or *Z)) -> single letter

(Example : hopp(ing) -> hop ; tann(ed) -> tan ; fall(ing) -> fall ; hiss(ing) -> hiss ; fizz(ed) -> fizz)

\textbf{(m=1 and *o) -> E (Example :}  fail(ing) -> fail ; fil(ing) -> file)

\textbf{Slide 15:}

The rule to map to a single letter causes the removal of one of the double letter pair.

The -E is put back on -AT, -BL and -IZ, so that the suffixes -ATE, -BLE and -IZE can be recognized later.

This E may be removed in step 4.

\textbf{Slide 16:}

\textbf{Step 1c :} 

\textbf{(\*v\*) Y -> I (Example :}  happy -> happi ; sky -> sky)

Step 1 deals with plurals and past participles. The subsequent steps are much more straightforward.

\textbf{Slide 17:}

\textbf{Step 2 :} 

\textbf{(m>0) ATIONAL -> ATE (Example :}  relational -> relate

(m>0) TIONAL -> TION (Example : conditional -> condition ; rational -> rational)

\textbf{(m>0) ENCI -> ENCE (Example :}  valenci -> valence)

\textbf{(m>0) ANCI -> ANCE (Example :}  hesitanci -> hesitance)

\textbf{(m>0) IZER -> IZE (Example :}  digitizer -> digitize)

\textbf{(m>0) ABLI -> ABLE (Example :}  conformabli -> conformable)

\textbf{(m>0) ALLI -> AL (Example :}  radicalli -> radical)

(m>0) ENTLI -> ENT (differentli -> different)

(m>0) ELI -> E (vileli -> vile)

(m>0) OUSLI -> OUS (analogousli -> analogous)

(m>0) IZATION -> IZE (vietnamization -> vietnamize)

\textbf{Slide 18:}

\textbf{Step 2:} Continue

(m>0) ATION -> ATE (predication -> predicate)

(m>0) ATOR -> ATE (operator -> operate)

(m>0) ALISM -> AL (feudalism -> feudal)

(m>0) IVENESS -> IVE (decisiveness -> decisive)

(m>0) FULNESS -> FUL (hopefulness -> hopeful)

(m>0) OUSNESS -> OUS (callousness -> callous)

(m>0) ALITI -> AL (formaliti -> formal)

(m>0) IVITI -> IVE (sensitiviti -> sensitive)

(m>0) BILITI -> BLE (sensibiliti -> sensible)

\textbf{Slide 19:}

\textbf{Step 2:} Continue

The test for the string S1 can be made fast by doing a program switch on the penultimate letter of the word being tested. This gives a fairly even breakdown of the possible values of the string S1. It will be seen in fact that the S1-strings in step 2 are presented here in the alphabetical order of their penultimate letter. Similar techniques may be applied in the other steps.

\textbf{Slide 20:}

\textbf{Step 3 :} 

\textbf{(m>0) ICATE -> IC (Example :}  triplicate -> triplic)

\textbf{(m>0) ATIVE -> (Example :}  formative -> form)

\textbf{(m>0) ALIZE -> AL (Example :}  formalize -> formal)

\textbf{(m>0) ICITI -> IC (Example :}  electriciti -> electric))

\textbf{(m>0) ICAL -> IC (Example :}  electrical -> electric)

\textbf{(m>0) FUL -> (Example :}  hopeful -> hope)

\textbf{(m>0) NESS -> (Example :}  goodness -> good)

\textbf{Slide 21:}

\textbf{Step 4 :} 

\textbf{(m>1) AL -> (Example :}  revival -> reviv)

\textbf{(m>1) ANCE -> (Example :}  allowance -> allow)

\textbf{(m>1) ENCE -> (Example :}  inference -> infer)

\textbf{(m>1) ER -> (Example :}  airliner -> airlin)

\textbf{(m>1) IC -> (Example :}  gyroscopic -> gyroscop)

\textbf{(m>1) ABLE -> (Example :}  adjustable -> adjust)

\textbf{(m>1) IBLE -> (Example :}  defensible -> defens)

\textbf{(m>1) ANT -> (Example :}  irritant -> irrit)

\textbf{(m>1) EMENT -> (Example :}  replacement -> replac)

\textbf{(m>1) MENT -> (Example :}  adjustment -> adjust)

\textbf{(m>1) ENT -> (Example :}  dependent -> depend)

\textbf{Slide 22:}

\textbf{Step 4 :} Continued

\textbf{(m>1 and (*S or *T)) ION -> (Example :}  adoption -> adopt)​

\textbf{(m>1) OU -> (Example :}  homologou -> homolog)​

\textbf{(m>1) ISM -> (Example :}  communism -> commun)​

\textbf{(m>1) ATE -> (Example :}  activate -> activ)​

\textbf{(m>1) ITI -> (Example :}  angulariti -> angular)​

\textbf{(m>1) OUS -> (Example :}  homologous -> homolog)​

\textbf{(m>1) IVE -> (Example :}  effective -> effect)​

\textbf{(m>1) IZE -> (Example :}  bowdlerize -> bowdler)​

The suffixes are now removed. All that remains is a little tidying up.​​

\textbf{Slide 23:}

\textbf{Step 5a :} 

\textbf{(m>1) E -> (Example :}  probate -> probat ; rate -> rate)

\textbf{(m=1 and not *o) E -> (Example :}  cease -> ceas)

\textbf{Slide 24:}

Step 5b

(m > 1 and *d and *L) -> single letter

\textbf{(Example :}  controll -> control ; roll -> roll)


\subsection{3_Syntax Analysis_2023 (2)}

\textbf{Slide 1:}

Syntax Analysis

\textbf{Slide 2:}

Steps in NLP

John Ate the Apple

Ate  the Apple John

\textbf{Slide 3:}

Tagging

Tagging is a kind of classification that may be defined as the automatic assignment of description to the tokens.

Here the descriptor is called tag, which may represent one of the part-of-speech, semantic information and so on.

parts of speech are defined instead based on their grammatical relationship with neighboring words or the morphological properties about their affixes

\textbf{Slide 4:}

Part-of-Speech (PoS) tagging

the process of assigning one of the parts of speech to the given word

POS tagging is a task of labelling each word in a sentence with its appropriate part of speech

Examples of Parts Of Speech – nouns, verb, adverbs, adjectives, pronouns, conjunction and their sub-categories

\textbf{Slide 5:}

POS Tagging Example

Annotate each word in a sentence with a part-of-speech.

A POS tagger considers surrounding words while assigning a tag.

[(“I”, “Pronoun”), (“like”, “Verb”), (“to”, “To”), (“read”, “Verb”), (“books”, “Noun”)]

\textbf{Slide 6:}

Categories of Parts of Speech

Closed class

Open class

\textbf{Slide 7:}

Open Class (Mostly Content word)

Called open class because we will keep on adding new words in the language example google,skype,photoshop

Noun,verbs,adjective,adverbs

Mostly content bearing –they refer to objects,actions,and features in the world

\textbf{Slide 8:}

Close class (function words)

Pronoun, determiner, prepositions, connectives

Limited in number

Used to tie the concepts of the sentences together we user these closed words

\textbf{Slide 9:}

\textbf{POS tags:}  level of Details

How much detail we want to specify the particular category

Level of details of POS Tag

Coarse grained                          Fine grained

Decision in dependent on the application

\textbf{Slide 10:}

Coarse Grained(Small Tag Set)

Here we do not provide the pos tag in much details(Grammatical Details)

For example –we can identify a word as NOUN but we can not tell whether it Singular Noun or Plural Noun

\textbf{Slide 11:}

Fines grained (Large Tag set)

In fine grained level we provide grammatical details of a word

For example if Verb then is it past tense ,future or present  tense

\textbf{NOTE :} May Leads to confusion due to too many POS tags

\textbf{Slide 12:}

POS TAG / Word Classes

9  traditional word classes of parts of speech

Noun, verb, adjective, preposition, adverb, article, pronoun, conjunction, interjection

N		noun		chair, bandwidth, pacing

V		verb		study, debate, munch

ADJ		adjective	purple, tall, ridiculous

ADV	adverb		unfortunately, slowly

P		preposition	of, by, to

PRO		pronoun		I, me, mine

DET		determiner	the, a, that, those

12

\textbf{Slide 13:}

Part-of-speech tagging

process of assigning a part-of-speech to each word in part-of-speech tagging a text.

input sequence x1, x2,..., xn of (tokenized) words and a tagset,

Output  is a sequence y1, y2,..., yn of tags,

each output yi corresponding exactly to one input xi as shown   below

\textbf{Slide 14:}

POS Tagging is disambiguation task

Words are ambiguous(more than one possible part-of-speech)

the goal is to find the correct tag for the situation

For example,

book can be a verb (book that flight)

Or

a noun (hand me that book)

The goal of POS-tagging is to resolve these ambiguity resolution ambiguities, choosing the proper tag for the context

\textbf{Slide 15:}

Part-of-Speech Tag sets

There are various tag sets to choose.

•The choice of the tag set depends on the nature of the application.

–We may use small tag set (more general tags) or

–large tag set (finer tags).

\textbf{•Some of widely used part-of-speech tag sets:} 

–Penn Treebank has 45 tags

–Brown Corpus has 87 tags

–C7 tag set has 146 tags

•In a tagged corpus, each word is associated with a tag from the used tag set.

\textbf{Slide 16:}

Penn Treebank Tagset

16

\textbf{Slide 17:}

POS Tagging Approaches

Rule-based tagging

E.g. EnCG ENGTWOL tagger

Stochastic, or, Probabilistic tagging

HMM (Hidden Markov Model) tagging

Transformation-based tagging

Learned rules (statistic and linguistic)

E.g., Brill tagger

17

\textbf{Slide 18:}

POS Tagging Approaches

\textbf{Slide 19:}

Rule Based POS Tagging Approaches

\textbf{Input:} - Sentence

Tag dictionary /Lexicon-Using Tag dictionary POS tags will be assigned to the words in the given sentence

\textbf{For Example :} -

I play cricket everyday (Play is Verb)

I want to perform a play  (Play is Noun )

(Here play has two tags Noun and Verb Ambiguous POS )

Handwritten rules-

For example an Ambiguous word is noun if it follows  a determiner

\textbf{Slide 20:}

Rule-Based Tagging

The rule-based approach uses handcrafted sets of rules to tag input sentence.

\textbf{Stage 1:} 

Typically…start with a dictionary of words and possible tags

Assign all possible tags to words using the dictionary

\textbf{Stage 2:} 

Write rules by hand to selectively remove tags

Stop when each word has exactly one (correct) tag

20

\textbf{Slide 21:}

Rules based on context

Please book that flight

I bought a Book.

Book

Wn-1

Wn

Wn+1

Book

Tn-1

Tn

Tn+1

\textbf{•Rule-2:}  if (the previous tag is an article)

then eliminate all verb tags

\textbf{Rule-1:}  if (the previous word is “to”)

then eliminate all Noun tags

\textbf{Slide 22:}

Apply Rules Eliminating Some POS

E.g., Eliminate VBN if VBD is an option when VBN|VBD follows “<start> PRP”

NN

RB

JJ			VB

PRP	VBD		TO	VB		DT	NN

She	promised	to	back 		the	bill

22

\textbf{Slide 23:}

\textbf{Rule-Based Part-of-Speech Tagging:}  Example

Pronoun

Verb

Article

Verb     Noun

Pronoun

Verb

Article

Verb     Noun

\textbf{Slide 24:}

Rule Base Tagger

It is oldest approach and uses hand written rules for tagging

Depends on dictionary or lexicon to get possible tags for each word to tagged

Language experts will sit together to find out what are the symbol patterns.

Hand written rules are used to identify the correct tag when the word has more than one possible tag.

Disambiguation is done by analysing the linguistic features of the word, its preceding word following word and other aspects .

\textbf{Slide 25:}

Rule Base Tagger

For example, if the preceding word is article,  then the word must be a noun. ​

​We can also understand Rule-based POS tagging by its two-stage architecture −

The first stage − In the first stage, create a dictionary to assign each word a list of potential parts-of-speech.

The second stage − In the second stage, create a large list of hand-written disambiguation rules to sort down the list to a single part-of-speech for each word as per requirements.

TAGGIT the first large rule-based tagger used context pattern rules.​

​

TAGGIT used the set of 71 tags and 3300 disambiguation rules.​

These rules disambiguated 77% of words in the million –word brown.​

\textbf{Slide 26:}

Example

University Corpus

I want to read a book.

Book-Noun or Verb

\textbf{Rule :} -Before book ,determiner (a) is there ,therefore it is a noun.

\textbf{Slide 27:}

Limitation of Rule based Approach

Hard coded rules are required

Rules has to be updated based on language change

It needs to check lexicon every time which is not efficient approach

Strong language expert team is required.

\textbf{Slide 28:}

Hidden Markov Model

\textbf{An HMM is a probabilistic sequence model:} 

given a sequence of units (words, letters, morphemes, sentences, whatever), it computes a probability distribution over possible sequences of labels and chooses the best label sequence.

\textbf{Slide 29:}

Markav Chain

Markov chains, named after Andrey Markov, a stochastic model that depicts a sequence of possible events

where predictions or probabilities for the next state are based solely on its previous event state, not the states before.

In simple words, the probability that n+1th steps will be x depends only on the nth steps not the complete sequence of steps that came before n.

This property is known as Markov Property or Memorylessness.

Let us explore our Markov chain with the help of a diagram,

\textbf{Slide 30:}

A diagram representing a two-state(here, E and A) Markov process.

Here the arrows originated from the current state and point to the future state and the number associated with the arrows indicates the probability of the Markov process changing from one state to another state.

For instance, if the Markov process is in state E, then the probability it changes to state A is 0.7, while the probability it remains in the same state is 0.3. Similarly, for any process in state A, the probability to change to E state is 0.4 and the probability to remain in the same state is 0.6.

\textbf{Slide 31:}

Markov Chain

tells us something about the probabilities of sequences of random variables, states, each of which can take on values from some set

These sets can be words or tags or symbol representing anything for example weather

Markov Assumption

if we want to predict the future in the sequence, all that matters is the current state.

All the states before the current state have no impact on the future except via the current state.

\textbf{Slide 33:}

\textbf{Markov chain is specified by the following components:} 

\textbf{Slide 34:}

The Hidden Markov Model

A Markov chain is useful when we need to compute a probability for a sequence of observable events

In many cases, however, the events we are interested in are hidden : we don’t observe them directly

For example we don’t normally observe part-of-speech tags in a text. Rather, we see words, and must infer the tags from the word sequence. We call the tags hidden because they are not observed

\textbf{Slide 35:}

A hidden Markov model

A hidden Markov model (HMM) allows us to talk about both

observed events (like words that we see in the input) and

hidden events (like part-of-speech tags) that we think of as causal factors in our probabilistic model.

\textbf{Slide 36:}

HMM is specified by the following components

\textbf{Slide 37:}

A first-order hidden Markov model instantiates two simplifying assumptions.

First, as with a first-order Markov chain, the probability of a particular state depends only on the previous state:

Second, the probability of an output observation oi depends only on the state that produced the observation qi and not on any other states or any other observations:

\textbf{Slide 43:}

Transition probability

The transition probability is the likelihood of a particular sequence for example, how likely is that a noun is followed by a model and a model by a verb and a verb by a noun.

This probability is known as Transition probability. It should be high for a particular sequence to be correct.

\textbf{Emission probability :} 

These sets of probabilities are Emission probabilities and should be high for our tagging to be likely.

Now, what is the probability that the word Ted is a noun, will is a model, spot is a verb and Will is a noun.

\textbf{Slide 50:}

The components of an HMM tagger

An HMM has two components, the A and B probabilities

The A matrix contains the tag transition probabilities P(ti |ti−1) which represent the probability of a tag occurring given the previous tag.

For example, modal verbs like will are very likely to be followed by a verb in the base form, a VB, like race, so we expect this probability to be high.

We compute the maximum likelihood estimate of this transition probability by counting, out of the times we see the first tag in a labeled corpus, how often the first tag is followed by the second:

\textbf{Slide 51:}

The B emission probabilities, P(wi |ti), represent the probability, given a tag (say MD), that it will be associated with a given word (say will).

The MLE of the emission probability is

\textbf{Slide 52:}

HMM tagging as decoding

HMM contains hidden variables, the task of determining the hidden variables sequence corresponding to the sequence of observations decoding is called decoding.

\textbf{Slide 53:}

For part-of-speech tagging, the goal of HMM decoding is to choose the tag sequence t1 ...tn that is most probable given the observation sequence of n words w1 ...wn:

\textbf{The way we’ll do this in the HMM is to use Bayes’ rule to instead compute:} 

\textbf{Furthermore, we simplify Eq. Above  by dropping the denominator P(w n 1 ):} 

\textbf{Slide 54:}

HMM taggers make two further simplifying assumptions.

The first is that the probability of a word appearing depends only on its own tag and is independent of neighboring words and tags:

The second assumption, the bigram assumption, is that the probability of a tag is dependent only on the previous tag, rather than the entire tag sequence;

\textbf{Slide 55:}

Plugging the simplifying assumptions from Eq. X  and Eq. Y  into Eq. Z results in the following equation for the most probable tag sequence from a bigram tagger:

\textbf{Slide 56:}

\textbf{Find the probability of the following sentence :} Jane will spot Will

\textbf{Slide 59:}

In general, we  from all the possible combination of parts of speech tags we will pick the one that generates the above sentence Jane will spot Will with the highest probability

How  to find the correct path ?

\textbf{Slide 62:}

Which one of this generates the sentence with highest probability ?

\textbf{Slide 68:}

For every possible state we require the highest probability path that ends in that state .

This is known as Viterbi HMM Algorithm

\textbf{Slide 70:}

Emission Probability Matrix

Transition  Probability Matrix

\textbf{Slide 81:}

We keep the edge which gives us a high value and remembering the edge which gives us a highest value .

We have all the probabilities and edges which gives those probabilities.

To get the optimal path, we start from end and trace backwards

Each state has only one incoming edge ,this will give us optimal path with highest probabilities .

\textbf{Slide 83:}

Consider the following data set

Apply HMM to find Tag Sequence for Will can spot Mary

\textbf{Slide 84:}

Transition Probability

\textbf{Slide 85:}

Emission  Probability


\subsection{4_Constituency Grammars(1) (2)}

\textbf{Slide 1:}

Constituency Grammars

\textbf{Slide 2:}

Context-free grammars

CFGS the backbone of many formal models of the syntax of natural language (and, for that matter, of computer languages).

they play a role in many computational applications, including grammar checking, semantic interpretation, dialogue understanding, and machine translation.

express sophisticated relations among the words in a sentence, yet computationally tractable enough that efficient algorithms exist for parsing sentences with them

\textbf{Slide 3:}

Constituency

Syntactic constituency is the idea that groups of words can behave as single units, or constituents.

For example, group of words come together to form a Noun Phrase.

Part of developing a grammar involves building an inventory of the constituents in the language.

\textbf{Slide 4:}

How do words group together in English?

Consider noun phrase the noun phrase, a sequence of words surrounding at least one noun.

Here are some examples of noun phrases

\textbf{Slide 5:}

How do words group together in English?

How these words for constituents or they group together ?

One evidence is that they can all appear in similar syntactic environments, for example, before a verb.

\textbf{Slide 6:}

But while the whole noun phrase can occur before a verb, this is not true for each of the individual words that make up a noun phrase.

The following are not grammatical sentences of English (we use an asterisk (*) to mark fragments that are not grammatical English sentences):

\textbf{Slide 7:}

Pre-posed or post-posed constructions

For example, the prepositional phrase on September seventeenth can be placed in a number of different locations in the following examples, including at the beginning (preposed) or at the end (postposed):

But again, while the entire phrase can be placed differently, the individual words making up the phrase cannot be

\textbf{Slide 8:}

Verb Phase

\textbf{Slide 9:}

Adjective Phrase

\textbf{Slide 10:}

Context-Free Grammars

formal system for modeling constituent structure in English and other natural languages is the Context-Free Grammar, or CFG.

It is also called Phrase-Structure Grammars

\textbf{Slide 11:}

Context-Free Grammars

A context-free grammar consists of a set of rules or productions

Which expresses the ways that symbols of the language can be grouped and ordered together, and a lexicon of words and symbols

Following is the Example of production rules of - CFG

\textbf{Slide 12:}

Context-free Grammars

Context-free rules can be hierarchically embedded, so we can combine the previous rules with others, like the following, that express facts about the lexicon:

\textbf{Slide 13:}

Symbols in CFG

symbols that are used in a CFG are divided into two classes

Terminal                                    non-terminal

The symbols

terminal that correspond to words in the language (“the”, “nightclub”)

symbols that express abstractions over these terminals are called non-terminals.

the item to the right of the arrow (→) is an ordered list of one or more terminals and non-terminals;

to the left of the arrow is a single non-terminal symbol expressing some cluster or generalization.

\textbf{Slide 14:}

Two ways of CFG

a device for generating sentences and

a device for assigning a structure to a given sentence.

\textbf{Slide 15:}

CFG as a generator

we can read the → arrow as “rewrite the symbol on the left with the string of symbols on the right”.

a flight can be derived from the non-terminal NP

\textbf{Slide 16:}

CFG as a generator

Thus, a CFG can be used to generate a set of strings. ]

This sequence of rule expansions is called a derivation of the string of words.

It is common to represent a derivation by a parse tree (commonly shown inverted with the root at the top).

Figure above “A parse tree for “a flight”” shows the tree representation of this derivation

\textbf{Slide 17:}

Start Symbol

The formal language defined by a CFG is the set of strings that are derivable start symbol from the designated start symbol.

Each grammar must have one designated start symbol, which is often called S.

Since context-free grammars are often used to define sentences, S is usually interpreted as the “sentence” node, and the set of strings that are derivable from S is the set of sentences in some simplified version of English.

\textbf{Slide 18:}

Few additional rules

The following rule expresses the fact that a sentence can consist of a noun phrase followed by a verb phrase:verb phrase

S → NP VP        I prefer a morning flight

A verb phrase in English consists of a verb followed by assorted other things;for example, one kind of verb phrase consists of a verb followed by a noun phrase:

VP → Verb NP prefer a morning flight

\textbf{Slide 19:}

Few additional rules

\textbf{the verb may be followed by a noun phrase and a prepositional phrase:} 

VP → Verb NP PP     leave Boston in the morning

\textbf{the verb phrase may have a verb followed by a prepositional phrase alone:} 

VP → Verb PP          leaving on Thursday

A prepositional phrase generally has a preposition followed by a noun phrase.

For example, a common type of prepositional phrase in the ATIS corpus is used to indicate location or direction:

PP → Preposition NP     from Los Angeles

\textbf{Slide 20:}

Example of lexicons

\textbf{Slide 21:}

The grammar for L0, with example phrases for each rule

\textbf{Slide 22:}

The parse tree for “I prefer a morning flight” according to grammar L0

\textbf{Slide 23:}

Formal Definition of Context-Free Grammar

A context-free grammar G is defined by four parameters: N, Σ, R, S (technically this is a “4-tuple”).

The problem of mapping from a string of words to its parse tree is called syntactic parsing;

\textbf{Slide 24:}

Constituency Parsing

Syntactic parsing is the task of assigning a syntactic structure to a sentence

Parse trees can be used in applications such as grammar checking sentence that cannot be parsed may have grammatical errors (or at least be hard to read).

Parse trees can be an intermediate stage of representation for semantic analysis

\textbf{Slide 25:}

Ambiguity

\textbf{Slide 28:}

Explanantion

Two parse trees for an ambiguous sentence.

The parse on the left corresponds to the humorous reading in which the elephant is in the pajamas,

the parse on the right corresponds to the reading in which Captain Spaulding did the shooting in his pajamas

\textbf{Slide 29:}

Types of ambiguity

Two common  kinds of ambiguity are attachment ambiguity and coordination ambiguity.

A  sentence has an attachment ambiguity if a particular constituent can be attached to the parse tree at more than one place.

The Groucho Marx sentence is an example of PP-attachment ambiguity. Various kinds of adverbial phrases are also subject to this kind of ambiguity.

For instance, in the following example the gerundive-VP flying to Paris can be part of a gerundive sentence whose subject is the Eiffel Tower or it can be an adjunct modifying the VP headed by saw:

We saw the Eiffel Tower flying to Paris.

attachment ambiguity

In coordination ambiguity phrases can be conjoined by a conjunction like and coordination ambiguity

For example, the phrase old men and women can be bracketed as [old [men and women]], referring to old men and old women, or as [old men] and [women], in which case it is only the men who are old.

coordination ambiguity

\textbf{Slide 30:}

\textbf{CKY Parsing:}  A Dynamic Programming Approach

\textbf{Slide 31:}

How to Convert CFG to CNF

\textbf{Slide 34:}

CKYAlgorithm

\textbf{Slide 35:}

CKYAlgorithm

\textbf{Slide 37:}

Cocke-Younger-Kasami Algorithm

It is used to solves the membership problem using a dynamic programming approach.

The algorithm is based on the principle that the solution to problem [i, j] can constructed from solution to subproblem [i, k] and solution to sub problem [k, j].

The algorithm requires the Grammar G to be in Chomsky Normal Form (CNF). Note that any Context-Free Grammar can be systematically converted to CNF.

This restriction is employed so that each problem can only be divided into two subproblems and not more – to bound the time complexity.

\textbf{Slide 38:}

How does the CYK Algorithm work?

For a string of length N, construct a table T of size N x N. Each cell in the table T[i, j] is the set of all constituents that can produce the substring spanning from position i to j.

The process involves filling the table with the solutions to the subproblems encountered in the bottom-up parsing process.

Therefore, cells will be filled from left to right and bottom to top.

\textbf{Slide 39:}

In T[i, j], the row number i denotes the start index, and the column number j denotes the end index.

\textbf{Slide 40:}

The algorithm considers every possible subsequence of letters and adds K to T[i, j] if the sequence of letters starting from i to j can be generated from the non-terminal K.

For subsequences of length 2 and greater, it considers every possible partition of the subsequence into two parts, and checks if there is a rule of the form A -> BC in the grammar where B and C can generate the two parts respectively, based on already existing entries in T.

The sentence can be produced by the grammar only if the entire string is matched by the start symbol, i.e, if S is a member of T[1, n].

\textbf{Slide 45:}

CYK Algorithm


\subsection{5_Probablilstic Context Free Grammar}

\textbf{Slide 1:}

Probablilstic Context Free Grammar

\textbf{Slide 2:}

USE OF PCFG

\textbf{Slide 3:}

Probabilistic CFG

\textbf{Slide 6:}

Parse Tree

\textbf{Slide 7:}

Probabilities of The above Tree


\subsection{6_Dependency Parsing (2)}

\textbf{Slide 1:}

Dependency Parsing

\textbf{Slide 2:}

Introduction

In some languages, the same statement can be written in such a way that on translating in English (word to word translation), it might yield ‘He is standing wall behind’ i.e the word order is free.

It isn’t necessary the adjective has to appear just before the noun it like ‘Handsome boy’.

It can be ‘Boy handsome’.

Hence, using rules for parsing such languages is problematic!!

And that is why dependency parsing is required.

\textbf{Slide 3:}

Dependency Parsing

Dependency parsing helps us build a parsing tree with the tags used determining the relationship between words

in the sentence rather than using any Grammar rule as used for syntactic parsing which gives a lot of flexibility even when the order of words (like ‘boy handsome’ or ‘handsome boy’) get changed.

\textbf{A sample dependency tree can be observed below:} 

\textbf{Slide 4:}

\textbf{Below is an example of Dependency parsing for:} 

‘I prefer the morning flight through Denver’

\textbf{Slide 5:}

terminologies as well used for Dependency parsing

\textbf{Head-Dependent:} 

In the arrows representing relationship, the origin word is the Head & the destination word is Dependent.

\textbf{For example:}  (I, ‘nsubj’, prefer), ‘prefer’ is Head \& ‘I’ is Dependent.

Root: Word which is the root of our parse tree. It is ‘prefer’ in the above example.

\textbf{Slide 6:}

terminologies as well used for Dependency parsing

Grammar Functions and Arcs: Tags between each Head-Dependent pair is a grammar function determining the relation between the Head & Dependent.

The arrowhead carrying the tag is called an Arc. Below is the table that interprets the meaning of some common grammar functions:

\textbf{Slide 7:}

\textbf{Some examples for Grammar Functions are given below:} 

\textbf{Slide 8:}

\textbf{Some examples for Grammar Functions are given below:} 

\textbf{Slide 9:}

Understanding the Dependency parse tree

we represent dependencies as a directed graph G= (V, A)

where V(set of vertices) represents words (and punctuation marks as well) in the sentence &

A( set of arcs) represent the grammar relationship between elements of V.

\textbf{Slide 10:}

Understanding the Dependency parse tree

A dependency parse tree is the directed graph mentioned above which has the below features:

Root has no Incoming arcs (can only be Head in Head-Dependent pair)

Vertices(except Root) should have only one incoming arc (Only one Parent/Head)

A Unique path should exist between Root & each vertex in the tree.

\textbf{Slide 11:}

Projectivity

Projective arc: An arc/arrows\_with\_tag are projective when ‘Head’ associated with the arc has a path to reach each word that lies between ‘Head’ \& ‘Dependent’.

\textbf{Slide 12:}

Explanation

Here, the arc between ‘the’ & ‘flights’ is projective as the ‘Head’ i.e ‘flights’ has a path to ‘morning’ as well which lies between ‘the’ & ‘flights. Similarly, arc between ‘canceled’ & ‘flights’ is also projective as ‘canceled’ has a path to ‘the’(canceled →flights →the) & ‘morning’ (canceled →flights →morning) which lies between Head ‘canceled’ & Dependent ‘flights’

\textbf{Slide 13:}

Projective and Non-Projective Parse Tress

Projective parse tree: A parse tree with all its arcs projective. The above tree is projective.

Non-projective parse tree: A tree with at least one of the arcs as non-projective. Have a look at the below example

\textbf{Slide 14:}

Non-Projective Parse Tree Example

Here, the arc between ‘flight’ & ‘was’ is non-projective as there exists no path between ‘flight’ & ‘morning’

\textbf{Slide 15:}

\textbf{how to create a dependency parse tree-Shift Reduce Parsing (Arc standard):} 

\textbf{We need some setup for this algorithm:} 

A Stack

Buffer of input words

Oracle function/parser & set of dependency relations (like obj, iobj, etc tags that are used to represent arcs).

\textbf{Slide 17:}

Initially, the Stack has [Root], a dummy variable.

All the words of the sentence to be parsed are in the buffer.

A word is popped from buffer & pushed in Stack.

As a word is pushed in Stack, the Oracle/parser takes the top two elements of the stack & outputs/predict any of the below 3 transition with relationship type:

\textbf{LEFTARC\_FUNCTION:} 

RIGHTARC\_FUNCTION

\textbf{SHIFT:} 

\textbf{Slide 18:}

\textbf{Hence, the actual prediction has two parts:} 

Determining Head & Dependent using LEFTARC/RIGHTARC

Determining their relation type i.e functions like conj, iobj, etc.

Hence the actual predictions done by Oracle are something like this LEFTARC\_IOBJ, RIGHTARC\_CONJ, etc.

For simplicity purposes, I will be using just LEFTARC & RIGHTARC notations in the below example.

We will stop as & when the buffer is empty & stack has only [Root] left.

\textbf{Slide 19:}

\textbf{Some restrictions are also considered:} 

No LEFTARC, if we have [Root] as 2nd element of the stack & buffer, is non-empty

At least 2 elements should be Stack for LEFTARC/RIGHTARC transition.

[Root] can never be a dependent (as it is a dummy variable)

\textbf{Slide 20:}

A Dependency Treebank is a text corpus in which each sentence has a corresponding dependency tree which is usually either extracted using Syntactic Treebanks(like PennTreebank) or manually marked by humans.

\textbf{Slide 21:}

Examples

consider the statement: ‘book me the morning flight’. We will be creating a dependency parse tree for this sentence

\textbf{Slide 22:}

Step-0: As Stack has only one element [root], SHIFT is predicted(exception cases mentioned above) & ‘book’ is pushed.

Step-1:[root, book] are popped out from Stack, Oracle predicts SHIFT hence they are put back, and ‘me’ is pushed in Stack.

Step-2:[book, me](top 2 elements of the Stack currently) are popped, Oracle predicts RIGTHARC, hence declaring ‘book’ as Head & ‘me’ as dependent & removing ‘me’ from Stack; pushing back ‘book’ & adding relation:(Book→me)

At last, we can pop ‘root’ & ‘book’ as the buffer is empty(which we avoided in Step-0 due to the exception mentioned). RIGTHARC is the relation detected & with this the parse tree is obtained. Using the relations found, a directed graph can be easily established with ‘book’ as the root node.



\end{document}
